\documentclass[11pt]{article}
\usepackage[margin=1in]{geometry}
\usepackage{amsmath,amssymb,amsthm}
\usepackage{mathrsfs}
\usepackage[hidelinks]{hyperref}
\usepackage{enumitem}
\setlength{\emergencystretch}{3em}

\newtheorem{theorem}{Theorem}
\newtheorem{lemma}[theorem]{Lemma}
\newtheorem{proposition}[theorem]{Proposition}
\newtheorem{corollary}[theorem]{Corollary}
\newtheorem{conjecture}[theorem]{Conjecture}
\newtheorem*{lemmaA}{Lemma A}
\theoremstyle{definition}
\newtheorem{definition}[theorem]{Definition}
\theoremstyle{remark}
\newtheorem{remark}[theorem]{Remark}

\newcommand{\R}{\mathbb{R}}
\newcommand{\Z}{\mathbb{Z}}
\newcommand{\N}{\mathbb{N}}
\newcommand{\Q}{\mathbb{Q}}
\newcommand{\ee}{\mathrm{e}}
\newcommand{\dist}{\operatorname{dist}}
\newcommand{\dimH}{\dim_{\mathrm{H}}}
\newcommand{\Lam}{\Lambda}

\title{Complete spectrum discretization, exact Hausdorff gauges,\\
and dimension phase transitions for Diophantine approximation\\
along the convergent denominators of \texorpdfstring{$e$}{e}}
\author{Anonymous author}
\date{Integrated draft, November 21, 2026}

\begin{document}
\maketitle

\begin{abstract}
Let $(Q_j)_{j\ge1}$ be the selected convergent denominators of $\ee$
immediately preceding the partial quotients $2j$
($Q_1=1,Q_2=7$, $Q_{j+1}=(4j+2)Q_j+Q_{j-1}$), let $\|y\|=\dist(y,\Z)$, and
define the weighted approximation exponent
$\Lam(x)=\limsup_{j\to\infty} j\,\|Q_j x\|$. We prove that the full spectrum is
completely discrete:
$\Lam(\R)=\frac14\Z_{\ge0}\cup\{\infty\}=\{0,1/4,1/2,3/4,1,\dots\}\cup\{\infty\}$,
so every open interval $(k/4,(k+1)/4)$ ($k\in\Z_{\ge0}$) is an unconditional
empty spectral gap, and $\Lam(x)=0$ exactly on the countable group
$L=\Z+\Z\frac{\ee+1}{2}$.

Our main metric results measure the critical levels and their fine structure.
For an integer $k\ge1$ let $h_k$ be the gauge with $h_k(1/Q_n)=(2k+1)^{-n}$,
interpolated linearly. On the circle $\R/\Z$, the exact level $\{\Lam=k/4\}$
has infinite but $\sigma$-finite $\mathcal H^{h_k}$ measure, while
$\{\Lam<k/4\}$ is $\mathcal H^{h_k}$-null. The gauge $h_k$ is exact: the level
has infinite $\mathcal H^{h_k}$ measure in every nonempty open arc, every gauge
of smaller order gives $\{\Lam\le k/4\}$ measure zero, and for every gauge of
larger order $\{\Lam=k/4\}$ is not $\sigma$-finite. For arbitrary gauges,
including oscillating ones, the three measure regimes are determined by
$\liminf_n (2k+1)^n g(1/Q_n)$, or by $\liminf_n 2^n g(1/Q_n)$ for the
unsigned exact-limit set. For finite-state tail constraints the base is
the largest Perron root among components of a right-resolving
presentation supporting infinitely many extreme defects. We give
a finite algorithm for arbitrary labelled graphs and two-sided
quantitative gauge bounds on a fixed-start coded set. Furthermore, at each
positive level $c=k/4$ ($k\ge1$) we establish a sharp three-way splitting
between the limsup level (digit alphabet of size $2k+1$, critical gauge base
$2k+1$), the unsigned exact limit set $\{x:j\|Q_jx\|\to k/4\}$ (universal binary
alphabet $\{\pm k\}$, critical gauge base $2$ \emph{independently} of
$k\ge1$), and each signed exact limit set
$\{x:j(Q_jx-\mathrm{round}(Q_jx))\to\pm k/4\}$ (which collapses to a single
countable coset of $L$). More generally, for any polynomial-growth recurrence
$Q_{j+1}=(cj^r+d)Q_j+Q_{j-1}$ ($r,c\in\Z_{\ge1}$, $d\in\Z$), the critical
exponent $\Lam^{(r,c,d)}(x)=\limsup_{j\to\infty}j^r\|Q_jx\|$ has full
spectrum $\frac1c\Z_{\ge0}\cup\{\infty\}$, intrinsic rank-$2$ zero-level
lattice $L^{(r,c,d)}$, critical gauges $h_k^{(r,c,d)}(1/Q_n)=(2k+1)^{-n}$
with $1/r$-scaled logarithmic order and positive-level three-way splitting, and
subcritical Hausdorff dimension phase transition $1-\alpha/r$ ($0\le\alpha<r$).

Next, every finite level set $\{\Lam\le c\}$ has Hausdorff dimension $0$.
Hence $\{\Lam<\infty\}$ has dimension $0$, and $\Lam=\infty$ Lebesgue-almost
everywhere (with Hausdorff dimension $1$). For $0\le\alpha<1$, the sublevel set $\{\Lam_\alpha\le c\}$ for $c\ge0$
has Hausdorff dimension $1-\alpha$. The same holds for the set where
$j^\alpha\|Q_jx\|\to c$ for $c\ge0$ if $\alpha>0$, and for
$c\in[0,1/2]$ if $\alpha=0$. The unweighted case also follows from work
of Bugeaud, Liao and Rams. For ordinary continued fractions,
we determine the Hausdorff dimensions of Hermite window-count constraints and
show that eventual periodicity of the partial quotients modulo four has the same
Hausdorff dimension as the fixed alphabet $1\bmod 4$, strictly between $357/512$
and $359/512$.

The classification of $\{\Lam=0\}$, the initial and full spectral-gap ladder,
Theorem~\ref{thm:gauge},
Corollaries~\ref{cor:exact-gauge} and~\ref{cor:liminf-gauges},
Remark~\ref{rem:gauge-scale}, all three parts (i)--(iii) of
Theorem~\ref{thm:trichotomy} on both $\R$ and $\R/\Z$ (\texttt{UnitAddCircle},
including the binary exact-limit gauge regime and the signed-coset collapse
$E_{\mathrm{lim}}^\pm(k/4)=x_k^\pm+L$), the dimension-zero theorem for all
finite sublevels (and full Lebesgue measure of $\{\Lam=\infty\}$), both the upper
and lower bounds for the subcritical dimension $1-\alpha$, and—for general
polynomial-growth recurrences ($r\ge1$)—Lemmas~\ref{lem:poly-growth}--\ref{lem:poly-lattice},
Theorem~\ref{thm:general-rec} (full spectrum $\frac1c\Z_{\ge0}\cup\{\infty\}$),
the critical-gauge conclusions of Theorem~\ref{thm:poly-gauge-trichotomy}
and Corollary~\ref{cor:poly-exact-gauge}, and the critical and supercritical
claims and both subcritical dimension bounds in Theorem~\ref{thm:poly-phase-transition}
are machine-checked in Lean~4 using the base axioms \texttt{propext},
\texttt{Classical.choice} and \texttt{Quot.sound}. The continued-fraction
window and residue theorems
(Theorems~\ref{thm:windows}--\ref{thm:residues}, beyond the $14$ core lemmas in
\texttt{WindowAndResidueCore.lean}), and the finite-state and frequency
theorems (Theorems~\ref{thm:finite-state-tail}, \ref{thm:finite-state-quant},
\ref{thm:pressure-frequency}, and~\ref{thm:edge-frequency-circulation})
are paper proofs (with an exact
integer-arithmetic certificate for the residue interval). The classification and
the first gap are a subsequence analogue of a criterion of Kraaikamp and
Liardet. The work is a spin-off of a project on the irrationality of $e+\pi$ and
does \emph{not} resolve that open problem.

For strongly connected finite-state constraints we also study exact
frequencies of the extreme labels. A tilted Perron root computes
the critical exponential gauge rate. A block construction determines
the critical case away from the maximal-entropy frequency, including
the full-alphabet endpoints.
Prescribed edge frequencies are governed by probability circulations
and conditional entropy. At extreme scalar frequencies the exact
rate is read from a tight-cycle subgraph, and a component rule
extends the critical classification to reducible presentations.
\end{abstract}

\section{Introduction}

This article studies the metric theory of Diophantine approximation along a
\emph{sparse, super-exponentially growing} subsequence of denominators: the
selected convergent denominators $(Q_j)_{j\ge1}$ of $\ee$ (notation in the
Setup below). We first show in Theorem~\ref{thm:full-spectrum} that the integer
three-term recurrence for $(Q_j)$ forces complete discretization of the weighted
approximation spectrum:
\[
\Lam(\R)=\frac14\Z_{\ge0}\cup\{\infty\}
=\left\{0,\frac14,\frac12,\frac34,1,\dots\right\}\cup\{\infty\},
\]
creating an infinite family of equidistant empty spectral gaps
$(k/4,(k+1)/4)$ of width $1/4$ and forcing the sublevel sets
$E_c=\{x:\Lam(x)\le c\}$ to be constant across each gap
($E_c=E_{k/4}$ for $c\in[k/4,(k+1)/4)$, Corollary~\ref{thm:sub}).

Our main metric result at the critical weight, Theorem~\ref{thm:gauge},
measures each positive quantized level of $\Lam$: for each integer $k\ge1$
there is an explicit Hausdorff gauge $h_k$ with $h_k(1/Q_n)=(2k+1)^{-n}$, built
from the recurrence itself, for which the exact level $\{\Lam=k/4\}$ has
infinite but $\sigma$-finite Hausdorff measure $\mathcal H^{h_k}$ on the circle
$\R/\Z$, while the strict sublevel $\{\Lam<k/4\}$ is $\mathcal H^{h_k}$-null.
This differs both from the $h$-sets of Cabrelli, Mendivil, Molter and
Shonkwiler \cite[\S5, Theorem~3]{CMMS}, whose gauge gives a finite positive
measure, and from the classical Lagrange level at value $3$ studied by Moreira,
Erazo and Angelini \cite{MEA}, which has either zero or non-$\sigma$-finite
measure for every dimension function. Corollary~\ref{cor:exact-gauge} shows
that $h_k$ is the exact critical gauge of the level $k/4$: smaller gauges give
measure zero even on the non-strict sublevel $\{\Lam\le k/4\}$, larger gauges
destroy $\sigma$-finiteness on $\{\Lam=k/4\}$, and the $\mathcal H^{h_k}$
measure of $\{\Lam=k/4\}$ is infinite in every nonempty open arc.
Corollary~\ref{cor:liminf-gauges} classifies every nondecreasing gauge by
its lower limit along the recurrence scales, even when the normalized
cost oscillates. Theorem~\ref{thm:finite-state-tail} shows how this
classification changes under constraints on consecutive defects.
The largest Perron root is taken only over components in which a path
can revisit an extreme defect infinitely often.
Furthermore, Theorem~\ref{thm:trichotomy} establishes a sharp three-way
splitting at each positive spectral level $k/4$ ($k\ge1$) comparing the limsup
level $\{\Lam=k/4\}$ (gauge base $2k+1$), the unsigned exact limit set
$\{x:j\|Q_jx\|\to k/4\}$ (universal critical gauge base $2$ for \emph{all}
$k\ge1$), and the signed exact limit sets
$\{x:j(Q_jx-\mathrm{round}(Q_jx))\to\pm k/4\}$ (each of which collapses to a
single countable $L$-coset). Theorems~\ref{thm:general-rec}--\ref{thm:poly-phase-transition}
extend the quantization, reference gauges, three-way splitting and
dimension results to general polynomial-growth recurrences
$Q_{j+1}=(cj^r+d)Q_j+Q_{j-1}$ ($r,c\in\Z_{\ge1}$, $d\in\Z$), replacing the
$\ee$-convergent lattice $L$ by an intrinsic rank-$2$ lattice $L^{(r,c,d)}$
spanned by fundamental solutions and establishing the discretized spectrum
$\frac1c\Z_{\ge0}\cup\{\infty\}$, exact gauges with $1/r$-scaled logarithmic
asymptotics, positive-level trichotomy, and subcritical dimension phase
transition $1-\alpha/r$ ($0\le\alpha<r$).

On the coarse scale of Hausdorff dimension, Theorem~\ref{thm:finite} shows that
every finite level set $E_c$ has Hausdorff dimension $0$, so $\{\Lam<\infty\}$
has dimension $0$ while its complement $\{\Lam=\infty\}$ has full Lebesgue
measure and Hausdorff dimension $1$. At subcritical weights
$\Lam_\alpha(x)=\limsup_{j\to\infty}j^\alpha\|Q_jx\|$ ($0\le\alpha<1$),
Theorems~\ref{thm:upper} and~\ref{conj:weight} give exact Hausdorff dimension $1-\alpha$ for
$\{x:\Lam_\alpha(x)\le c\}$ for every $c\ge0$, and for
$\{x\in[0,1]:j^\alpha\|Q_jx\|\to c\}$ for every admissible target $c$.
For $\alpha=0$ such targets are $0\le c\le 1/2$.
The proof uses a covering bound and a shifted-grid Moran construction;
the case $\alpha=0$ also follows from Bugeaud, Liao and Rams \cite{BLR}, and
both bounds for $0\le\alpha<1$ are machine-checked in Lean~4. Finally, in
Section~\ref{sec:windows-residues}, Theorems~\ref{thm:windows}--\ref{thm:no-short-dim}
determine the Hausdorff dimensions and almost-sure frequency of Hermite
continued-fraction window constraints, and Theorem~\ref{thm:residues} proves
that eventual periodicity of ordinary continued-fraction digits modulo $4$ has
the same Hausdorff dimension as the fixed alphabet $1\bmod 4$, bounded strictly
between $357/512$ and $359/512$.

These results build on the zero-level classification and initial spectral gap
(Section~\ref{sec:class}), which arose as machinery in a separate formalization
project on the (open) irrationality of $e+\pi$ \cite{Phase1}.

\subsection{Setup}

The continued fraction of $\ee$ is $[2;1,2,1,1,4,1,1,6,\dots]$. Let $P_j/Q_j$
be the convergent immediately preceding the partial quotient $2j$; the first
few are $3/1,19/7,193/71,2721/1001$. The denominators satisfy
\begin{equation}\label{eq:rec}
Q_{j+1}=(4j+2)Q_j+Q_{j-1}\quad(j\ge2),\qquad Q_1=1,\ Q_2=7,
\end{equation}
so $Q_{j+1}/Q_j\sim 4j$ and $\log Q_j\sim j\log j$ (super-exponential growth).
Write $\|y\|=\dist(y,\Z)$ and define
\[
\Lam(x):=\limsup_{j\to\infty} j\,\|Q_j x\|\in[0,\infty].
\]
More generally, for $\alpha\ge0$ put
$\Lam_\alpha(x):=\limsup_{j\to\infty} j^\alpha\|Q_j x\|\in[0,\infty]$,
so that $\Lam=\Lam_1$.
Let $L:=\Z+\Z\frac{\ee+1}{2}$, a countable dense subgroup of $\R$.

\subsection{Relation to prior work}\label{sec:prior}

The \emph{phenomenon} (a spectral gap plus zero-dimensional sub-threshold level
sets) is classical: in the classical Lagrange spectrum, the spectrum below
$\sqrt5$ (and below $3$) is countable (Hausdorff dimension $0$), and Hall's ray
gives a half-line. Moreira \cite{Moreira} studies the Hausdorff dimension of
portions of the classical Markov and Lagrange spectra below a threshold. Lima
and Moreira \cite{LimaMoreira} prove phase transitions for the \emph{dynamical}
Markov and Lagrange spectra of generic horseshoes; the corresponding question
for the classical spectra is left open there. Note that our level sets $E_c$ are
subsets of the space of inputs $x\in\R$.

The closest prior continued-fraction result is a theorem of Kraaikamp and
Liardet \cite{KL}, quoted as Lemma~4.1 in \cite{BDS}. Let $\alpha$ be irrational
with convergents $p_n/q_n$. Then a real $\beta$ lies outside
$\Z+\Z\alpha$ if and only if
$\|q_n\beta\|>\tfrac14\,q_n|q_n\alpha-p_n|$ for infinitely many $n$. This
is a classification together with a gap, with constant $1/4$, for the
normalized quantity $\|q_n\beta\|/(q_n|q_n\alpha-p_n|)$ along \emph{all}
convergent denominators. Theorems~\ref{thm:class} and~\ref{thm:gap} can be
compared with this criterion in two complementary ways:
\begin{itemize}[leftmargin=1.8em]
\item \emph{As a subsequence along the convergents of $\ee$:} Passing from all
convergent denominators of $\ee$ to the subsequence $(Q_j)$ of every third
convergent denominator changes two things. First, the zero group grows from
$\Z+\Z\ee$ to $L$: for $x=(\ee+1)/2$ the quantity $j\|Q_jx\|$ tends to $0$, so
the infinitely many $n$ supplied by Kraaikamp and Liardet for $\alpha=\ee$ lie
outside the subsequence $(Q_j)$. Second, the normalization changes. Since
$2j+1<1/(Q_j|Q_j\ee-P_j|)<2j+2$, one has
$\limsup_{j\to\infty}\|Q_jx\|/(Q_j|Q_j\ee-P_j|)=2\Lam(x)$. Thus the gap
$\Lam\ge1/4$ reads $\limsup\ge1/2$ in their normalization, and the value
$\Lam(x_0)=1/4$ is attained by the explicit point $x_0$ in
Remark~\ref{rem:unweighted}. The two constants $1/4$ therefore refer to
different normalizations.
\item \emph{As the full convergent sequence of $\beta=(\ee+1)/2$:} Put
$R_j=(P_j+Q_j)/2\in\Z$. The fractions $R_j/Q_j$ are the convergents of
$\beta=(\ee+1)/2=[1;1,6,10,14,\dots]$, whose initial convergent pairs are
$(1,1)$ and $(2,1)$ and whose remaining partial quotients are $4j+2$ for
$j\ge1$. Applying the Kraaikamp--Liardet criterion \cite{KL} (in the form
recorded on p.~106 of \cite{BDS}) directly to $\beta$ gives
$\|Q_jx\|>\tfrac14 Q_j|Q_j\beta-R_j|$ infinitely often for $x\notin L$. Since
the continued-fraction error formula gives $jQ_j|Q_j\beta-R_j|\to1/4$, this
black-box application yields $\Lam(x)\ge1/16$ outside $L$, which already
recovers the zero-level classification $\{\Lam=0\}=L$, whereas the direct
three-term recurrence analysis in Theorems~\ref{thm:gap}
and~\ref{thm:full-spectrum} yields the sharp initial gap $1/4$ and the complete
quantization $\Lam(\R)=\frac14\Z_{\ge0}\cup\{\infty\}$.
\end{itemize}
For a general irrational $\alpha$ whose partial quotients tend to infinity, an
analogous quantization of the spectrum should follow from the Ostrowski
expansion of $\beta$ with respect to $\alpha$; we do not pursue that direction
here, and instead treat general polynomial-growth denominator recurrences
$Q_{j+1}=(cj^r+d)Q_j+Q_{j-1}$ in
Theorems~\ref{thm:general-rec}--\ref{thm:poly-phase-transition}. Accordingly, we
regard Theorems~\ref{thm:class} and~\ref{thm:gap} as a variant of known results.
The main new contributions are the complete spectrum discretization and
three-way positive-level splitting (Theorems~\ref{thm:full-spectrum},
\ref{thm:trichotomy}, and~\ref{thm:general-rec}--\ref{thm:poly-phase-transition}), the critical gauge theorems
(Theorem~\ref{thm:gauge} and Corollary~\ref{cor:exact-gauge}), the machine-checked
dimension-zero and subcritical dimension theorems (Theorems~\ref{thm:finite},
\ref{thm:upper}, and~\ref{conj:weight}), and the continued-fraction window and
residue theorems (Theorems~\ref{thm:windows}--\ref{thm:residues}).

For the selected sequence $(Q_j)$, Corollary~8 of Bugeaud, Liao and Rams
\cite{BLR} applies directly since $Q_{j+1}/Q_j\to\infty$. It gives dimension
one for every constant target $c\in[0,1/2]$ at weight $\alpha=0$ in
Theorem~\ref{conj:weight} (in particular for $\{x:\|Q_jx\|\to0\}$), by taking
the centres equal to $c$ and using $|\|y\|-\|c\||\le\|y-c\|$; integer
translation invariance transfers full dimension between $\R$ and $[0,1]$. For
nonzero weights $0<\alpha<1$, the logarithmic growth condition (1.6) of their
Theorem~7 fails for the linear maps $x\mapsto Q_jx$ because
$\log Q_{j+1}/\log Q_j\to1$, and classical mass-transference methods for
$\limsup$ shrinking-target sets \cite{BV} do not apply to the eventual
(liminf-type) bounds defining $E_c$ and the prescribed-limit sets. We establish
the exact dimension $1-\alpha$ here by an elementary covering upper bound and a
shifted-grid Moran construction.

Regarding Hausdorff gauges at critical levels, Cabrelli, Mendivil, Molter and
Shonkwiler \cite[\S5, Theorem~3]{CMMS} construct a gauge giving positive
\emph{finite} measure to a Cantor set from its gap sequence, whereas Moreira,
Erazo and Angelini \cite{MEA} show that the classical Lagrange level at value
$3$ has either zero or non-$\sigma$-finite measure for every dimension function.
Here the recurrence specifies the gauge $h_k(1/Q_n)=(2k+1)^{-n}$, and the entire
critical level $\{\Lam=k/4\}$ on $\R/\Z$ has \emph{infinite, $\sigma$-finite}
$\mathcal H^{h_k}$ measure, proved by combining a countable decomposition by
initial recurrence data with translation invariance under the dense subgroup
$L$.

Relative to earlier literature, our setting has four distinguishing features:
(i) a \emph{subsequence} of convergent denominators of $\ee$ (every third),
(ii) the critical \emph{weight} $j$, (iii) the lattice
$L=\Z+\Z\frac{\ee+1}{2}$ (note $(\ee+1)/2$, not $\ee$), and (iv) the
\emph{sharp} quantum $1/4$. The weight $j$ is \emph{critical} for the covering
argument: it makes the per-parent branching number uniformly bounded, yielding
dimension $0$ for \emph{all} finite level sets (Theorem~\ref{thm:finite}),
whereas without the weight the branching grows linearly and the unweighted
zero-set $\{x:\|Q_jx\|\to0\}$ has full dimension $1$, while the weighted
zero-set $\{x:j\|Q_jx\|\to0\}$ equals the countable group $L$.

\section{Classification and initial spectral gap}\label{sec:class}

The two results of this section were first proved, and Lean-kernel-checked at
Layer~A, in the Phase-1 project \cite{Phase1}. Here Layer~A means that the
audited declaration in Lean~4 depends only on Lean's three standard base axioms
\texttt{propext}, \texttt{Classical.choice} and \texttt{Quot.sound}. For
completeness we give short self-contained proofs. They use three facts about the
selected convergents $P_j/Q_j$:
\begin{enumerate}[label=(\arabic*),leftmargin=2em]
\item $P_j$ and $Q_j$ are odd, and $P_j$ satisfies the same recurrence
\eqref{eq:rec} as $Q_j$.
\item $P_jQ_{j+1}-P_{j+1}Q_j=2(-1)^{j+1}$ for all $j\ge1$.
\item The approximation error $\delta_j:=Q_j\ee-P_j$ satisfies $\delta_j\to0$
and $|\delta_j|<1/(2jQ_j)$ for all $j\ge1$.
\end{enumerate}
The last estimate is not in mathlib v4.35, and earlier drafts of the Phase-1
formalization assumed it as an axiom. Both $\delta_j\to0$ and
$|\delta_j|<1/(2jQ_j)$ are now kernel-checked Lean theorems
(\texttt{deltaSel\_tendsto\_zero} and \texttt{abs\_deltaSel\_lt}, proved via
\texttt{DeltaSign.lean}), derived from the Hermite--Pad\'e integral
representation
\[
(-1)^n n!\,\delta_n=\int_0^1 \ee^t t^n(1-t)^n\,dt>0.
\]

\paragraph{Digits.} For $x\in\R$ and $j\ge1$ let $M_j=\mathrm{round}(Q_jx)\in\Z$
be a nearest integer to $Q_jx$, put $u_j=Q_jx-M_j\in[-1/2,1/2]$ (so
$|u_j|=\|Q_jx\|\le1/2$), and for $j\ge2$ define the integer recurrence defect
(digit)
\begin{equation}\label{eq:defect-id}
 d_j:=M_{j+1}-A_jM_j-M_{j-1}=A_ju_j+u_{j-1}-u_{j+1}\in\Z,\qquad A_j:=4j+2,
\end{equation}
where the second equality follows immediately from \eqref{eq:rec}.

\begin{lemmaA}
If $d_j=0$ for all $j\ge N$, then $x\in L$.
\end{lemmaA}

\begin{proof}
Put $v_j=(P_j+Q_j)/2$, an integer since $P_j$ and $Q_j$ are odd. The
vectors $(v_N,v_{N+1})$ and $(Q_N,Q_{N+1})$ have determinant
$\tfrac12(P_NQ_{N+1}-P_{N+1}Q_N)=\pm1$, so they form a basis of $\Z^2$.
Choose integers $a,b$ with $M_N=av_N+bQ_N$ and $M_{N+1}=av_{N+1}+bQ_{N+1}$.
The sequences $(M_j)$, $(v_j)$ and $(Q_j)$ all satisfy \eqref{eq:rec}
from index $N$ on, hence $M_j=av_j+bQ_j$ for all $j\ge N$. With
$y:=x-b-a\frac{\ee+1}{2}$ this gives
\[
 u_j=Q_jx-M_j=Q_jy+\tfrac a2\,\delta_j .
\]
Since $|u_j|\le1/2$, $\delta_j\to0$ and $Q_j\to\infty$, we get $y=0$, that
is, $x=b+a\frac{\ee+1}{2}\in L$.
\end{proof}

\begin{theorem}[Classification, {\cite[Thm.~2]{Phase1}}]\label{thm:class}
For $x\in\R$, $\quad j\|Q_jx\|\to0\iff x\in L.$
\end{theorem}

\begin{proof}
If $x=m+n\frac{\ee+1}{2}$ with $m,n\in\Z$, then
$Q_jx=Q_jm+nv_j+\tfrac n2\delta_j$ with $v_j$ as above, so
$\|Q_jx\|\le|n|\,|\delta_j|/2<|n|/(4jQ_j)$ and $j\|Q_jx\|\to0$.
Conversely, if $j|u_j|\to0$, then
$|d_j|\le A_j|u_j|+|u_{j-1}|+|u_{j+1}|\to0$. The $d_j$ are integers, so
$d_j=0$ eventually, and $x\in L$ by Lemma~A.
\end{proof}

\begin{theorem}[Gap, {\cite[Cor.~3]{Phase1}}]\label{thm:gap}
If $x\notin L$ and $c<1/4$, then it is \emph{not} the case that
$j\|Q_jx\|\le c$ eventually. Equivalently $\Lam(x)\ge1/4$ for $x\notin L$.
\end{theorem}

\begin{proof}
Suppose $j|u_j|\le c<1/4$ for all $j\ge J$. For $j\ge J+1$,
\[
 |d_j|\le\frac{(4j+2)c}{j}+\frac{c}{j-1}+\frac{c}{j+1}
 \le4c+\frac{4c}{j-1},
\]
which is less than $1$ for all large $j$. Hence $d_j=0$ eventually, and
$x\in L$ by Lemma~A.
\end{proof}

\paragraph{Lean status.} Theorems~\ref{thm:class} and~\ref{thm:gap} are
kernel-checked at Layer~A in the module
\texttt{EpiFormal/\allowbreak Classification.lean} as
\texttt{theorem2\_\allowbreak classification} and
\texttt{corollary3\_\allowbreak gap}. Their axiom audit reports
only \texttt{propext}, \texttt{Classical.choice} and \texttt{Quot.sound}.

\section{Complete spectrum discretization and step-function collapse}\label{sec:spectrum}

\begin{corollary}[Initial spectral gap, Layer~A in Lean]\label{thm:spectral}
For every $x\in\R$, either $\Lam(x)=0$ or $\Lam(x)\ge1/4$. Thus
$\Lam(\R)\subseteq\{0\}\cup[1/4,\infty]$, and $\Lam(x)=0\iff x\in L$.
\end{corollary}

This is an immediate consequence of Theorems~\ref{thm:class}
and~\ref{thm:gap}; we state it separately because it is the form used
below and the form directly machine-checked in \texttt{SpectralGap.lean}.

\begin{proof}
If $x\in L$, Theorem~\ref{thm:class} gives $j\|Q_jx\|\to0$, hence $\Lam(x)=0$.
If $x\notin L$, Theorem~\ref{thm:gap} gives $\Lam(x)\ge1/4$. The dichotomy and
the characterization of the zero value follow.
\end{proof}

\begin{remark}\label{rem:linear-indep}
In particular $\Lam(\ee)=0$, since $\ee=2\cdot\frac{\ee+1}{2}-1\in L$; $\ee$ is
\emph{not} a generic point for this exponent. For $x=B\pi$ ($B=1,\dots,10$),
$\sqrt2$, the golden ratio and random reals, the values $j\|Q_jx\|$ for
$j\le60$ are large (Section~\ref{sec:num}), but such finite data do not
determine $\Lam(x)$. By Corollary~\ref{thm:spectral}, the statement
``$\Lam(B\pi)>0$ for every integer $B\ge1$'' is equivalent to
$B\pi\notin L$ for all $B\ge1$, i.e.\ to the linear independence of
$1,\ee,\pi$ over $\Q$, which is open. Nothing in this article decides it.
\end{remark}

We now show that the initial gap $(0,1/4)$ is merely the bottom rung of an
infinite ladder of equidistant spectral gaps: the entire finite spectrum of
$\Lam$ is quantized in multiples of $1/4$.

\begin{theorem}[Complete discretization of the spectrum, Layer~A in Lean]\label{thm:full-spectrum}
For every $x\in\R$ with $\Lam(x)<\infty$, one has $u_j=O(1/j)\to0$ and
\begin{equation}\label{eq:digit-asymp}
\lim_{j\to\infty}\bigl(d_j-4ju_j\bigr)=0,\qquad
\lim_{j\to\infty}\bigl(|d_j|-4j\|Q_jx\|\bigr)=0.
\end{equation}
Consequently $4\Lam(x)=\limsup_{j\to\infty}|d_j|\in\Z_{\ge0}$ whenever
$\Lam(x)<\infty$, and the full spectrum of $\Lam$ is given by the exact identity
\begin{equation}\label{eq:full-spectrum}
\Lam(\R)=\frac14\Z_{\ge0}\cup\{\infty\}
=\left\{0,\frac14,\frac12,\frac34,1,\dots\right\}\cup\{\infty\}.
\end{equation}
In particular, for every integer $k\ge0$, the open interval
$\bigl(\frac{k}{4},\frac{k+1}{4}\bigr)$ is an unconditional empty spectral gap
for $\Lam$.
(Lean, module \texttt{FullSpectrumDiscrete}:
\texttt{euler\_\allowbreak signed\_\allowbreak defect\_\allowbreak approx},
\texttt{euler\_\allowbreak abs\_\allowbreak defect\_\allowbreak approx}, and
\texttt{spectral\_\allowbreak gap\_\allowbreak open\_\allowbreak interval}.)
\end{theorem}

\begin{proof}
If $\Lam(x)<\infty$, then $\|Q_jx\|=|u_j|\le(\Lam(x)+1)/j$ for all large $j$,
so $u_j\to0$. Rearranging \eqref{eq:defect-id} gives
$d_j-4ju_j=2u_j+u_{j-1}-u_{j+1}\to0$, and the reverse triangle inequality
$\bigl||d_j|-4j|u_j|\bigr|\le|d_j-4ju_j|$ yields \eqref{eq:digit-asymp}.
Since $|d_j|\in\Z_{\ge0}$, the sequence $(|d_j|)$ can only have a nonnegative
integer or $\infty$ as its limit superior. By \eqref{eq:digit-asymp},
$4\Lam(x)=\limsup_{j\to\infty}4j\|Q_jx\|=\limsup_{j\to\infty}|d_j|\in\Z_{\ge0}\cup\{\infty\}$,
which excludes every open interval $(k/4,(k+1)/4)$ from $\Lam(\R)$.
Conversely, $\Lam(0)=0$, $\Lam=\infty$ holds on a set of full Lebesgue measure
(Section~\ref{sec:finite-levels}), and for each integer $k\ge1$ choosing a
sufficiently large starting index $N$ and constant digits $d_j\equiv k$ for
$j\ge N$ (as in the proof of Theorem~\ref{thm:gauge} below) produces a point
$x_k\in\R$ with $|u_j|<1/8$ for $j\ge N$, $M_j=\mathrm{round}(Q_jx_k)$, and
$4j\|Q_jx_k\|\to k$, so $\Lam(x_k)=k/4$. Thus every value in
$\frac14\Z_{\ge0}\cup\{\infty\}$ is attained.
\end{proof}

For $c\ge0$ write $E_c:=\{x\in\R:\Lam(x)\le c\}$.

\begin{corollary}[Sub-threshold level sets and gap constancy, Layer~A in Lean]\label{thm:sub}
For every integer $k\ge0$ and every $c\in\bigl[\frac{k}{4},\frac{k+1}{4}\bigr)$,
one has $E_c=E_{k/4}=\{x\in\R:u_j(x)\to0\text{ and }\limsup_{j\to\infty}|d_j(x)|\le k\}$.
In particular, if $c<1/4$ ($k=0$), then $E_c=E_0=L$, so $E_c$ is countable and
$\dimH E_c=0$.
(Lean, modules \texttt{FullSpectrumDiscrete} and \texttt{SpectralGap}:
\texttt{eventually\_\allowbreak abs\_\allowbreak defect\_\allowbreak le\_\allowbreak of\_\allowbreak apprExponent\_\allowbreak le},
\texttt{levelSet\_\allowbreak eq\_\allowbreak on\_\allowbreak gap},
\texttt{levelSet\_\allowbreak below\_\allowbreak eq\_\allowbreak lattice}, and
\texttt{dimH\_\allowbreak levelSet\_\allowbreak below}.)
\end{corollary}

\begin{proof}
Clearly $E_{k/4}\subseteq E_c$ for $c\ge k/4$. Conversely, if $x\in E_c$ with
$c<(k+1)/4$, then $\Lam(x)\le c<(k+1)/4$ cannot lie in the empty spectral gap
$(k/4,(k+1)/4)$ by Theorem~\ref{thm:full-spectrum}, so $\Lam(x)\le k/4$ and
eventually $|d_j(x)|\le k$. For $k=0$, $\Lam(x)=0\iff x\in L$ by
Corollary~\ref{thm:spectral}, and $L$ is a countable image of $\Z^2$, hence has
Hausdorff dimension $0$ \cite{Falconer}.
\end{proof}

\paragraph{Lean status.} Corollaries~\ref{thm:spectral} and~\ref{thm:sub} (for
$c<1/4$), together with the countability of $L$ and $\dimH E_c=0$ for $c<1/4$,
are kernel-checked at Layer~A in \texttt{EpiFormal/\allowbreak SpectralGap.lean}
(five declarations: \texttt{apprExponent\_eq\_zero\_iff},
\texttt{spectral\_gap}, \texttt{levelSet\_below\_eq\_lattice},
\texttt{latticeL\_countable}, and \texttt{dimH\_levelSet\_below}; no
\texttt{sorryAx}). The full spectrum discretization, open-interval spectral-gap
exclusion $(k/4,(k+1)/4)$, and the step-function constancy $E_c=E_{k/4}$ across
each gap are kernel-checked at Layer~A in
\texttt{EpiFormal/\allowbreak FullSpectrumDiscrete.lean} (15 declarations, base
axioms only, no \texttt{sorryAx}). The integrated revision reran \texttt{Audit.lean} on the existing
dependency checkout. It reports only the three base axioms for the
classification and spectral-discretization theorems. The two accompanying
files for the arbitrary-gauge and finite-state results were checked directly.
The experimental module
\texttt{Trichotomy\allowbreak And\allowbreak Poly\allowbreak Extension.lean}
has unfinished proofs and is outside this verification.

\section{All finite level sets have dimension zero, and the complementary infinite level}\label{sec:finite-levels}

The sub-threshold result ($c<1/4$) uses the sharp classification. For
$c\ge1/4$ the level set $E_c$ is uncountable (Theorem~\ref{thm:gauge}), but we
show it still has Hausdorff dimension $0$. The entire covering argument ---
factorial growth of $Q_j$, uniform branching bound, surviving-center recursion,
periodicity reduction, Hausdorff-content estimate, and the countable-union
assembly --- is \emph{machine-checked at Layer~A}
(\texttt{EpiFormal/DimHCore.lean} + \texttt{EpiFormal/DimHFinite.lean}, base
axioms only, no \texttt{sorryAx}, and \emph{no} continued-fraction axioms).
The elementary proof below mirrors the formalization.

\begin{lemma}[Factorial growth, Layer~A in Lean]\label{lem:factorial}
For $j\ge1$, $Q_j\ge j!$. (Lean: \texttt{DimHCore.\allowbreak Qsel\_ge\_factorial}.)
\end{lemma}

\begin{proof}
We proceed by induction on $j\ge1$. For $j=1$, $Q_1=1=1!$. For $j=2$,
$Q_2=7\ge2=2!$. For $j\ge2$, every term $Q_i$ is positive, so the recurrence
\eqref{eq:rec} gives
\[
Q_{j+1}=(4j+2)Q_j+Q_{j-1}>(j+1)Q_j\ge(j+1)j!=(j+1)!,
\]
completing the induction.
\end{proof}

\begin{lemma}[Uniform branching bound, Layer~A in Lean]\label{lem:branch}
Fix $t\ge0$, $j\ge1$ and a parent center $c\in\Z$. The number of child centers
$k'\in\{0,1,\dots,Q_{j+1}\}$ with
\[
\left|\frac{c}{Q_j}-\frac{k'}{Q_{j+1}}\right|\le \frac{t}{jQ_j}+\frac{t}{(j+1)Q_{j+1}}
\]
is at most $B:=\lfloor16t\rfloor+1$, \emph{uniformly in $j$ and $c$}.
(Lean: \texttt{DimHCore.\allowbreak fiber\_card\_le}.)
\end{lemma}

\begin{proof}
Put $R_j:=t/(jQ_j)+t/((j+1)Q_{j+1})$. Multiplying the distance inequality
$|c/Q_j-k'/Q_{j+1}|\le R_j$ by $Q_{j+1}>0$ shows that every qualifying integer
$k'$ lies in the closed real interval $[A,A+W]$ with left endpoint
$A:=Q_{j+1}(c/Q_j-R_j)$ and width
\[
W:=2R_jQ_{j+1}=2t\left(\frac{Q_{j+1}}{jQ_j}+\frac{1}{j+1}\right).
\]
Since $Q_0=Q_1=1$ and $Q_{i-1}\le Q_i$ for all $i\ge1$, the recurrence
\eqref{eq:rec} implies $Q_{j+1}=(4j+2)Q_j+Q_{j-1}\le(4j+3)Q_j\le7jQ_j$ for all
$j\ge1$, while $1/(j+1)\le1$. Hence
\[
W\le 2t(7+1)=16t.
\]
For any qualifying $k'$, we have $0\le k'-A\le W\le16t$, so the map
$k'\mapsto\lfloor k'-A\rfloor$ takes values in $\{0,1,\dots,\lfloor16t\rfloor\}$
and is injective on integers (if $k'_1<k'_2$, then $k'_2-A\ge(k'_1-A)+1$, so
$\lfloor k'_2-A\rfloor\ge\lfloor k'_1-A\rfloor+1$). Therefore the number of
qualifying integers $k'$ is at most $\lfloor16t\rfloor+1=B$.
\end{proof}

\begin{theorem}[Finite level sets, Layer~A in Lean]\label{thm:finite}
For every $c\in\R$, $\dimH E_c=0$. Consequently
$F:=\{x\in\R:\Lam(x)<\infty\}=\bigcup_{n\ge1}E_n$ has $\dimH F=0$.
(Lean: \texttt{DimHFinite.\allowbreak dimH\_levelSet\_finite}; \texttt{\#print axioms}
returns \texttt{[propext, Classical.choice, Quot.sound]} only.)
\end{theorem}

\begin{proof}
Fix $c\in\R$ and choose a rational or real threshold $t>\max(c,0)$ (in Lean,
$t=\max(c,0)+1>0$). First note that for every $m\in\Z$ and $j\ge1$, since
$Q_j\in\Z$, we have $\|Q_j(x+m)\|=\|Q_jx+Q_jm\|=\|Q_jx\|$, and hence
$\Lam(x+m)=\Lam(x)$. Thus $E_c$ is invariant under integer translations and
decomposes as the countable union
\[
E_c=\bigcup_{m\in\Z}\bigl(m+(E_c\cap[0,1])\bigr).
\]
Because integer translations are isometries and Hausdorff dimension is stable
under countable unions ($\dimH(\bigcup_m X_m)=\sup_m\dimH X_m$), it suffices to
prove $\dimH(E_c\cap[0,1])=0$.

For each $j\ge1$, define $r_j:=t/(jQ_j)$ and
$A_j(t):=\{x\in\R:\|Q_jx\|\le t/j\}$. If $x\in E_c\cap[0,1]$, then
$\limsup_{j\to\infty}j\|Q_jx\|=\Lam(x)\le c<t$, which by definition of the limit
superior means that $j\|Q_jx\|\le t$ \emph{eventually} (a liminf-type tail
condition, \emph{not} a limsup set). Therefore there exists an integer $J\ge1$
such that $x\in A_j(t)$ for all $j\ge J$. Consequently,
\[
E_c\cap[0,1]\subseteq\bigcup_{J=1}^\infty C_J,\qquad
C_J:=[0,1]\cap\bigcap_{j=J}^\infty A_j(t).
\]
By countable stability again, it suffices to fix $J\ge1$ and prove that
$\dimH C_J=0$, i.e.\ $\mathcal H^s(C_J)=0$ for every $s>0$.

For $j\ge J$ and $k\in\Z$, let $I_j(k):=[k/Q_j-r_j,k/Q_j+r_j]$ be the closed
interval of center $k/Q_j$ and radius $r_j$ (so $\operatorname{diam}I_j(k)=2r_j$).
For any $x\in C_J$ and $j\ge J$, let $k_j(x):=\mathrm{round}(Q_jx)\in\Z$ be a
nearest integer to $Q_jx$. Since $x\in[0,1]$ and $Q_j\ge1$, we have
$0\le Q_jx\le Q_j$, so $k_j(x)\in\{0,1,\dots,Q_j\}$; moreover, $x\in A_j(t)$
gives $|Q_jx-k_j(x)|=\|Q_jx\|\le t/j$, and dividing by $Q_j$ yields
$|x-k_j(x)/Q_j|\le r_j$, that is, $x\in I_j(k_j(x))$. Furthermore, at two
consecutive levels $j$ and $j+1$, the same point $x\in C_J$ belongs to both
$I_j(k_j(x))$ and $I_{j+1}(k_{j+1}(x))$, so the triangle inequality gives the
parent--child center compatibility condition
\begin{equation}\label{eq:parent-child-tri}
\left|\frac{k_j(x)}{Q_j}-\frac{k_{j+1}(x)}{Q_{j+1}}\right|
\le\left|\frac{k_j(x)}{Q_j}-x\right|+\left|x-\frac{k_{j+1}(x)}{Q_{j+1}}\right|
\le r_j+r_{j+1}=\frac{t}{jQ_j}+\frac{t}{(j+1)Q_{j+1}}.
\end{equation}

Define the surviving center sets $S_n\subseteq\{0,1,\dots,Q_{J+n}\}$ for
$n\ge0$ inductively by setting $S_0:=\{0,1,\dots,Q_J\}$ and, for $n\ge0$,
letting $S_{n+1}$ be the set of all $k'\in\{0,1,\dots,Q_{J+n+1}\}$ for which
there exists $k\in S_n$ satisfying
$|k/Q_{J+n}-k'/Q_{J+n+1}|\le r_{J+n}+r_{J+n+1}$. By induction on $n\ge0$ using
\eqref{eq:parent-child-tri}, every $x\in C_J$ satisfies $k_{J+n}(x)\in S_n$ for
all $n\ge0$, and hence
\[
C_J\subseteq\bigcup_{k\in S_n} I_{J+n}(k)\qquad\text{for all }n\ge0.
\]
By Lemma~\ref{lem:branch}, for each parent $k\in S_n$ the number of children
$k'\in S_{n+1}$ is at most $B=\lfloor16t\rfloor+1$. Since $|S_0|=Q_J+1$,
induction gives $|S_n|\le(Q_J+1)B^n$ for all $n\ge0$. Each interval
$I_{J+n}(k)$ has diameter
\[
\operatorname{diam}I_{J+n}(k)=2r_{J+n}=\frac{2t}{(J+n)Q_{J+n}}
\le\frac{2t}{(J+n)!}\longrightarrow0\qquad(n\to\infty),
\]
where we used $Q_{J+n}\ge(J+n)!$ from Lemma~\ref{lem:factorial} and $J+n\ge1$.

Now fix any $s>0$ and any gauge scale $\delta>0$. For all sufficiently large
$n$, the intervals $\{I_{J+n}(k):k\in S_n\}$ form a finite $\delta$-cover of
$C_J$, and by Lemma~\ref{lem:factorial} ($Q_{J+n}\ge(J+n)!$) their $s$-dimensional
Hausdorff covering cost satisfies
\[
\sum_{k\in S_n}\bigl(\operatorname{diam}I_{J+n}(k)\bigr)^s
\le |S_n|\left(\frac{2t}{(J+n)Q_{J+n}}\right)^s
\le(Q_J+1)(2t)^s B^{-J}\left(\frac{(B^{1/s})^{J+n}}{(J+n)!}\right)^s.
\]
Using the elementary factorial lower bound $m!\ge(m/\ee)^m$ for $m=J+n\ge1$
(or the convergent ratio $(B^{1/s})^{m+1}/(m+1)!\div((B^{1/s})^m/m!)=B^{1/s}/(m+1)<1/2$
for $m+1>2B^{1/s}$, formalized in mathlib as
\texttt{FloorSemiring.tendsto\_pow\_div\_factorial\_atTop}), we have
\[
0\le\frac{(B^{1/s})^m}{m!}\le\left(\frac{\ee B^{1/s}}{m}\right)^m\longrightarrow0
\qquad(m=J+n\to\infty),
\]
and raising to the positive power $s>0$ shows that the covering cost tends to
$0$ as $n\to\infty$. Thus $\mathcal H_\delta^s(C_J)=0$ for every $\delta>0$,
so $\mathcal H^s(C_J)=\lim_{\delta\downarrow0}\mathcal H_\delta^s(C_J)=0$ for
every $s>0$, which gives $\dimH C_J=0$. Taking the countable union over
$J\ge1$ and $m\in\Z$ yields $\dimH E_c=0$. Finally,
$F=\{x\in\R:\Lam(x)<\infty\}=\bigcup_{n=1}^\infty E_n$ is a countable union of
sets of Hausdorff dimension $0$, so $\dimH F=0$.
\end{proof}

\begin{corollary}[The complementary infinite level, Layer~A in Lean]\label{cor:infinite-level}
The set $F=\{x\in\R:\Lam(x)<\infty\}$ has Lebesgue measure zero, and its
complement $F_\infty:=\{x\in\R:\Lam(x)=\infty\}=\R\setminus F$ has full
Lebesgue measure on $\R$. Moreover, for every nondegenerate interval
$[a,b]\subset\R$ ($a<b$), one has
$\operatorname{Leb}(F_\infty\cap[a,b])=b-a>0$ and
\[
\dimH(F_\infty\cap[a,b])=\dimH F_\infty=1.
\]
(Lean, module \texttt{EpiFormal/\allowbreak GaugeExact.lean}: Layer~A, base axioms only.)
\end{corollary}

\begin{proof}
That $\operatorname{Leb}(F)=0$ follows in two independent ways. First, by
Theorem~\ref{thm:finite}, $\dimH F=0<1$, which implies $\mathcal H^1(F)=0$;
since one-dimensional Hausdorff measure $\mathcal H^1$ on $\R$ coincides with
outer Lebesgue measure $\operatorname{Leb}$, we obtain $\operatorname{Leb}(F)=0$.

Second, one can see $\operatorname{Leb}(F)=0$ directly by an elementary
measure-preserving calculation without invoking Hausdorff dimension. For an
integer $m\ge1$ and $j>2m$, the circle endomorphism $T_{Q_j}:\R/\Z\to\R/\Z$,
$x\mapsto Q_jx\bmod1$, preserves normalized Lebesgue measure on $\R/\Z$ (each
arc of length $\ell$ has $Q_j$ disjoint preimage arcs of length $\ell/Q_j$).
The set $\bar A_j(m):=\{x\in\R/\Z:\|Q_jx\|\le m/j\}=T_{Q_j}^{-1}([-m/j,m/j]\bmod1)$
therefore has circle Lebesgue measure $2m/j$. For any fixed starting index
$J\ge1$ and every $j\ge\max(J,2m+1)$, monotonic inclusion gives
\[
\operatorname{Leb}_{\R/\Z}\left(\bigcap_{i=J}^\infty \bar A_i(m)\right)
\le\operatorname{Leb}_{\R/\Z}\bigl(\bar A_j(m)\bigr)=\frac{2m}{j}\longrightarrow0
\qquad(j\to\infty),
\]
so $\operatorname{Leb}_{\R/\Z}(\bigcap_{i\ge J}\bar A_i(m))=0$. Every point
$x\in\R/\Z$ with $\Lam(x)<\infty$ belongs to $\bigcap_{i\ge J}\bar A_i(m)$ for
some integers $m>\Lam(x)$ and $J\ge1$. Countable subadditivity over $m\ge1$ and
$J\ge1$ yields $\operatorname{Leb}_{\R/\Z}(\{\Lam<\infty\})=0$ on the circle,
and countable union over integer translates $[k,k+1]$ ($k\in\Z$) yields
$\operatorname{Leb}(F)=0$ on $\R$.

Now let $[a,b]\subset\R$ be any nondegenerate closed interval ($a<b$). Since
$[a,b]=(F_\infty\cap[a,b])\cup(F\cap[a,b])$ and $\operatorname{Leb}(F\cap[a,b])\le\operatorname{Leb}(F)=0$,
additivity of Lebesgue measure gives
\[
\mathcal H^1(F_\infty\cap[a,b])=\operatorname{Leb}(F_\infty\cap[a,b])
=\operatorname{Leb}([a,b])-\operatorname{Leb}(F\cap[a,b])=b-a>0.
\]
By the definition of Hausdorff dimension, $\mathcal H^1(F_\infty\cap[a,b])>0$
implies $\dimH(F_\infty\cap[a,b])\ge1$. On the other hand, any subset of $\R$
has Hausdorff dimension at most $\dimH\R=1$, so
\[
1\le\dimH(F_\infty\cap[a,b])\le\dimH F_\infty\le\dimH\R=1.
\]
Hence $\dimH(F_\infty\cap[a,b])=\dimH F_\infty=1$. Together with
Theorem~\ref{thm:finite}, this establishes the sharp metric dichotomy: $\Lam$
is infinite Lebesgue-almost everywhere (with local Hausdorff dimension~$1$ on
every interval), while the set where $\Lam$ is finite has Hausdorff
dimension~$0$.
\end{proof}

\paragraph{Lean status of the assembly.} All steps of Theorem~\ref{thm:finite}
are kernel-checked at Layer~A. The surviving centers are encoded as a
\texttt{Finset} recursion (\texttt{survSet}), not as paths:
$S_0=\{0,\dots,Q_J\}$ and $S_{n+1}=\bigcup_{k\in S_n}\mathrm{branchSet}(k)$,
with $|S_n|\le(Q_J+1)B^n$ by \texttt{Finset.card\_biUnion\_le\_card\_mul} and
Lemma~\ref{lem:branch} (\texttt{card\_survSet\_le}). Membership of the center
sequence $\mathrm{round}(Q_{J+n}x)$ in $S_n$ is proved by induction via the
triangle inequality (\texttt{centerOf\_mem\_survSet}), giving the cover
(\texttt{tailSet\_cover}). The content bound
$|S_n|(2t/((J+n)Q_{J+n}))^s\to0$ uses $Q_j\ge j!$ and
\texttt{FloorSemiring.\allowbreak tendsto\_pow\_div\_factorial\_atTop} with
$c=B^{1/s}$ (\texttt{tendsto\_contentBound}); no asymptotic notation is used.
Feeding this into \texttt{hausdorffMeasure\_le\_liminf\_sum} yields
$\mu H^s=0$ for all $s>0$ (\texttt{hausdorffMeasure\_tailSet\_zero}), hence
$\dimH=0$ (\texttt{dimH\_tailSet\_zero}). Periodicity
$\|Q_j(x+m)\|=\|Q_jx\|$ (\texttt{apprSeq\_add\_intCast}) reduces general $x$ to
$[0,1]$; integer shifts are $1$-Lipschitz and countable unions preserve
dimension $0$ (\texttt{dimH\_iUnion}). The sub-threshold case $c<1/4$ is
\emph{not} used: a monotonicity reduction to $c'=\max(c,0)$ lets the covering
argument handle all $c$ uniformly, so the theorem needs no continued-fraction
axioms.

\section{Subcritical weights and exact Hausdorff dimensions}\label{sec:weights}

\begin{remark}\label{rem:unweighted}
The inequality $\|Q_jx\|\le j\|Q_jx\|$ gives the inclusion from the
weighted zero-set $\{x:j\|Q_jx\|\to0\}=L$ into the unweighted zero-set
$\{x:\|Q_jx\|\to0\}$. To see strictness explicitly, the point $x_0$ constructed
in \texttt{QuarterPointFinal.lean} (with constant digits $d_j\equiv1$) satisfies
$j\|Q_jx_0\|\to1/4$ (\texttt{QuarterPointFinal.\allowbreak lambda\_x0\_eq\_quarter}).
Its unweighted error $\|Q_jx_0\|$ tends to zero, while $x_0\notin L$ by the
weighted classification (Theorem~\ref{thm:class}). The combined witness is
kernel-checked at Layer~A in \texttt{EpiFormal/\allowbreak UnweightedCorrection.lean}
(\texttt{UnweightedCorrection.\allowbreak witness}) using only the three Lean
base axioms.
\end{remark}

Recall that for $0\le\alpha\le1$ we write
$\Lam_\alpha(x)=\limsup_{j\to\infty}j^\alpha\|Q_jx\|$, so that $\Lam_1=\Lam$.

\begin{theorem}[Weighted upper bound, Layer~A in Lean]\label{thm:upper}
For $0\le\alpha\le1$ and every finite $c\in\R$,
$\dimH\{x\in\R:\Lam_\alpha(x)\le c\}\le1-\alpha$.
(Lean: \texttt{WeightedFinal.\allowbreak dimH\_levelSet\ensuremath{\alpha}\_le};
\texttt{\#print axioms} returns \texttt{[propext, Classical.choice, Quot.sound]} only.)
\end{theorem}

\begin{proof}
Fix $0\le\alpha\le1$ and $c\in\R$, and choose $t>\max(c,0)$ (in Lean,
$t=\max(c,0)+1>0$). Put $E_{\alpha,c}:=\{x\in\R:\Lam_\alpha(x)\le c\}$. Since
$Q_j\in\Z$, we have $\|Q_j(x+m)\|=\|Q_jx\|$ for every $m\in\Z$ and $j\ge1$,
hence $\Lam_\alpha(x+m)=\Lam_\alpha(x)$. Thus
\[
E_{\alpha,c}=\bigcup_{m\in\Z}\bigl(m+(E_{\alpha,c}\cap[0,1])\bigr).
\]
By isometry invariance and countable stability of Hausdorff dimension, it
suffices to show that $\dimH(E_{\alpha,c}\cap[0,1])\le1-\alpha$.

For $j\ge1$, set $r_{j,\alpha}:=t/(j^\alpha Q_j)$ and
$A_{j,\alpha}(t):=\{x\in\R:\|Q_jx\|\le t/j^\alpha\}$. Any $x\in E_{\alpha,c}\cap[0,1]$
satisfies $\limsup_{j\to\infty}j^\alpha\|Q_jx\|\le c<t$, so $j^\alpha\|Q_jx\|\le t$
eventually; hence
\[
E_{\alpha,c}\cap[0,1]\subseteq\bigcup_{J=1}^\infty C_{J,\alpha},\qquad
C_{J,\alpha}:=[0,1]\cap\bigcap_{j=J}^\infty A_{j,\alpha}(t).
\]
Fix $J\ge1$. For each $j\ge J$ and $k\in\Z$, let
$I_{j,\alpha}(k):=[k/Q_j-r_{j,\alpha},k/Q_j+r_{j,\alpha}]$ be the closed
interval of center $k/Q_j$ and radius $r_{j,\alpha}$. For any $x\in C_{J,\alpha}$
and $j\ge J$, the nearest integer $k_j(x):=\mathrm{round}(Q_jx)$ belongs to
$\{0,1,\dots,Q_j\}$ and satisfies $|x-k_j(x)/Q_j|=\|Q_jx\|/Q_j\le r_{j,\alpha}$,
so $x\in I_{j,\alpha}(k_j(x))$. At consecutive levels $j$ and $j+1$, the
triangle inequality gives
\begin{equation}\label{eq:weighted-parent-child}
\left|\frac{k_j(x)}{Q_j}-\frac{k_{j+1}(x)}{Q_{j+1}}\right|
\le\left|\frac{k_j(x)}{Q_j}-x\right|+\left|x-\frac{k_{j+1}(x)}{Q_{j+1}}\right|
\le r_{j,\alpha}+r_{j+1,\alpha}=\frac{t}{j^\alpha Q_j}+\frac{t}{(j+1)^\alpha Q_{j+1}}.
\end{equation}
Since $0\le\alpha\le1$ and $j\ge1$, we have $1-\alpha\ge0$ and
$j=j^\alpha j^{1-\alpha}\le j^\alpha(j+1)^{1-\alpha}$, so
$1/j^\alpha\le(j+1)^{1-\alpha}/j$, while
$1/(j+1)^\alpha=(j+1)^{1-\alpha}/(j+1)$. Putting $t'_j:=t(j+1)^{1-\alpha}\ge0$,
we obtain the pointwise radius comparison (\texttt{WeightedFiber.wrad\_le\_rad})
\[
r_{j,\alpha}+r_{j+1,\alpha}
\le\frac{t'_j}{jQ_j}+\frac{t'_j}{(j+1)Q_{j+1}}.
\]
Applying Lemma~\ref{lem:branch} with threshold $t'_j$ shows that for each fixed
parent center $k\in\{0,1,\dots,Q_j\}$, the number of child centers
$k'\in\{0,1,\dots,Q_{j+1}\}$ satisfying \eqref{eq:weighted-parent-child} is at
most
\[
\lfloor16t'_j\rfloor+1\le 16t(j+1)^{1-\alpha}+1\le K(j+1)^{1-\alpha},
\qquad K:=16t+1,
\]
since $(j+1)^{1-\alpha}\ge1$.

Define the surviving center sets $S_{n,\alpha}\subseteq\{0,1,\dots,Q_{J+n}\}$
($n\ge0$) inductively by $S_{0,\alpha}:=\{0,1,\dots,Q_J\}$ and letting
$S_{n+1,\alpha}$ consist of all $k'\in\{0,1,\dots,Q_{J+n+1}\}$ within
$r_{J+n,\alpha}+r_{J+n+1,\alpha}$ of some $k/Q_{J+n}$ with $k\in S_{n,\alpha}$.
By induction on $n\ge0$, $C_{J,\alpha}\subseteq\bigcup_{k\in S_{n,\alpha}}I_{J+n,\alpha}(k)$
and the cardinality $|S_{n,\alpha}|$ satisfies
(\texttt{WeightedDimH.card\_wsurvSet\_le})
\[
|S_{n,\alpha}|
\le(Q_J+1)\prod_{i=0}^{n-1}\bigl(K(J+i+1)^{1-\alpha}\bigr)
\le(Q_J+1)K^n\bigl((J+n)!\bigr)^{1-\alpha},
\]
because $\prod_{i=0}^{n-1}(J+i+1)=(J+n)!/J!\le(J+n)!$. By
Lemma~\ref{lem:factorial}, $Q_{J+n}\ge(J+n)!$, and $(J+n)^\alpha\ge1$, so each
interval $I_{J+n,\alpha}(k)$ has diameter
\[
\operatorname{diam}I_{J+n,\alpha}(k)=2r_{J+n,\alpha}
=\frac{2t}{(J+n)^\alpha Q_{J+n}}\le\frac{2t}{(J+n)!}\longrightarrow0
\qquad(n\to\infty).
\]
Now fix any exponent $s>1-\alpha$, and put $\beta:=s-(1-\alpha)>0$. For every
$\delta>0$ and all sufficiently large $n$, $\{I_{J+n,\alpha}(k):k\in S_{n,\alpha}\}$
is a finite $\delta$-cover of $C_{J,\alpha}$ whose $s$-dimensional Hausdorff
covering cost satisfies
\[
\sum_{k\in S_{n,\alpha}}\bigl(\operatorname{diam}I_{J+n,\alpha}(k)\bigr)^s
\le |S_{n,\alpha}|\left(\frac{2t}{(J+n)!}\right)^s
\le(Q_J+1)(2t)^s K^{-J}\cdot K^{J+n}\bigl((J+n)!\bigr)^{1-\alpha-s}.
\]
With $m=J+n$ and $K^{1/\beta}>0$, the factor $K^m(m!)^{1-\alpha-s}=K^m(m!)^{-\beta}=\bigl((K^{1/\beta})^m/m!\bigr)^\beta$
satisfies, via the elementary factorial lower bound $m!\ge(m/\ee)^m$,
\[
0\le K^m(m!)^{-\beta}=\left(\frac{(K^{1/\beta})^m}{m!}\right)^\beta
\le\left(\frac{\ee K^{1/\beta}}{m}\right)^{\beta m}\longrightarrow0
\qquad(m=J+n\to\infty)
\]
(formalized in \texttt{WeightedPhase.\allowbreak content\_\allowbreak core\_\allowbreak tendsto} via
\texttt{FloorSemiring.\allowbreak tendsto\_\allowbreak pow\_\allowbreak div\_\allowbreak factorial\_\allowbreak atTop} and continuity of
$u\mapsto u^\beta$ at $0$). Hence $\mathcal H_\delta^s(C_{J,\alpha})=0$ for all
$\delta>0$, so $\mathcal H^s(C_{J,\alpha})=0$ and $\dimH C_{J,\alpha}\le s$.
Since $s>1-\alpha$ was arbitrary, letting $s\downarrow1-\alpha$ gives
$\dimH C_{J,\alpha}\le1-\alpha$ for every $J\ge1$, and countable stability over
$J\ge1$ and $m\in\Z$ yields $\dimH E_{\alpha,c}\le1-\alpha$.
\end{proof}

\paragraph{Lean status.} Theorem~\ref{thm:upper} is kernel-checked at Layer~A
as \texttt{WeightedFinal.\allowbreak dimH\_levelSet\ensuremath{\alpha}\_le},
assembled through \texttt{WeightedFiber.lean}, \texttt{WeightedDimHCore.lean},
\texttt{WeightedUpper.lean}, and \texttt{WeightedFinal.lean}. Its axiom audit
reports only the three base axioms.

\begin{theorem}[Prescribed limits at subcritical weights, Layer~A in Lean]\label{conj:weight}
Let $0\le\alpha<1$. For $\alpha>0$ take any $c\ge0$; for $\alpha=0$
take $0\le c\le1/2$. Then
\[
 \dimH\{x\in[0,1]: j^\alpha\|Q_jx\|\longrightarrow c\}=1-\alpha .
\]
In particular $\dimH\{x\in\R:\Lam_\alpha(x)\le c\}=1-\alpha$ for every
finite $c\ge0$, including $c=0$.
(Lean, module \texttt{WeightedExactDim}:
\texttt{dimH\_\allowbreak exactLimitSet\ensuremath{\alpha}\_eq},
\texttt{dimH\_\allowbreak levelSet\ensuremath{\alpha}\_eq}, and
\texttt{dimH\_\allowbreak apprExponent\ensuremath{\alpha}\_eq\_\allowbreak zero};
\texttt{\#print axioms} returns \texttt{[propext, Classical.choice, Quot.sound]} only.)
\end{theorem}

\begin{proof}
Write $E_{\alpha}^{\mathrm{lim}}(c):=\{x\in[0,1]:j^\alpha\|Q_jx\|\to c\}$. Since
$j^\alpha\|Q_jx\|\to c$ implies $\Lam_\alpha(x)=\limsup_{j\to\infty}j^\alpha\|Q_jx\|=c$,
we have $E_{\alpha}^{\mathrm{lim}}(c)\subseteq\{x\in\R:\Lam_\alpha(x)\le c\}$, so
Theorem~\ref{thm:upper} immediately gives the upper bound
$\dimH E_{\alpha}^{\mathrm{lim}}(c)\le1-\alpha$.

We now give the complete paper proof of the lower bound
$\dimH E_{\alpha}^{\mathrm{lim}}(c)\ge1-\alpha$. For $j\ge1$, put
\[
w_j:=\frac{j^{-\alpha}}{\log(j+2)},\qquad r_j:=\frac{w_j}{Q_j}.
\]
Choose an integer $N\ge3$ large enough that $w_j<1/8$ for all $j\ge N$, and
also $cj^{-\alpha}+w_j\le1/2$ for all $j\ge N$ whenever either ($\alpha>0$,
$c\ge0$) or ($\alpha=0$, $0\le c<1/2$). Define the target fractional shift
$\theta_j\in[w_j,1/2-w_j]$ for $j\ge N$ by clamping $cj^{-\alpha}$ away from the
half-integer boundary:
\[
\theta_j:=\begin{cases}
w_j,&\text{if }c=0,\\[2pt]
cj^{-\alpha},&\text{if }c>0\text{ and }(\alpha>0\text{ or }\alpha=0,c<1/2),\\[2pt]
\tfrac12-w_j,&\text{if }\alpha=0\text{ and }c=\tfrac12.
\end{cases}
\]
(Note why the boundary adjustment $\theta_j=1/2-w_j$ is needed in the endpoint
case $\alpha=0,c=1/2$, and $\theta_j=w_j$ when $c=0$: if one centered directly
at $(k+1/2)/Q_j$, points on the right half of the interval would cross the
half-integer $k+1/2$ where $y\mapsto\|y\|$ folds back, whereas with center
$(k+\theta_j)/Q_j$ and radius $r_j=w_j/Q_j$, every point $x$ in the closed
interval $I_j(k):=[(k+\theta_j-w_j)/Q_j,(k+\theta_j+w_j)/Q_j]$ satisfies
$k\le k+\theta_j-w_j\le Q_jx\le k+\theta_j+w_j\le k+1/2$. On $[k,k+1/2]$ the
distance to the nearest integer is given without folding by
$\|Q_jx\|=Q_jx-k\in[\theta_j-w_j,\theta_j+w_j]$.) In particular, in the case
$\alpha=0,c=1/2$, for $x\in I_j(k)$ we have
$Q_jx\in[k+1/2-2w_j,k+1/2]$, so
\[
\|Q_jx\|=\frac12-\left(k+\frac12-Q_jx\right)\in\left[\frac12-2w_j,\frac12\right]
\quad\Longrightarrow\quad
\bigl|\|Q_jx\|-\tfrac12\bigr|\le2w_j=\frac{2}{\log(j+2)}\longrightarrow0,
\]
and in all three cases every $x\in I_j(k)$ satisfies
\[
\bigl|j^\alpha\|Q_jx\|-c\bigr|
\le j^\alpha\bigl|\|Q_jx\|-\theta_j\bigr|+\bigl|j^\alpha\theta_j-c\bigr|
\le j^\alpha w_j+j^\alpha w_j=\frac{2}{\log(j+2)}\longrightarrow0
\qquad(j\to\infty).
\]
Moreover, for fixed $j\ge N$ and distinct integers $k_1\ne k_2$, the distance
between the centers of $I_j(k_1)$ and $I_j(k_2)$ is
$|k_1-k_2|/Q_j\ge1/Q_j>2w_j/Q_j=2r_j$ (since $w_j<1/8<1/2$), so the level-$j$
intervals $\{I_j(k):k\in\Z\}$ are pairwise disjoint.

Next we bound the number of strictly nested level-$(j+1)$ intervals inside a
given parent interval $I_j(k)$ ($j\ge N$). A child interval
$I_{j+1}(k')=[c_{j+1}(k')-r_{j+1},c_{j+1}(k')+r_{j+1}]$ with center
$c_{j+1}(k')=(k'+\theta_{j+1})/Q_{j+1}$ is contained in
$I_j(k)=[c_j(k)-r_j,c_j(k)+r_j]$ ($c_j(k)=(k+\theta_j)/Q_j$) if and only if
$|c_{j+1}(k')-c_j(k)|\le r_j-r_{j+1}$, which is equivalent to requiring the
integer $k'$ to lie in the closed interval $[A_j(k),B_j(k)]$ where
\[
A_j(k):=Q_{j+1}\bigl(c_j(k)-(r_j-r_{j+1})\bigr)-\theta_{j+1},\qquad
B_j(k):=Q_{j+1}\bigl(c_j(k)+(r_j-r_{j+1})\bigr)-\theta_{j+1}.
\]
The length of $[A_j(k),B_j(k)]$ is independent of the parent index $k$ and of
the shifts $\theta_j,\theta_{j+1}$:
\[
W_j:=B_j(k)-A_j(k)=2(r_j-r_{j+1})Q_{j+1}
=2\left(w_j\frac{Q_{j+1}}{Q_j}-w_{j+1}\right).
\]
Since any closed real interval $[A,B]$ of length $W=B-A\ge1$ contains at least
$\lfloor B\rfloor-\lceil A\rceil+1>W-1$ integers, the number of eligible child
centers $k'\in\Z$ with $I_{j+1}(k')\subseteq I_j(k)$ is at least
$2(r_j-r_{j+1})Q_{j+1}-1$. Using $Q_{j+1}=(4j+2)Q_j+Q_{j-1}>(4j+2)Q_j$ from
\eqref{eq:rec} and $w_{j+1}\le w_j$, we have
\[
2(r_j-r_{j+1})Q_{j+1}-1
=2w_j\frac{Q_{j+1}}{Q_j}-2w_{j+1}-1
\ge 2w_j(4j+2)-2w_j-1
=\frac{8j^{1-\alpha}+2j^{-\alpha}}{\log(j+2)}-1.
\]
Since $0\le\alpha<1$, $j^{1-\alpha}/\log(j+2)\to\infty$ as $j\to\infty$.
Increasing $N$ if necessary so that $7j^{1-\alpha}/\log(j+2)\ge1$ and
$b_j:=\lfloor j^{1-\alpha}/\log(j+2)\rfloor\ge2$ for all $j\ge N$, we obtain
\[
2(r_j-r_{j+1})Q_{j+1}-1\ge\frac{j^{1-\alpha}}{\log(j+2)}\ge b_j\ge2
\qquad\text{for all }j\ge N,
\]
so every parent interval $I_j(k)$ ($j\ge N$) contains at least $b_j$ pairwise
disjoint child intervals $I_{j+1}(k')$.

Choose $N$ large enough that $I_N(1)\subseteq(0,1)$. Construct a homogeneous
Cantor tree by starting at level $N$ with the single root interval $I_N(1)$
($\mathcal E_N:=\{1\}$), and for each $j\ge N$ and each retained parent
$k\in\mathcal E_j$, keeping exactly $b_j$ child indices $k'\in\mathcal E_{j+1}$
with $I_{j+1}(k')\subseteq I_j(k)$ (for instance, the first $b_j$ integers
$\lceil A_j(k)\rceil+m$, $0\le m<b_j$). Then at each level $n\ge N$ there are
exactly
\[
M_n:=\prod_{j=N}^{n-1}b_j
\]
pairwise disjoint closed intervals $\{I_n(k):k\in\mathcal E_n\}$ of radius
$r_n=w_n/Q_n$, and their compact intersection
\[
K:=\bigcap_{n=N}^\infty\bigcup_{k\in\mathcal E_n}I_n(k)\subseteq[0,1]
\]
is nonempty and satisfies $K\subseteq E_{\alpha}^{\mathrm{lim}}(c)$. Let $\mu$
be the unique Borel probability measure on $K$ assigning equal mass
$\mu(I_n(k))=M_n^{-1}$ to each of the $M_n$ retained level-$n$ intervals ($n\ge N$).

To estimate the asymptotic growth of $\log M_n$ and $\log Q_{n+1}$, observe
that for $j\ge N$,
\[
\frac{j^{1-\alpha}}{2\log(j+2)}\le\frac{j^{1-\alpha}}{\log(j+2)}-1
<b_j\le\frac{j^{1-\alpha}}{\log(j+2)},
\]
so $\log b_j=(1-\alpha)\log j-\log\log(j+2)+O(1)$. By integral comparison for
the monotone function $t\mapsto\log t$,
$\sum_{j=N}^{n-1}\log j=\int_1^n\log t\,dt+O(\log n)=n\log n-n+O(\log n)$, and
$\sum_{j=N}^{n-1}\log\log(j+2)\le n\log\log(n+2)$, which gives
\begin{equation}\label{eq:log-Mn-asymp}
\log M_n=\sum_{j=N}^{n-1}\log b_j
=(1-\alpha)n\log n+O(n+n\log\log n)=(1-\alpha)n\log n+o(n\log n).
\end{equation}
Similarly, from $4jQ_j<(4j+2)Q_j<Q_{j+1}\le(4j+3)Q_j$ for all $j\ge1$, taking
logarithms and summing from $j=1$ to $n$ gives
\[
\log Q_{n+1}=\sum_{j=1}^n\log\frac{Q_{j+1}}{Q_j}
=\sum_{j=1}^n\bigl(\log j+\log 4+O(1/j)\bigr)=n\log n+O(n).
\]

Now fix any $0<s<1-\alpha$. By \eqref{eq:log-Mn-asymp} and
$\log Q_{n+1}=n\log n+O(n)$, we have
\begin{align*}
\log\left(\frac{3}{M_n Q_{n+1}^{-s}}\right)
&=s\log Q_{n+1}-\log M_n+\log 3\\
&=-\bigl((1-\alpha)-s\bigr)n\log n+o(n\log n)\longrightarrow-\infty
\qquad(n\to\infty).
\end{align*}
Hence we may choose an index $N_s\ge N$ such that
\begin{equation}\label{eq:mass-scale-bound}
\frac{3}{M_n}\le Q_{n+1}^{-s}\qquad\text{for all }n\ge N_s.
\end{equation}
Let $\{U_i\}_{i\ge1}$ be any countable cover of $K$ by nonempty sets of diameter
$0<h_i:=\operatorname{diam}(U_i)<Q_{N_s}^{-1}$ (if $\operatorname{diam}(U_i)=0$,
then $U_i$ is a singleton and $\mu(U_i)=0$ since $M_n^{-1}\to0$). For each $i$,
since $Q_n$ is strictly increasing to $\infty$, there exists a unique integer
$n_i\ge N_s$ such that
\[
Q_{n_i+1}^{-1}\le h_i<Q_{n_i}^{-1}.
\]
The set $U_i$ is contained in a closed interval $[x_0,x_0+h_i]$ of length
$h_i<Q_{n_i}^{-1}$. If a level-$n_i$ interval $I_{n_i}(k)$ ($k\in\Z$) intersects
$[x_0,x_0+h_i]$, then its center $c_{n_i}(k)=(k+\theta_{n_i})/Q_{n_i}$ must lie
in $[x_0-r_{n_i},x_0+h_i+r_{n_i}]$, which has length
$h_i+2r_{n_i}<Q_{n_i}^{-1}+2w_{n_i}Q_{n_i}^{-1}<(1+1/4)Q_{n_i}^{-1}<2Q_{n_i}^{-1}$.
Since the centers $c_{n_i}(k)$ ($k\in\Z$) form a grid of spacing exactly
$Q_{n_i}^{-1}$, an interval of length $<2Q_{n_i}^{-1}$ can contain at most $3$
grid points $c_{n_i}(k)$ (if it contained $4$ points $k_1<k_2<k_3<k_4$, its
length would be at least $(k_4-k_1)Q_{n_i}^{-1}\ge3Q_{n_i}^{-1}$). Thus $U_i$
intersects at most $3$ retained level-$n_i$ intervals $I_{n_i}(k)$
($k\in\mathcal E_{n_i}$). Because $\mu$ is supported on
$\bigcup_{k\in\mathcal E_{n_i}}I_{n_i}(k)$ and each $I_{n_i}(k)$ has measure
$M_{n_i}^{-1}$, it follows from \eqref{eq:mass-scale-bound} and
$Q_{n_i+1}^{-1}\le h_i$ that
\[
\mu(U_i)\le\frac{3}{M_{n_i}}\le Q_{n_i+1}^{-s}\le h_i^s=\bigl(\operatorname{diam}(U_i)\bigr)^s.
\]
Summing over the cover $\{U_i\}$ yields
\[
\sum_{i}\bigl(\operatorname{diam}(U_i)\bigr)^s
\ge\sum_{i}\frac{3}{M_{n_i}}\ge\sum_{i}\mu(U_i)\ge\mu(K)=1.
\]
Taking the infimum over all countable $\delta$-covers with
$0<\delta<Q_{N_s}^{-1}$ and letting $\delta\downarrow0$ gives
$\mathcal H^s(E_{\alpha}^{\mathrm{lim}}(c))\ge\mathcal H^s(K)\ge1>0$, and therefore
$\dimH E_{\alpha}^{\mathrm{lim}}(c)\ge s$. Letting $s\uparrow1-\alpha$ proves
$\dimH E_{\alpha}^{\mathrm{lim}}(c)\ge1-\alpha$, establishing the exact identity
$\dimH E_{\alpha}^{\mathrm{lim}}(c)=1-\alpha$.

Finally, for any finite $c\ge0$ (including $c=0$), we have the inclusions
\[
E_{\alpha}^{\mathrm{lim}}(0)
\subseteq\{x\in\R:\Lam_\alpha(x)=0\}
\subseteq\{x\in\R:\Lam_\alpha(x)\le c\}.
\]
The lower bound $\dimH E_{\alpha}^{\mathrm{lim}}(0)=1-\alpha$ just proved and
the upper bound $\dimH\{x\in\R:\Lam_\alpha(x)\le c\}\le1-\alpha$ from
Theorem~\ref{thm:upper} squeeze both sets to have exact Hausdorff dimension
$1-\alpha$.

The Lean formalization in \texttt{EpiFormal/\allowbreak WeightedExactDim.lean}
certifies the same lower bound at Layer~A for every $0<s<1-\alpha$ by taking
polynomial window widths $w_j=j^{-\alpha-\varepsilon}$ with
$\varepsilon=(1-\alpha-s)/2>0$, clamping the target shift into
$[w_j,1/2-w_j]$, and applying compactness plus finite-tree counting directly to
the definition of Hausdorff measure to establish
$\mathcal H^s(K)\ge1$ (\texttt{one\_le\_hausdorffMeasure\_cantorSet}).
\end{proof}

For $\alpha=0$ the dimension assertion also follows from Corollary~8 of
\cite{BLR} with constant target $c$, by the distance inequality used
above. For $0<\alpha<1$, condition~(1.6) of their Theorem~7 fails because
$\log Q_{j+1}/\log Q_j\to1$. The lower bound in the paper proof above can also
be read from the homogeneous Moran bounds of Feng, Wen and Wu \cite{FWW}: the
retained level has $\log M_n\sim(1-\alpha)n\log n$ and both bounding logarithmic
length scales are asymptotic to $n\log n$. The independent covering upper bound
(Theorem~\ref{thm:upper}) applies to the full prescribed-limit set and sublevel
set.

\section{Critical levels, exact Hausdorff gauges, and the positive-level trichotomy}\label{sec:gauge}

\subsection{Critical levels and exact gauges}

For an integer $k\ge1$, set $b=2k+1$. Define an increasing continuous
gauge $h_k$ by $h_k(0)=0$, $h_k(1/Q_n)=b^{-n}$ ($n\ge1$), and
linear interpolation on each intervening interval $[1/Q_{n+1},1/Q_n]$ (with
$h_k(t)=h_k(1/Q_1)$ for $t\ge1/Q_1$). Only its values near zero matter. Write
$\mathcal H^{h_k}$ for Hausdorff measure with covering cost
$h_k(\operatorname{diam} U)$.

\begin{theorem}[Size of a critical spectral level, Layer~A in Lean]\label{thm:gauge}
On the circle $\R/\Z$, the exact level $\{x:\Lam(x)=k/4\}$ has
infinite but $\sigma$-finite $\mathcal H^{h_k}$ measure.
The strict sublevel $\{x:\Lam(x)<k/4\}$ has zero $\mathcal H^{h_k}$ measure.
In the family of gauges $g_\beta$ given by $g_\beta(1/Q_n)=\beta^{-n}$
and the same linear interpolation ($\beta>1$), the exact level has zero
$\mathcal H^{g_\beta}$ measure for $\beta>2k+1$, whereas for $1<\beta<2k+1$ its
$\mathcal H^{g_\beta}$ measure is not $\sigma$-finite.
\end{theorem}

\begin{proof}
\emph{Digits.} Use $M_j$, $u_j=Q_jx-M_j$, $A_j=4j+2$ and
$d_j=M_{j+1}-A_jM_j-M_{j-1}=A_ju_j+u_{j-1}-u_{j+1}$ as in
Section~\ref{sec:class}, with a fixed Borel rule at ties (for instance,
$M_j=\lfloor Q_jx+1/2\rfloor$). Ties occur only on the
countable set $\bigcup_{j\ge1}\{x:Q_jx\in\Z+\tfrac12\}$, which is null for every
non-atomic measure and is avoided by the Cantor construction below.
If $\Lam(x)\le k/4$, eventually $|u_j|\le(k+1)/j$, so
$d_j-4ju_j=2u_j+u_{j-1}-u_{j+1}\to0$ and
$\limsup_{j\to\infty}|d_j|=4\Lam(x)\le k$. Since $d_j\in\Z$, the integer digits
eventually satisfy $|d_j|\le k$. For the \emph{strict} sublevel $\Lam(x)<k/4$,
Corollary~\ref{thm:sub} improves this to $|d_j|\le k-1$ eventually.

\emph{Upper bounds.} Partition the non-strict sublevel $\{\Lam\le k/4\}$ on
$[0,1]$ (or $\R/\Z$) into countably many classes
$\mathcal C(N,m_{N-1},m_N)$ indexed by a common starting index $N\ge2$ and the
two starting integers $(M_{N-1},M_N)=(m_{N-1},m_N)\in\Z^2$, requiring
$|u_j|\le(k+1)/j$ and $|d_j|\le k$ for all $j\ge N$. With a fixed Borel rounding
rule these classes are Borel sets. Because $M_{j+1}=A_jM_j+M_{j-1}+d_j$ is
determined inductively from $(M_{N-1},M_N)$ by the $n-N$ digits
$d_N,\dots,d_{n-1}\in\{-k,\dots,k\}$, by level $n>N$ a class has at most
$b^{\,n-N}$ possible values of $M_n$ ($b=2k+1$). Each value of $M_n$ confines
$x$ to the interval $[M_n/Q_n-(k+1)/(nQ_n),M_n/Q_n+(k+1)/(nQ_n)]$ of diameter
$2(k+1)/(nQ_n)\le1/Q_n$ for all $n\ge2(k+1)$ (or directly inside
$[M_n/Q_n-1/(2Q_n),M_n/Q_n+1/(2Q_n)]$ of diameter $1/Q_n$ using $|u_n|\le1/2$).
Since $1/Q_n\to0$, for every $\delta>0$ and all large $n$ this gives a finite
$\delta$-cover of $\mathcal C(N,m_{N-1},m_N)$ with $\mathcal H_\delta^{h_k}$
cost at most
\[
b^{\,n-N}h_k(1/Q_n)=b^{\,n-N}b^{-n}=b^{-N}<\infty.
\]
Letting $\delta\downarrow0$ gives
$\mathcal H^{h_k}(\mathcal C(N,m_{N-1},m_N))\le b^{-N}<\infty$. For the strict
sublevel $\{\Lam<k/4\}$, the allowed tail digits satisfy
$d_j\in\{-(k-1),\dots,k-1\}$, so the count at level $n$ uses at most
$(2k-1)^{n-N}$ values of $M_n$ with $2k-1<b=2k+1$. The $\mathcal H_\delta^{h_k}$
covering cost of each class is therefore at most
\[
(2k-1)^{n-N}h_k(1/Q_n)=b^{-N}\left(\frac{2k-1}{2k+1}\right)^{n-N}\longrightarrow0
\qquad(n\to\infty),
\]
so every class is $\mathcal H^{h_k}$-null, and countable subadditivity yields
\begin{equation}\label{eq:strict-null}
 \mathcal H^{h_k}\bigl(\{x:\Lam(x)<k/4\}\bigr)=0 .
\end{equation}

\emph{Positivity.} Fix a large integer $N\ge2$ (chosen below), set
$M_{N-1}=M_N=0$, and prescribe arbitrary digits
$\omega=(d_N,d_{N+1},\dots)\in\{-k,\ldots,k\}^{\{j\ge N\}}$ via the linear
recurrence
\[
M_{j+1}=A_jM_j+M_{j-1}+d_j\qquad(j\ge N).
\]
Define the determinant sequence $W_j:=M_{j+1}Q_j-M_jQ_{j+1}$ for $j\ge N-1$, so
that $W_{N-1}=M_NQ_{N-1}-M_{N-1}Q_N=0$. Using $M_{j+1}=A_jM_j+M_{j-1}+d_j$ and
$Q_{j+1}=A_jQ_j+Q_{j-1}$ from \eqref{eq:rec}, the $A_jM_jQ_j$ terms cancel and
we obtain the first-order recurrence
\[
W_j=(A_jM_j+M_{j-1}+d_j)Q_j-M_j(A_jQ_j+Q_{j-1})
=d_jQ_j-(M_jQ_{j-1}-M_{j-1}Q_j)=d_jQ_j-W_{j-1}
\]
for all $j\ge N$. We claim by induction on $j\ge N-1$ that
\begin{equation}\label{eq:Wj-bound}
|W_j|\le 2kQ_j\qquad\text{for all }j\ge N-1.
\end{equation}
Indeed, for $j=N-1$, $|W_{N-1}|=0\le2kQ_{N-1}$. For $j\ge N$, since $Q_j\ge2Q_{j-1}$
(in fact $Q_j\ge(4j-2)Q_{j-1}\ge6Q_{j-1}$ for $j\ge2$), the inductive
hypothesis $|W_{j-1}|\le2kQ_{j-1}\le kQ_j$ and $|d_j|\le k$ yield
\[
|W_j|\le|d_j|Q_j+|W_{j-1}|\le kQ_j+2kQ_{j-1}\le kQ_j+kQ_j=2kQ_j,
\]
proving \eqref{eq:Wj-bound}. Dividing $W_j$ by $Q_jQ_{j+1}$ gives the adjacent
rational difference
\[
\frac{M_{j+1}}{Q_{j+1}}-\frac{M_j}{Q_j}=\frac{W_j}{Q_jQ_{j+1}},\qquad
\left|\frac{M_{j+1}}{Q_{j+1}}-\frac{M_j}{Q_j}\right|\le\frac{2k}{Q_{j+1}}.
\]
Because $Q_{j+2}\ge(4j+6)Q_{j+1}\ge2Q_{j+1}$ for all $j\ge1$, we have
$Q_{m+1+r}\ge2^r Q_{m+1}$ for all $m\ge N$ and $r\ge0$. Consequently, the
telescoping series converges uniformly on $\{-k,\dots,k\}^{\{j\ge N\}}$ to a
real limit
\[
x(\omega):=\lim_{j\to\infty}\frac{M_j}{Q_j}
=\frac{M_m}{Q_m}+\sum_{j=m}^\infty\frac{W_j}{Q_jQ_{j+1}}\qquad(m\ge N),
\]
with the uniform tail bound
\begin{equation}\label{eq:telescoping-tail}
\left|x(\omega)-\frac{M_m}{Q_m}\right|
\le\sum_{j=m}^\infty\frac{|W_j|}{Q_jQ_{j+1}}
\le\sum_{j=m}^\infty\frac{2k}{Q_{j+1}}
\le\frac{2k}{Q_{m+1}}\sum_{r=0}^\infty 2^{-r}
=\frac{4k}{Q_{m+1}}\qquad\text{for all }m\ge N.
\end{equation}
Multiplying \eqref{eq:telescoping-tail} by $Q_m$ gives, for every $m\ge N$,
\[
|Q_mx(\omega)-M_m|\le\frac{4kQ_m}{Q_{m+1}}\le\frac{4k}{4m+2}<\frac{k}{m}.
\]
Choose $N>8k$ (so $4k/Q_{N+1}<1/8$). Then for every $j\ge N$,
$|Q_jx(\omega)-M_j|<k/j<1/8<1/2$, which ensures that $M_j$ is the \emph{unique}
nearest integer to $Q_jx(\omega)$, that $u_j(x(\omega))=Q_jx(\omega)-M_j$ avoids
all half-integer ties, and that the true recurrence digit of $x(\omega)$ at
index $j$ is $M_{j+1}-A_jM_j-M_{j-1}=d_j$. Moreover,
$|4ju_j(x(\omega))-d_j|\le2|u_j|+|u_{j-1}|+|u_{j+1}|\to0$, so
$4\Lam(x(\omega))=\limsup_{j\to\infty}|d_j|\le k$. Hence the compact image
$K:=x(\{-k,\dots,k\}^{\{j\ge N\}})$ lies in the \emph{non-strict} sublevel
$\{x\in\R:\Lam(x)\le k/4\}$; because $M_N=0$ and $|x(\omega)|\le4k/Q_{N+1}$ by
\eqref{eq:telescoping-tail} at $m=N$, $K$ is contained in
$[-4k/Q_{N+1},4k/Q_{N+1}]$, an interval of length $8k/Q_{N+1}<1/4<1/2$ on which
the canonical projection $\R\to\R/\Z$ is an isometric embedding.

Equip $\{-k,\dots,k\}^{\{j\ge N\}}$ with the uniform Bernoulli product
probability measure assigning independent weight $b^{-1}=(2k+1)^{-1}$ to each
digit $d_j$, and let $\mu$ be its pushforward probability measure on $K$. Take
any interval $U\subset\R$ of length $h=\operatorname{diam}(U)$ with
$1/Q_{n+1}\le h<1/Q_n$ for some $n>N$. If $x,y\in K\cap U$, then for every index
$j$ with $N\le j<n$, we have $|x-y|\le h<1/Q_n$, $|Q_jx-M_j(x)|<1/8$, and
$|Q_jy-M_j(y)|<1/8$, so the triangle inequality gives
\[
|M_j(x)-M_j(y)|
\le|M_j(x)-Q_jx|+Q_j|x-y|+|Q_jy-M_j(y)|
<\frac18+\frac{Q_j}{Q_n}+\frac18=\frac14+\frac{Q_j}{Q_n}\le\frac14+\frac16<1.
\]
Since $M_j(x)$ and $M_j(y)$ are integers, $M_j(x)=M_j(y)$ for all $N\le j<n$
(and $M_{N-1}=0$ by definition). At level $n$, the same estimate gives
\[
|M_n(x)-M_n(y)|<\frac14+Q_n h<\frac14+1=\frac54<2,
\]
so for all $x\in K\cap U$ the integer $M_n(x)$ can take at most $3$ consecutive
integer values (differing by at most $1$ from $\min_{x\in K\cap U}M_n(x)$).
Because the integer sequence $M_{N-1},M_N,\dots,M_n$ uniquely determines the
first $n-N$ digits $d_N,\dots,d_{n-1}$ via $d_j=M_{j+1}-A_jM_j-M_{j-1}$, the
intersection $K\cap U$ is contained in at most $3$ cylinder sets of depth $n-N$,
each of $\mu$-measure $b^{-(n-N)}$. Since $h_k$ is nondecreasing and
$h\ge1/Q_{n+1}$, we have $h_k(h)\ge h_k(1/Q_{n+1})=b^{-(n+1)}$, and therefore
\[
\mu(U)=\mu(K\cap U)\le 3b^{-(n-N)}=3b^{N+1}b^{-(n+1)}\le 3b^{N+1}h_k(h).
\]
Summing over any countable cover $\{U_i\}$ of $K$ with
$\operatorname{diam}(U_i)<1/Q_{N+1}$ gives
$\sum_i h_k(\operatorname{diam}U_i)\ge(3b^{N+1})^{-1}\sum_i\mu(U_i)\ge(3b^{N+1})^{-1}>0$
(the mass distribution principle for $\mathcal H^{h_k}$), so
$\mathcal H^{h_k}(K)\ge(3b^{N+1})^{-1}>0$. By the isometric embedding of
$[-4k/Q_{N+1},4k/Q_{N+1}]$ into $\R/\Z$, the projection of $K$ to the circle has
the same positive $\mathcal H^{h_k}$ measure. Since $K\subseteq\{\Lam\le k/4\}$
and \eqref{eq:strict-null} shows that the strict sublevel $\{\Lam<k/4\}$ has
$\mathcal H^{h_k}$ measure zero, we conclude that
\[
 \mathcal H^{h_k}\bigl(\{x\in\R/\Z:\Lam(x)=k/4\}\bigr)
 \ \ge\ \mathcal H^{h_k}(K)-\mathcal H^{h_k}\bigl(\{\Lam<k/4\}\bigr)
 \ =\ \mathcal H^{h_k}(K)\ >\ 0 .
\]

\emph{Measurability, infinitude and $\sigma$-finiteness.} The functions
$x\mapsto j\|Q_jx\|$ are continuous on the circle $\R/\Z$. Hence $\Lam$ is Borel
measurable as their pointwise limit superior, every sublevel $\{\Lam\le c\}$ and
superlevel $\{\Lam\ge c\}=\bigcap_{m\ge1}\bigcup_{j\ge m}\{x:j\|Q_jx\|\ge c-1/m\}$
is Borel, and so is the exact level
$S_k:=\{x\in\R/\Z:\Lam(x)=k/4\}=\{\Lam\le k/4\}\cap\{\Lam\ge k/4\}$. By
Corollary~\ref{cor:infinite-level} (or Theorem~\ref{thm:finite}), $S_k$ has
Lebesgue measure zero. For any $y\in L=\Z+\Z\frac{\ee+1}{2}$, the triangle
inequality for $\|\cdot\|$ gives
$\bigl|j\|Q_j(x+y)\|-j\|Q_jx\|\bigr|\le j\|Q_jy\|\to0$ as $j\to\infty$ by
Theorem~\ref{thm:class} (the critical weight $j$ is essential here), so
$\Lam(x+y)=\Lam(x)$ and $S_k$ is invariant under translation by the subgroup
$\bar L:=L/\Z\subset\R/\Z$. Since $(\ee+1)/2$ is irrational, Kronecker's theorem
implies that $\bar L$ is dense in $\R/\Z$.

Suppose for a contradiction that $0<\mathcal H^{h_k}(S_k)<\infty$. Then the
normalized restriction
\[
\nu(B):=\frac{\mathcal H^{h_k}(S_k\cap B)}{\mathcal H^{h_k}(S_k)}
\qquad(B\subseteq\R/\Z\text{ Borel})
\]
is a Borel probability measure on $\R/\Z$ concentrated on $S_k$ ($\nu(S_k)=1$).
Because $\mathcal H^{h_k}$ is translation-invariant on $\R/\Z$ and $S_k+\ell=S_k$
for every $\ell\in\bar L$, $\nu$ is invariant under rotation by every
$\ell\in\bar L$: $\nu(B+\ell)=\nu(B)$ for all Borel $B\subseteq\R/\Z$. We show
that $\nu$ is invariant under rotation by \emph{every} $t\in\R/\Z$. Let
$f:\R/\Z\to\R$ be any continuous function. By uniform continuity on the compact
group $\R/\Z$, for any $t\in\R/\Z$ we may choose a sequence $\ell_m\in\bar L$
with $\ell_m\to t$ in $\R/\Z$, and then
$\sup_{x\in\R/\Z}|f(x-\ell_m)-f(x-t)|\to0$ as $m\to\infty$. By $\bar L$-invariance
of $\nu$ and bounded convergence,
\[
\int_{\R/\Z}f(x-t)\,d\nu(x)
=\lim_{m\to\infty}\int_{\R/\Z}f(x-\ell_m)\,d\nu(x)
=\lim_{m\to\infty}\int_{\R/\Z}f(x)\,d\nu(x)
=\int_{\R/\Z}f(x)\,d\nu(x).
\]
Thus the rotated Borel probability measure $\nu_t(B):=\nu(B+t)$ has the same
integral as $\nu$ on every continuous function $f\in C(\R/\Z)$ (or equivalently
the same Fourier coefficients $\int_{\R/\Z}\ee^{2\pi i m x}\,d\nu(x)=0$ for all
$m\in\Z\setminus\{0\}$), so $\nu_t=\nu$ for all $t\in\R/\Z$ by uniqueness of
Haar probability measure on $\R/\Z$: concretely, for every integer $m\ge1$,
partitioning $\R/\Z$ into $m$ disjoint half-open arcs of equal length $1/m$
gives mass $1/m$ to each arc, Hence every arc of rational length $p/m$ has
$\nu$-measure $p/m$, and monotonicity extends $\nu(I)=\operatorname{Leb}_{\R/\Z}(I)$
to all open arcs $I\subset\R/\Z$, which generate the Borel $\sigma$-algebra.
Thus $\nu=\operatorname{Leb}_{\R/\Z}$, contradicting
$1=\nu(S_k)=\operatorname{Leb}_{\R/\Z}(S_k)=0$. Therefore
$\mathcal H^{h_k}(S_k)=\infty$. Finally, $S_k$ has $\sigma$-finite
$\mathcal H^{h_k}$ measure because $S_k\subseteq\{\Lam\le k/4\}$, which is
covered by the countably many Borel classes $\mathcal C(N,m_{N-1},m_N)$ above,
each of finite $\mathcal H^{h_k}$ measure $\le b^{-N}<\infty$.

For the power-gauge family $g_\beta(1/Q_n)=\beta^{-n}$, on each linear
interpolation interval $[1/Q_{n+1},1/Q_n]$ both $g_\beta$ and $h_k$ are affine
with positive values at the endpoints, so the ratio $g_\beta(r)/h_k(r)$ is
monotone on $[1/Q_{n+1},1/Q_n]$ and therefore lies between the endpoint values
$(\beta/(2k+1))^{-n}$ and $(\beta/(2k+1))^{-n-1}$. Thus
$g_\beta(r)/h_k(r)\to0$ when $\beta>2k+1$ and $h_k(r)/g_\beta(r)\to0$ when
$1<\beta<2k+1$ as $r\downarrow0$, which are special cases of
Corollary~\ref{cor:exact-gauge}(ii)--(iii).

For comparison, Cabrelli, Mendivil, Molter and Shonkwiler
\cite[\S5, Theorem~3]{CMMS} construct, for a Cantor set whose gap
sequence is nonincreasing, a gauge $h$ for which the set is an $h$-set
($1/32\le\mathcal H^h\le1$, a finite positive measure), whereas Moreira, Erazo
and Angelini \cite{MEA} show that the classical Lagrange level at value $3$ is
never $\sigma$-finite with positive measure for any gauge. Here the gauge comes
from the recurrence's own scale, and the entire exact level has \emph{infinite}
but $\sigma$-finite $\mathcal H^{h_k}$ measure.
\end{proof}

\paragraph{Lean status.}
The null part \eqref{eq:strict-null} is kernel-checked at Layer~A in
\texttt{EpiFormal/\allowbreak GaugeNull.lean} as
\texttt{GaugeNull.\allowbreak mkMetric\_\allowbreak strictLevel\_\allowbreak null}
for any nondecreasing function $m:\R_{\ge0\infty}\to\R_{\ge0\infty}$ satisfying
$m(1/Q_n)\le C(2k+1)^{-n}$ ($n\ge1$), without requiring continuity of $m$; the
same module kernel-checks Corollary~\ref{cor:exact-gauge}(ii) on $\R$ as
\texttt{GaugeNull.\allowbreak mkMetric\_\allowbreak levelSet\_\allowbreak null\_\allowbreak of\_\allowbreak subgauge}.
In \texttt{EpiFormal/\allowbreak GaugeExact.lean}, $L$-translation invariance
(\texttt{exactLevel\_\allowbreak add\_\allowbreak latticeL}), the countable
finite-measure covering
(\texttt{exactLevel\_\allowbreak countable\_\allowbreak cover\_\allowbreak finite\_\allowbreak mkMetric}),
and---for arbitrary monotone continuous $m$ satisfying matching two-sided
bounds---the Cantor-tree lower bound
(\texttt{mkMetric\_\allowbreak cantorTree\_\allowbreak pos}), exact-level
positivity (\texttt{mkMetric\_\allowbreak exactLevel\_\allowbreak pos}),
exact-level infinitude on $\R$
(\texttt{mkMetric\_\allowbreak exactLevel\_\allowbreak eq\_\allowbreak top}),
and the non-$\sigma$-finiteness criterion
(\texttt{not\_\allowbreak countable\_\allowbreak union\_\allowbreak finite\_\allowbreak of\_\allowbreak ratio\_\allowbreak tendsto\_\allowbreak zero})
are likewise kernel-checked.

The external continuity hypothesis, the explicit piecewise-linear gauge
construction, and the native quotient-metric circle statements on
\texttt{UnitAddCircle} ($\R/\Z$) are completely discharged in
\texttt{EpiFormal/\allowbreak GaugeCircle.lean}:
the gauges $h_k$ (\texttt{critGauge k}), $g_2$ (\texttt{binaryGauge}), and
$g_\beta$ (\texttt{interpGauge \ensuremath{\beta}}) are constructed explicitly
with kernel-checked monotonicity (\texttt{critGauge\_\allowbreak mono}),
continuity on $\R_{\ge0\infty}$ including $0^+$ and $\infty$
(\texttt{critGauge\_\allowbreak continuous}), and exact endpoint and segment
values (\texttt{critGauge\_\allowbreak at\_\allowbreak inv\_\allowbreak Qsel},
\texttt{critGauge\_\allowbreak eq\_\allowbreak segVal}).
This yields the assumption-free instantiations of Theorem~\ref{thm:gauge},
Corollary~\ref{cor:exact-gauge}(i)--(iii), and Remark~\ref{rem:gauge-scale} both
on $\R$
(\texttt{mkMetric\_\allowbreak critGauge\_\allowbreak strictLevel\_\allowbreak null},
\texttt{mkMetric\_\allowbreak critGauge\_\allowbreak exactLevel\_\allowbreak pos},
\texttt{mkMetric\_\allowbreak critGauge\_\allowbreak exactLevel\_\allowbreak eq\_\allowbreak top},
\texttt{critGauge\_\allowbreak exactLevel\_\allowbreak sigma\_\allowbreak finite\_\allowbreak cover},
\texttt{mkMetric\_\allowbreak critGauge\_\allowbreak exactLevel\_\allowbreak inter\_\allowbreak Ioo\_\allowbreak eq\_\allowbreak top},
\texttt{mkMetric\_\allowbreak powerGauge\_\allowbreak exactLevel\_\allowbreak null\_\allowbreak of\_\allowbreak gt},
\texttt{mkMetric\_\allowbreak powerGauge\_\allowbreak exactLevel\_\allowbreak not\_\allowbreak sigma\_\allowbreak finite\_\allowbreak of\_\allowbreak lt},
and \texttt{critGauge\_\allowbreak adjacent\_\allowbreak level\_\allowbreak separation})
and natively on \texttt{UnitAddCircle} ($\R/\Z$)
(\texttt{mkMetric\_\allowbreak critGauge\_\allowbreak circleStrictLevel\_\allowbreak null},
\texttt{mkMetric\_\allowbreak circleLevelSet\_\allowbreak null\_\allowbreak of\_\allowbreak subgauge},
\texttt{mkMetric\_\allowbreak critGauge\_\allowbreak circleExactLevel\_\allowbreak inter\_\allowbreak open\_\allowbreak eq\_\allowbreak top},
\texttt{mkMetric\_\allowbreak critGauge\_\allowbreak circleExactLevel\_\allowbreak eq\_\allowbreak top},
\texttt{critGauge\_\allowbreak circleExactLevel\_\allowbreak sigma\_\allowbreak finite\_\allowbreak cover},
\texttt{mkMetric\_\allowbreak critGauge\_\allowbreak circleExactLevel\_\allowbreak not\_\allowbreak sigma\_\allowbreak finite\_\allowbreak of\_\allowbreak ratio},
\texttt{mkMetric\_\allowbreak powerGauge\_\allowbreak circleExactLevel\_\allowbreak null\_\allowbreak of\_\allowbreak gt},
and
\texttt{mkMetric\_\allowbreak powerGauge\_\allowbreak circleExactLevel\_\allowbreak not\_\allowbreak sigma\_\allowbreak finite\_\allowbreak of\_\allowbreak lt}),
proved constructively by splitting the Cantor tree into depth-$M$ cylinders,
packing $\Omega\bigl((b-a)Q_{J+M+1}\bigr)$ disjoint $L$-translates using
$\delta_j/2\in L$, and transferring across the local isometry
$\R\to\R/\Z$ on intervals of length $\le1/2$. For the binary exact-limit
regime of Theorem~\ref{thm:trichotomy}(ii), the binary Cantor-tree lower bound,
global and open-set $\mathcal H^{g_2}$-infinitude, $\sigma$-finite covering,
$\mathcal H^{h_k}$-nullity, and the $g_\beta$ threshold dichotomy ($\beta>2$ vs.\
$1<\beta<2$) are likewise kernel-checked on both $\R$ and \texttt{UnitAddCircle}
(\texttt{mkMetric\_\allowbreak binaryGauge\_\allowbreak unsigned\allowbreak ExactLimit\_\allowbreak pos},
\texttt{mkMetric\_\allowbreak binaryGauge\_\allowbreak unsigned\allowbreak ExactLimit\_\allowbreak inter\_\allowbreak Ioo\_\allowbreak eq\_\allowbreak top},
\texttt{mkMetric\_\allowbreak binaryGauge\_\allowbreak circle\allowbreak Unsigned\allowbreak ExactLimit\_\allowbreak inter\_\allowbreak open\_\allowbreak eq\_\allowbreak top},
\texttt{mkMetric\_\allowbreak binaryGauge\_\allowbreak circle\allowbreak Unsigned\allowbreak ExactLimit\_\allowbreak eq\_\allowbreak top},
\texttt{binaryGauge\_\allowbreak circle\allowbreak Unsigned\allowbreak ExactLimit\_\allowbreak sigma\_\allowbreak finite\_\allowbreak cover},
\texttt{critGauge\_\allowbreak circle\allowbreak Unsigned\allowbreak ExactLimit\_\allowbreak null},
\texttt{mkMetric\_\allowbreak powerGauge\_\allowbreak circle\allowbreak Unsigned\allowbreak ExactLimit\_\allowbreak null\_\allowbreak of\_\allowbreak gt},
and
\texttt{mkMetric\_\allowbreak powerGauge\_\allowbreak circle\allowbreak Unsigned\allowbreak ExactLimit\_\allowbreak not\_\allowbreak sigma\_\allowbreak finite\_\allowbreak of\_\allowbreak lt}).
All 53 audited declarations in \texttt{GaugeCircle.lean} report only
\texttt{propext}, \texttt{Classical.choice}, and \texttt{Quot.sound}, with zero
\texttt{sorryAx}.

The gauge $h_k$ is exact for the level $k/4$ in the following sense.
Here a gauge is a nondecreasing function $g:[0,\infty)\to[0,\infty)$
with $g(0)=0$, and $\mathcal H^g$ is the Hausdorff measure on the circle
with covering cost $g(\operatorname{diam}U)$.

\begin{corollary}[Exact critical gauge]\label{cor:exact-gauge}
Let $k\ge1$ and $S_k=\{x\in\R/\Z:\Lam(x)=k/4\}$.
\begin{enumerate}[label=(\roman*)]
\item For every nonempty open arc $I\subseteq\R/\Z$,
$\mathcal H^{h_k}(S_k\cap I)=\infty$.
\item If $g$ is a gauge with $g(1/Q_n)/h_k(1/Q_n)\to0$ as $n\to\infty$, then
$\mathcal H^{g}(\{\Lam\le k/4\})=0$; in particular $\mathcal H^g(S_k)=0$.
(Lean on $\R$ and $\R/\Z$: \texttt{GaugeNull.\allowbreak mkMetric\_\allowbreak levelSet\_\allowbreak null\_\allowbreak of\_\allowbreak subgauge} and \texttt{GaugeCircle.\allowbreak mkMetric\_\allowbreak circleLevelSet\_\allowbreak null\_\allowbreak of\_\allowbreak subgauge}.)
\item If $g$ is a gauge with $h_k(t)/g(t)\to0$ as $t\to0^+$, then $S_k$ is
not $\sigma$-finite for $\mathcal H^g$; in particular
$\mathcal H^g(S_k)=\infty$.
\end{enumerate}
\end{corollary}

\begin{proof}
(i) Hausdorff measures on the circle are invariant under rotations, and
$S_k$ is invariant under rotation by the elements of $L$ (proof of
Theorem~\ref{thm:gauge}). Suppose $\mathcal H^{h_k}(S_k\cap I)<\infty$ for
some nonempty open arc $I$. Since $L$ is dense modulo one, the arcs $I+\ell$
($\ell\in L$) cover the circle, and by compactness finitely many of them,
say $I+\ell_1,\ldots,I+\ell_N$, already do. Then
$\mathcal H^{h_k}(S_k)\le\sum_{i=1}^N\mathcal H^{h_k}(S_k\cap(I+\ell_i))
=N\,\mathcal H^{h_k}(S_k\cap I)<\infty$, contradicting
Theorem~\ref{thm:gauge}.

(ii) In the proof of Theorem~\ref{thm:gauge}, the non-strict sublevel
$\{\Lam\le k/4\}$ is a countable union of classes. For a class with
starting index $N$ and every large $n$, at most $b^{\,n-N}$ intervals of
diameter at most $1/Q_n$ cover it, where $b=2k+1$. Since $g$ is
nondecreasing, the $g$-cost of this cover is at most
\[
 b^{\,n-N}g(1/Q_n)=b^{-N}\,\frac{g(1/Q_n)}{h_k(1/Q_n)}\longrightarrow0 .
\]
So every class is $\mathcal H^g$-null, and so is their countable union.

(iii) Fix $\varepsilon>0$ and choose $\delta>0$ with
$h_k(t)\le\varepsilon g(t)$ for $0\le t<\delta$. Every $\delta$-cover then
satisfies $\sum h_k(\operatorname{diam}U)\le
\varepsilon\sum g(\operatorname{diam}U)$, so
$\mathcal H^{h_k}(E)\le\varepsilon\,\mathcal H^g(E)$ for every set $E$.
Hence $\mathcal H^g(E)<\infty$ implies $\mathcal H^{h_k}(E)=0$. If $S_k$
were a countable union of sets $E_i$ with $\mathcal H^g(E_i)<\infty$, we
would get $\mathcal H^{h_k}(S_k)=0$, contradicting Theorem~\ref{thm:gauge}.
\end{proof}

\begin{corollary}[Arbitrary gauges at the exact and unsigned levels]\label{cor:liminf-gauges}
Let $g$ be a Hausdorff gauge, $r_n=1/Q_n$, and
\[
\rho_b(g)=\liminf_{n\to\infty} b^n g(r_n).
\]
For $S_k=\{\Lam=k/4\}$ take $b=2k+1$, and for
$T_k=\{x:j\|Q_jx\|\to k/4\}$ take $b=2$, where $k\ge1$ and both sets
are taken on the circle. In each case the named set is
$\mathcal H^g$-null if $\rho_b(g)=0$. If $0<\rho_b(g)<\infty$, it is
$\sigma$-finite and its intersection with every nonempty open set has
infinite $\mathcal H^g$ measure. If $\rho_b(g)=\infty$, its intersection
with every nonempty open set is non-$\sigma$-finite. For $b=2k+1$ the
zero assertion applies to the full sublevel $\{\Lam\le k/4\}$, and
$\rho_b(g)<\infty$ implies that $\{\Lam<k/4\}$ is $\mathcal H^g$-null.
\end{corollary}

\begin{proof}
The fixed initial-pair classes of Theorems~\ref{thm:gauge}
and~\ref{thm:trichotomy}(ii) have covers by at most $b^{n-N}$ intervals
of diameter $r_n$ for a fixed starting level $N$ and all large $n$.
Their gauge cost is bounded by $b^{-N}b^n g(r_n)$. Taking the lower
limit along permissible scales gives zero or finite measure for each
class, according as $\rho_b(g)$ is zero or finite, and the classes are
countable. For the strict sublevel the same argument uses $2k-1$
instead of $2k+1$. Bounded cost along a subsequence at base $2k+1$
then tends to zero along that subsequence at base $2k-1$.

Write $h_b(r_n)=b^{-n}$ with linear interpolation. If
$r_{n+1}\le r\le r_n$, then
\[
 h_b(r)\le b^{-n},\qquad g(r)\ge g(r_{n+1}).
\]
For any $0<A<\rho_b(g)$, eventually $b^{n+1}g(r_{n+1})\ge A$,
so $h_b(r)\le(b/A)g(r)$ for all sufficiently small $r$. The critical
reference measure is infinite on the intersection with every nonempty
open set, by Corollary~\ref{cor:exact-gauge}(i) and
Theorem~\ref{thm:trichotomy}(ii), respectively. The comparison gives the
same local infinitude for $g$. When $\rho_b(g)=\infty$, $A$ can be
arbitrarily large, so any set of finite $\mathcal H^g$ measure is
$\mathcal H^{h_b}$-null. A countable union of such sets cannot cover
a nonempty open intersection of the named set.
\end{proof}

\begin{remark}[Oscillating normalized costs]
Let $n_j=2^j$ and prescribe $g(r_n)=b^{-n_j}$ for
$n_j\le n<n_{j+1}$, interpolating linearly in $r$. The lower and upper
limits of $b^ng(r_n)$ are $1$ and $\infty$. Prescribing
$g(r_n)=b^{-(n_{j+1}-1)}$ instead gives lower and upper limits $0$ and $1$.
Both prescriptions are nondecreasing as functions of $r$ and tend to zero
at the origin. They can be made strictly increasing near zero by adding
$h_b^2$, which adds $b^{-n}$ to the normalized costs at the nodes.
\end{remark}

\paragraph{Lean verification.}
The accompanying \texttt{GaugeLiminfComplete.lean} verifies
Corollary~\ref{cor:liminf-gauges} on \texttt{UnitAddCircle}, including
the strict-sublevel assertion and the non-$\sigma$-finite conclusion on
each nonempty open set. Its three principal declarations are
\texttt{circleExactLevel\_\allowbreak liminf\_\allowbreak trichotomy},
\texttt{circleUnsigned\allowbreak ExactLimit\_\allowbreak liminf\_\allowbreak trichotomy}, and
\texttt{mkMetric\_\allowbreak circleStrictLevel\_\allowbreak null\_\allowbreak of\_\allowbreak liminf\_\allowbreak lt\_\allowbreak top}.
All eighteen audited declarations use only the three base axioms.

\begin{remark}[Level separation and interpolation-free logarithmic scale]\label{rem:gauge-scale}
The gauges $h_k$ strictly separate the critical levels. Indeed, for $k\ge2$,
$h_k(1/Q_n)/h_{k-1}(1/Q_n)=\bigl((2k-1)/(2k+1)\bigr)^n\to0$, and on each
interval $[1/Q_{n+1},1/Q_n]$ both gauges are linear, so their ratio lies
between its values at the endpoints. Hence $h_k(t)/h_{k-1}(t)\to0$ as
$t\to0^+$. By Corollary~\ref{cor:exact-gauge}(iii), the level $S_k$ is not
$\sigma$-finite for $\mathcal H^{h_{k-1}}$, whereas $S_{k-1}$ has infinite
but $\sigma$-finite $\mathcal H^{h_{k-1}}$ measure. By
Corollary~\ref{cor:exact-gauge}(ii), the level $S_{k-1}$ is
$\mathcal H^{h_k}$-null, consistent with \eqref{eq:strict-null}.

Moreover, the scale of the critical gauge $h_k$ can be read without reference
to the linear interpolation. For $1/Q_{n+1}\le r\le1/Q_n$ the quantity
$-\log h_k(r)$ lies between $n\log(2k+1)$ and $(n+1)\log(2k+1)$.
By \eqref{eq:rec} both endpoints for $\log(1/r)$ are asymptotic to
$n\log n$, giving
\[
 -\log h_k(r)\sim \log(2k+1)\,
             \frac{\log(1/r)}{\log\log(1/r)}\qquad(r\downarrow0).
\]
In particular, any gauge $g$ whose corresponding logarithmic rate exists and is
strictly greater than $\log(2k+1)$ satisfies $g(1/Q_n)/h_k(1/Q_n)\to0$ and
therefore gives zero measure to $S_k$, while any gauge with a strictly smaller
logarithmic rate satisfies $h_k(t)/g(t)\to0$ and gives a non-$\sigma$-finite
measure on $S_k$.
\end{remark}

\subsection{The three-way splitting at each positive spectral level}

While every subcritical weight $0\le\alpha<1$ assigns the same Hausdorff
dimension $1-\alpha$ to the sublevel set $\{\Lam_\alpha\le c\}$ and to the
exact limit set $\{j^\alpha\|Q_jx\|\to c\}$ (Theorem~\ref{conj:weight}), the
critical weight $\alpha=1$ exhibits a dramatic three-way splitting at each
positive quantized level $c=k/4$ ($k\in\Z_{\ge1}$), governed by the allowed
tail alphabet of the integer recurrence digits $d_j(x)$.

\begin{theorem}[Trichotomy at positive spectral levels, Layer~A in Lean]\label{thm:trichotomy}
Fix an integer $k\ge1$ and let $u_j(x)=Q_jx-\mathrm{round}(Q_jx)$. At the
spectral level $c=k/4$, the sublevel set, the unsigned exact limit set, and
the signed exact limit sets split into three distinct regimes:
\begin{enumerate}[label=(\roman*),leftmargin=2em]
\item \textbf{Limsup level (alphabet of size $2k+1$, gauge base $2k+1$):}
$x\in E_{k/4}=\{\Lam\le k/4\}$ if and only if $u_j(x)\to0$ and eventually
$d_j(x)\in\{-k,-k+1,\dots,k\}$. On $\R/\Z$, both $E_{k/4}$ and the exact
limsup level $\{\Lam=k/4\}$ have infinite, $\sigma$-finite $\mathcal H^{h_k}$
measure for the exact critical gauge $h_k(1/Q_n)=(2k+1)^{-n}$
(Theorem~\ref{thm:gauge} and Corollary~\ref{cor:exact-gauge}).
\item \textbf{Unsigned exact limit set (universal binary alphabet of size $2$, gauge base $2$):}
Define $E_{\mathrm{lim}}(k/4):=\{x\in\R:\lim_{j\to\infty}j\|Q_jx\|=k/4\}$.
Then $x\in E_{\mathrm{lim}}(k/4)$ if and only if $u_j(x)\to0$ and eventually
$|d_j(x)|=k$, i.e.\ $d_j(x)\in\{-k,+k\}$ (exactly two branches per step,
\emph{independently} of $k\ge1$). Consequently, in the interpolated gauge
family $g_\beta(1/Q_n)=\beta^{-n}$, the exact critical gauge for
$E_{\mathrm{lim}}(k/4)$ is $h_{\mathrm{lim}}:=g_2$ ($h_{\mathrm{lim}}(1/Q_n)=2^{-n}$)
for \emph{every} $k\ge1$: on $\R/\Z$, $E_{\mathrm{lim}}(k/4)$ has infinite,
$\sigma$-finite $\mathcal H^{g_2}$ measure (and infinite $\mathcal H^{g_2}$
measure in every nonempty open arc, whereas
$\mathcal H^{h_k}(E_{\mathrm{lim}}(k/4))=0$ since $2<2k+1$), and
$E_{\mathrm{lim}}(c)=\emptyset$ for all $c\notin\frac14\Z_{\ge0}$.
\item \textbf{Signed exact limit sets (single tail branch, single countable $L$-coset):}
For each sign $\pm$, define
$E_{\mathrm{lim}}^\pm(k/4):=\{x\in\R:\lim_{j\to\infty}j\,u_j(x)=\pm k/4\}$.
Then $x\in E_{\mathrm{lim}}^\pm(k/4)$ if and only if $u_j(x)\to0$ and
eventually $d_j(x)=\pm k$. For any $x,y\in E_{\mathrm{lim}}^\pm(k/4)$, their
difference satisfies $x-y\in L=\Z+\Z\frac{\ee+1}{2}$. Hence each signed limit
set is a single countable coset of $L$:
\[
E_{\mathrm{lim}}^+(k/4)=x_k^+ + L,\qquad
E_{\mathrm{lim}}^-(k/4)=x_k^- + L\quad(x_k^-=-x_k^+),
\]
and therefore is countable and has zero measure for every non-atomic Hausdorff
gauge.
\end{enumerate}
(Lean declarations in \texttt{Full\allowbreak Spectrum\allowbreak Discrete.lean},
\texttt{Gauge\allowbreak Circle.lean}, and
\texttt{Trichotomy\allowbreak And\allowbreak Poly\allowbreak Extension.lean}:
the binary digit characterization \texttt{unsigned\_\allowbreak exact\_\allowbreak limit\_\allowbreak iff\_\allowbreak defect\_\allowbreak binary},
the constant digit characterization \texttt{signed\_\allowbreak exact\_\allowbreak limit\_\allowbreak iff\_\allowbreak defect\_\allowbreak const},
the $L$-coset difference theorem \texttt{sub\_\allowbreak mem\_\allowbreak latticeL\_\allowbreak of\_\allowbreak same\_\allowbreak signed\_\allowbreak limit},
countability \texttt{signed\_\allowbreak exact\_\allowbreak limit\_\allowbreak set\_\allowbreak countable},
zero dimension \texttt{dimH\_\allowbreak signed\_\allowbreak exact\_\allowbreak limit\_\allowbreak set\_\allowbreak zero},
the native $\R$ and \texttt{UnitAddCircle} binary-gauge theorems for~(ii) in \texttt{GaugeCircle.lean},
and the full signed-coset collapse for~(iii) via
\texttt{xPlus\_\allowbreak mem\_\allowbreak signed\allowbreak Exact\allowbreak LimitSet},
\texttt{signed\allowbreak Exact\allowbreak LimitSet\_\allowbreak pos\_\allowbreak eq\_\allowbreak coset},
\texttt{signed\allowbreak Exact\allowbreak LimitSet\_\allowbreak neg\_\allowbreak eq\_\allowbreak coset}, and
\texttt{trichotomy\_\allowbreak part\_\allowbreak iii\_\allowbreak complete}.)
\end{theorem}

\begin{proof}
Statement~(i) is immediate from Theorem~\ref{thm:full-spectrum},
Corollary~\ref{thm:sub}, Theorem~\ref{thm:gauge}, and
Corollary~\ref{cor:exact-gauge}.

For~(ii), if $j\|Q_jx\|\to k/4$, then $u_j(x)=O(1/j)\to0$, and
\eqref{eq:digit-asymp} gives $|d_j(x)|\to k$. Since $|d_j(x)|\in\Z_{\ge0}$, we
must have $|d_j(x)|=k$, i.e.\ $d_j(x)\in\{-k,+k\}$, for all sufficiently large
$j$. Conversely, if $u_j(x)\to0$ and eventually $d_j(x)\in\{-k,+k\}$, then
$|d_j(x)|=k$ eventually, so \eqref{eq:digit-asymp} implies $4j\|Q_jx\|\to k$,
that is, $x\in E_{\mathrm{lim}}(k/4)$ (and if $c\notin\frac14\Z_{\ge0}$, no
integer sequence $|d_j|$ can converge to $4c$, so $E_{\mathrm{lim}}(c)=\emptyset$).
We now verify each metric assertion for $E_{\mathrm{lim}}(k/4)$ on $\R/\Z$ under
the binary gauge $g_2(1/Q_n)=2^{-n}$:
\begin{itemize}[leftmargin=1.8em]
\item \emph{$\sigma$-finiteness under $\mathcal H^{g_2}$ and nullity under $\mathcal H^{h_k}$:}
Partition $E_{\mathrm{lim}}(k/4)$ into countably many Borel classes
$\mathcal C_{\mathrm{bin}}(N,m_{N-1},m_N)$ indexed by a starting index $N\ge2$
and starting pair $(M_{N-1},M_N)=(m_{N-1},m_N)\in\Z^2$, requiring
$|u_j|\le(k+1)/j$ and $d_j\in\{-k,+k\}$ for all $j\ge N$. Since at each step
$j\ge N$ there are only $2$ choices for $d_j\in\{-k,+k\}$, at level $n>N$ a
class has at most $2^{n-N}$ possible values of $M_n$, each confining $x$ to an
interval of diameter $2(k+1)/(nQ_n)\le1/Q_n$ for all $n\ge2(k+1)$. Thus the
$\mathcal H_\delta^{g_2}$ covering cost of $\mathcal C_{\mathrm{bin}}(N,m_{N-1},m_N)$
is at most $2^{n-N}g_2(1/Q_n)=2^{-N}<\infty$, proving that
$E_{\mathrm{lim}}(k/4)$ is $\sigma$-finite for $\mathcal H^{g_2}$. For the
larger-base gauge $h_k=g_{2k+1}$, since $2<2k+1$ for every $k\ge1$, the
$\mathcal H_\delta^{h_k}$ covering cost of each class is at most
\[
2^{n-N}h_k(1/Q_n)=(2k+1)^{-N}\left(\frac{2}{2k+1}\right)^{n-N}\longrightarrow0
\qquad(n\to\infty),
\]
so every class is $\mathcal H^{h_k}$-null and
$\mathcal H^{h_k}(E_{\mathrm{lim}}(k/4))=0$. More generally, any gauge $g$ with
$g(1/Q_n)/g_2(1/Q_n)=2^ng(1/Q_n)\to0$ gives
$\mathcal H^g(E_{\mathrm{lim}}(k/4))=0$.
\item \emph{Positivity under $\mathcal H^{g_2}$:}
Fix $N>8k$, set $M_{N-1}=M_N=0$, and restrict the digit sequences in the
\emph{Positivity} step of Theorem~\ref{thm:gauge} to the binary space
$\{-k,+k\}^{\{j\ge N\}}$. By \eqref{eq:telescoping-tail}, every resulting point
$x(\omega)=\lim_{j\to\infty}M_j/Q_j$ satisfies $|Q_jx(\omega)-M_j|<k/j<1/8$ and
has true recurrence digits $d_j(x(\omega))=d_j\in\{-k,+k\}$ for all $j\ge N$, so
$|d_j(x(\omega))|=k$ and $j\|Q_jx(\omega)\|\to k/4$. Hence the compact binary
Cantor image $K_{\mathrm{bin}}:=x(\{-k,+k\}^{\{j\ge N\}})$ lies entirely inside
$E_{\mathrm{lim}}(k/4)\cap[-4k/Q_{N+1},4k/Q_{N+1}]$ (no strict sublevel needs to
be subtracted). Equip $\{-k,+k\}^{\{j\ge N\}}$ with the uniform Bernoulli
product measure assigning weight $1/2$ to $-k$ and $1/2$ to $+k$, and let
$\mu_{\mathrm{bin}}$ be its pushforward to $K_{\mathrm{bin}}$. As in the proof
of Theorem~\ref{thm:gauge}, any interval $U$ of length
$1/Q_{n+1}\le h<1/Q_n$ ($n>N$) determines a unique $M_j$ for each $N\le j<n$ and
at most $3$ values of $M_n$, so $K_{\mathrm{bin}}\cap U$ meets at most $3$ binary
cylinders of depth $n-N$, each of $\mu_{\mathrm{bin}}$-mass $2^{-(n-N)}$. Thus
\[
\mu_{\mathrm{bin}}(U)\le 3\cdot 2^{-(n-N)}=3\cdot 2^{N+1}\cdot 2^{-(n+1)}
\le 3\cdot 2^{N+1}g_2(h),
\]
and the mass distribution principle gives
$\mathcal H^{g_2}(E_{\mathrm{lim}}(k/4))\ge\mathcal H^{g_2}(K_{\mathrm{bin}})\ge(3\cdot 2^{N+1})^{-1}>0$.
\item \emph{Infinitude on every nonempty open arc and non-$\sigma$-finiteness for larger gauges:}
For any $y\in L$, $j\|Q_jy\|\to0$ by Theorem~\ref{thm:class}, so
$\bigl|j\|Q_j(x+y)\|-j\|Q_jx\|\bigr|\le j\|Q_jy\|\to0$; hence
$E_{\mathrm{lim}}(k/4)$ is Borel, Lebesgue-null (since
$E_{\mathrm{lim}}(k/4)\subseteq E_{k/4}$), and invariant under translation by
the dense subgroup $\bar L=L/\Z$ on $\R/\Z$. The Haar-uniqueness argument of
Theorem~\ref{thm:gauge} therefore forces
$\mathcal H^{g_2}(E_{\mathrm{lim}}(k/4))=\infty$ on $\R/\Z$, and the compactness
covering argument of Corollary~\ref{cor:exact-gauge}(i),(iii) with $g_2$ in
place of $h_k$ shows that $\mathcal H^{g_2}(E_{\mathrm{lim}}(k/4)\cap I)=\infty$
for every nonempty open arc $I\subseteq\R/\Z$ and that $E_{\mathrm{lim}}(k/4)$
is not $\sigma$-finite for $\mathcal H^g$ whenever $g_2(t)/g(t)\to0$ as
$t\to0^+$ (in particular for $g_\beta$ with $1<\beta<2$).
\end{itemize}

For~(iii), the signed asymptotic $d_j(x)-4ju_j(x)\to0$ in
\eqref{eq:digit-asymp} shows that $ju_j(x)\to\pm k/4$ if and only if
$u_j(x)\to0$ and the integer $d_j(x)$ is eventually constant and equal to
$\pm k$. The existence of a witness $x_k^+\in E_{\mathrm{lim}}^+(k/4)$ (and
$x_k^-=-x_k^+\in E_{\mathrm{lim}}^-(k/4)$) is given by the constant-digit
construction $d_j\equiv k$ in Theorem~\ref{thm:gauge}. If
$x,y\in E_{\mathrm{lim}}^\pm(k/4)$, then $j(u_j(x)-u_j(y))\to0$; since
$Q_j(x-y)=\bigl(M_j(x)-M_j(y)\bigr)+\bigl(u_j(x)-u_j(y)\bigr)$ with
$M_j(x)-M_j(y)\in\Z$, we have $\|Q_j(x-y)\|\le|u_j(x)-u_j(y)|$, hence
$j\|Q_j(x-y)\|\to0$. By Theorem~\ref{thm:class}, $x-y\in L$. Conversely, for
any $z\in L$, $j\|Q_jz\|\to0$ gives $u_j(z)=o(1/j)\to0$; for
$x\in E_{\mathrm{lim}}^\pm(k/4)$, since $u_j(x)=O(1/j)\to0$, we have
$|u_j(x)+u_j(z)|<1/2$ for all large $j$, so $Q_j(x+z)$ rounds to
$M_j(x)+M_j(z)$ and $u_j(x+z)=u_j(x)+u_j(z)$, whence
$j\,u_j(x+z)=j\,u_j(x)+j\,u_j(z)\to\pm k/4$. Thus $x+z\in E_{\mathrm{lim}}^\pm(k/4)$
and $E_{\mathrm{lim}}^\pm(k/4)=x_k^\pm+L$.
\end{proof}

\paragraph{Lean status for Theorem~\ref{thm:trichotomy}.}
All three parts~(i), (ii), and~(iii) of Theorem~\ref{thm:trichotomy} are now
kernel-checked at Layer~A (using only the three Lean base axioms
\texttt{propext}, \texttt{Classical.choice}, and \texttt{Quot.sound}, with zero
\texttt{sorryAx}). Specifically, the digit characterizations, signed-limit
$L$-difference, countability, and dimension-zero lemmas are proved in
\texttt{EpiFormal/\allowbreak Full\allowbreak Spectrum\allowbreak Discrete.lean};
all metric gauge statements for the limsup regime~(i) and the binary exact-limit
regime $E_{\mathrm{lim}}(k/4)$ in~(ii)---binary Cantor-tree positivity, global
and nonempty-open-subset $\mathcal H^{g_2}$-infinitude, $\sigma$-finite
covering, $\mathcal H^{h_k}$-nullity, and the $g_\beta$ threshold dichotomy
($\beta>2$ vs.\ $1<\beta<2$)---are proved natively on both $\R$ and
\texttt{UnitAddCircle} ($\R/\Z$) in
\texttt{EpiFormal/\allowbreak GaugeCircle.lean} (building on
\texttt{GaugeNull.lean} and \texttt{GaugeExact.lean}); and the full
signed-coset collapse $E_{\mathrm{lim}}^\pm(k/4)=x_k^\pm+L$ in part~(iii) is
proved in
\texttt{EpiFormal/\allowbreak TrichotomyAndPolyExtension.lean} across eight
audited declarations (\texttt{xPlus\_\allowbreak mem\_\allowbreak signedExactLimitSet},
\texttt{neg\_\allowbreak xPlus\_\allowbreak mem\_\allowbreak signedExactLimitSet},
\texttt{add\_\allowbreak latticeL\_\allowbreak mem\_\allowbreak signedExactLimitSet},
\texttt{signedExactLimitSet\_\allowbreak pos\_\allowbreak eq\_\allowbreak coset},
\texttt{signedExactLimitSet\_\allowbreak neg\_\allowbreak eq\_\allowbreak coset},
\texttt{signedExactLimitSet\_\allowbreak zero\_\allowbreak eq\_\allowbreak latticeL},
\texttt{signedExactLimitSet\_\allowbreak empty\_\allowbreak of\_\allowbreak not\_\allowbreak quarter\_\allowbreak int},
and \texttt{trichotomy\_\allowbreak part\_\allowbreak iii\_\allowbreak complete}).

\section{Finite-state tails, frequency spectra, and the cubic-root critical amplitude}\label{sec:finite-state}

This section develops the finite-state layer of the theory. It
uses the gauge functions and covering constructions of
Section~\ref{sec:gauge} (notably the fixed-start covering
argument in the proof of Theorem~\ref{thm:gauge}) together with
its own three-prefix covering estimate
(Theorem~\ref{thm:finite-state-quant}); none of its arguments
depend on the polynomial-growth recurrences of
Section~\ref{sec:poly-rec} or on the continued-fraction window
theory of Section~\ref{sec:windows-residues}.

\paragraph{Notation for this section.}
Let $G$ be a finite directed graph whose edges carry labels in
$\{-k,\ldots,k\}$, distinct on the outgoing edges at each
vertex, and let $\chi(e)\in\{0,1\}$ mark the extreme labels
$\pm k$. The tilted weighted adjacency matrix is
$A_t(v,u)=\sum_{e:v\to u}\exp(t\chi(e))$; its Perron root
$\rho(A_t)$ defines the \emph{pressure} $P(t)=\log\rho(A_t)$,
with equilibrium frequency $p_t=P'(t)$ and conditional entropy
$h_t=P(t)-tp_t$. The associated Markov kernel
$K_t(e)=\exp(t\chi(e)-P(t))\,w_t(u)/w_t(v)$, where $w_t>0$ is a
right Perron vector, has stationary law $\pi_t$ and is used
from the frequency-spectrum subsection onward. Along a
length-$m$ path, $J_m$ denotes the number of extreme labels,
and, for $t\ne0$, the \emph{signed displacement} is
$V_m=-\operatorname{sgn}(t)(J_m-mp_t)$. Implied constants may
depend on $G$, $k$, $t$, and on the endpoints of any interval
in the statements.

\subsection{Finite-state constraints on defect tails}\label{sec:finite-state-tail}

Let $G$ be a finite directed graph with edge labels in
$\{-k,\ldots,k\}$, with distinct labels on the outgoing edges at each
vertex. Mark the edges labelled $k$ or $-k$. Let $S_k(G)\subset S_k$ be
the points whose eventual defect sequence labels an infinite path in $G$.
A strongly connected component is called marked if it contains an internal
marked edge. Write $\lambda_C$ for the Perron root of the adjacency matrix
of a marked component $C$, counting edge multiplicities.

\begin{theorem}[Finite-state defect tails]\label{thm:finite-state-tail}
If $G$ has no marked component, then $S_k(G)=\varnothing$.
Otherwise put $\lambda_*=\max_C\lambda_C$, where $C$ ranges over
the marked components. If $\lambda_*=1$, then $S_k(G)$ is countable.
If $\lambda_*>1$, it satisfies the three measure conclusions of
Corollary~\ref{cor:liminf-gauges}, with $S_k(G)$ in place of the named
set and with $b=\lambda_*$. In the zero case the null conclusion refers
to $S_k(G)$.
\end{theorem}

\begin{proof}
First consider a strongly connected marked graph with Perron root
$\lambda>1$. Take a positive right Perron vector $w$. The sum of
terminal weights over paths of length $m$ from vertex $v$ equals
$\lambda^m w_v$. Since the coordinates of $w$ are bounded above and
away from zero, the path count is comparable to $\lambda^m$, uniformly
over the initial vertex. The fixed initial-pair covering argument
therefore gives at most $C\lambda^{n-N}$ intervals of diameter $r_n$
for each fixed starting position, pair and vertex. This gives the
zero and finite-cover conclusions when the normalized lower limit is
zero or finite, respectively.

To obtain positive critical measure, give an edge $v\to v'$ probability
$w_{v'}/(\lambda w_v)$. The path-cylinder mass telescopes to
$w_{v_m}/(\lambda^m w_{v_0})\le C\lambda^{-m}$. Fix one marked edge.
From every vertex there is a path crossing it in at most some fixed
number $L_0$ of steps, with probability at least $p>0$. Conditioning
successively in blocks of length $L_0$ bounds the probability of
avoiding the edge for $m$ blocks by $(1-p)^m$. Applying this after each
starting time shows that almost every path crosses a marked edge
infinitely often.

Feed the edge labels into the defect recurrence, with the two
initial integers zero and with starting index $N>8k$ sufficiently
large. The construction in the proof of Theorem~\ref{thm:gauge} gives
a point $x=\lim M_j/Q_j$ with
$|Q_jx-M_j|\le 4kQ_j/Q_{j+1}<1/8$ for $j\ge N$.
Thus the $M_j$ are its nearest integers and the prescribed labels
are its defects. Infinitely many marked labels give $\Lam(x)=k/4$
by the defect asymptotic. The pushforward of the path probability
therefore gives full mass to $S_k(G)$.

If a set of diameter $r$ satisfies $r_{n+1}\le r<r_n$, the difference
between two of its coded values of $M_j$ is less than
$1/4+Q_jr$. For $j<n$ with $j\ge N$ this is less than one, and
at $j=n$ there are at most three possible adjacent integers.
Consequently the set meets at most three label cylinders of length
$n-N$. From a prescribed start vertex each label word specifies at
most one path. Thus its mass is at most $3C\lambda^{-(n-N)}$, which is
bounded by a fixed multiple of $h_\lambda(r)$ for the interpolated gauge
$h_\lambda(r_n)=\lambda^{-n}$. Singleton mass is zero by the same
cylinder estimate. The construction is in a real interval of length
less than $1/2$ if $N$ is large, so projection to the circle is
isometric there. The mass-distribution bound gives
$\mathcal H^{h_\lambda}(S_k(G))>0$.

For Borel measurability, intersect with the Borel level $S_k$. For each
start and vertex, the condition that every finite prefix is admissible
is a countable intersection of Borel conditions on the rounding
functions. An infinite path exists by finite branching. Take the
countable union over starts and the finite union over vertices.
For $y\in L$, the defect asymptotic gives $d_j(y)=0$ eventually.
Because the nearest-integer errors of $x\in S_k$ and of $y$ tend to
zero, their nearest integers add eventually, so $d_j(x+y)=d_j(x)$
eventually. Hence $S_k(G)$ is invariant under the dense group $L$
and is Lebesgue-null. If its positive critical measure were finite,
its normalized restriction would be invariant under this dense group
and, by continuity against continuous test functions, under the
entire circle. It would then be Haar measure, a contradiction.
For a nonempty open set $U$, the family $U+\ell$, $\ell\in L\bmod 1$,
covers the circle by density and has a finite subcover by compactness.
Invariance under those translations would make the whole level have
finite measure if $S_k(G)\cap U$ had finite measure. This proves local
infinitude. The comparison of adjacent scales in
Corollary~\ref{cor:liminf-gauges} now gives the three conclusions for
every gauge. The middle case combines the countable covering by
fixed initial-pair classes with local infinitude on the full level.

For a general finite graph, each infinite path eventually remains
in one strongly connected component, since the component graph is
acyclic. If its point is in $S_k$, its defects have absolute value
$k$ infinitely often. Thus the terminal component is marked, and
\[
 S_k(G)=\bigcup_{C\ {\rm marked}} S_k(C).
\]
For a marked component with $\lambda_C=1$, the Perron count bounds the
number of intervals in each fixed-start cover by a uniform integer $B$.
A class with such covers is finite. If it contained $B+1$ distinct points,
their minimum pairwise distance would be positive. At a sufficiently
fine level each of the $B$ intervals could contain at most one of those
points. This contradiction gives finiteness for each fixed start index,
vertex and integer pair, so the union is countable.

Finally, when $1<\lambda_C<\lambda_*$ and
$\rho_{\lambda_*}(g)<\infty$, a bounded subsequence at base $\lambda_*$
has normalized cost tending to zero at base $\lambda_C$ by the
factor $(\lambda_C/\lambda_*)^n$. Such components are null.
Countable components are also null. At least one component attains
$\lambda_*$, giving local infinitude or local non-$\sigma$-finiteness
in the corresponding regimes. The finite union of the component
covers gives the asserted global finite-cover conclusion.
\end{proof}

\paragraph{Lean status for Theorem~\ref{thm:finite-state-tail}.}
\texttt{FiniteStateTail.lean} verifies the marked-component decomposition
and the emptiness conclusion, the Perron-root-one countability conclusion,
and the fixed-start Perron upper bound on the line and circle. It also
verifies the null and countable finite-cover consequences within each
component, given the positive Perron vector, and global nullity when the dominant normalized cost has lower limit zero. The fixed-start upper
constant is $w_v/(\min_u w_u\,\lambda^N)$.
The positive-measure lower bound for recurrent marked coding and the
resulting local infinitude are established by the measure argument
in the proof above.

\begin{remark}[Two graph examples]
Forbidding successive marked labels gives the adjacency matrix
\[
 \begin{pmatrix}2k-1&2\\2k-1&0\end{pmatrix}.
\]
For $k\ge2$ its characteristic polynomial has opposite signs at $2$
and $2k+1$, so its Perron root is strictly between those values.
For a reducible example, take a disjoint unmarked vertex with the
$2k-1$ unmarked loops and a marked vertex with the two loops labelled
$\pm k$. For $k\ge2$ the larger adjacency root comes from the unmarked
vertex, while $S_k(G)$ has critical base $2$. Thus the root of the
whole adjacency matrix need not determine the gauge of an exact level.
\end{remark}

For comparison, the graph-directed similarity results of Mauldin
and Williams \cite[Theorems~3 and~4]{MW} use a spectral radius of weighted
similarity matrices and study the full construction. Here the scales
are supplied by $Q_n$, and the insistence on infinitely many marked
edges arises from exact-level membership. The local infinitude uses
the arithmetic translation group of the level.

The kernel verification in Corollary~\ref{cor:liminf-gauges} covers the two full alphabets used there.

\subsection{Quantitative finite-state gauges and general label presentations}

\begin{theorem}[Quantitative gauge bounds for a fixed start]\label{thm:finite-state-quant}
Let $C$ be a marked strongly connected graph as in
Theorem~\ref{thm:finite-state-tail}, with Perron root $\lambda>1$.
Let $w$ be a positive right Perron vector and set
$R=\max_v w_v/\min_v w_v$.
Start the recurrence at $N$ sufficiently large, with
$M_{N-1}=M_N=0$, and start the graph at a prescribed vertex $v$.
Let $K^*_{N,v}$ consist of its coded points with infinitely many
marked labels, projected to the circle. Then for every Hausdorff gauge $g$,
\[
 \frac{\rho_\lambda(g)}{3R\lambda^{N+1}}
 \ \le\ \mathcal H^g(K^*_{N,v})
 \ \le\ R\lambda^{-N}\rho_\lambda(g).
\]
The constants are uniform in $g,N,v$ apart from the displayed powers
of $\lambda$.
\end{theorem}

\begin{proof}
The initial pair makes every path realizable by the recurrence code
used in Theorem~\ref{thm:finite-state-tail}, and the entire construction
lies in an isometric circle chart. Each length $n-N$ prefix confines
its points to one interval of diameter $1/Q_n$. The Perron count
bounds the number of intervals by $R\lambda^{n-N}$.
Taking a lower limit along scales gives the upper bound.

Take the same Perron probability as in Theorem~\ref{thm:finite-state-tail}.
The code is a continuous injection from the compact path space for the
fixed initial vertex, so the Borel set of paths visiting a marked edge
infinitely often has Borel image.
The probability gives the marked recurrent paths full mass and pushes
forward to a probability on $K^*_{N,v}$. At any diameter $r$ with
$r_{n+1}\le r<r_n$, a covering set meets at most three depth $n-N$
cylinders, whose masses are each at most $R\lambda^{-(n-N)}$.
Given any finite $0<a<\rho_\lambda(g)$, for all sufficiently large
$n$ we have $g(r_{n+1})\ge a\lambda^{-(n+1)}$. The set's mass is
therefore at most $(3R\lambda^{N+1}/a)g(r)$.
Singletons have zero probability. Applying this bound to a countable
cover and taking the small-diameter limit gives the lower bound with
$a$ in place of $\rho_\lambda(g)$. Take the supremum over $a$;
when $\rho_\lambda(g)=0$ the lower bound is automatic.
\end{proof}

\begin{corollary}[Gauge algorithm for arbitrary labelled graphs]\label{cor:finite-state-nondet}
The conclusions of Theorem~\ref{thm:finite-state-tail} hold for an
arbitrary finite graph labelled by $\{-k,\ldots,k\}$ if one computes
the marked-component roots in the following finite subset graph.
Its vertices are nonempty subsets of the original vertices. For a
label $a$, there is an edge from $U$ to the set of all endpoints of
$a$-edges starting in $U$, whenever that endpoint set is nonempty.
Mark the edges labelled $\pm k$. In particular the result is
independent of the graph presentation of the eventual label language.
\end{corollary}

\begin{proof}
The subset graph has at most $2^{|V|}-1$ vertices and gives at most
one outgoing edge per label at every vertex. Every infinite word
of the original graph is read in the subset graph starting from
$V$. Conversely, an infinite word read from a subset has original
paths of every finite length starting in that subset. The tree of
those paths is finitely branching, so it has an infinite branch.
The two graphs thus have the same infinite label language and
the same eventual language. Infinitely many marked labels depend
only on the word. Apply Theorem~\ref{thm:finite-state-tail}.

For a base $\lambda_*>1$, one recovers it uniquely from the
transition for gauges $g_b(r_n)=b^{-n}$, $b>1$, with linear
interpolation. Indeed, $\rho_{\lambda_*}(g_b)$ equals $0$, $1$
or $\infty$ according as $b>\lambda_*$, $b=\lambda_*$ or
$b<\lambda_*$. The empty and countable cases are likewise
intrinsic, since a critical set with $\lambda_*>1$ has positive
measure for $g_{\lambda_*}$, which vanishes on countable sets.
\end{proof}

For example, a vertex with two loops labelled $k$ and one labelled
$-k$ has raw adjacency root $3$ and critical base $2$. The subset
construction here is the standard right-resolving idea for labelled
graphs (see Marcus~\cite[\S II-F]{MarcusSofic}). The arithmetic
part is the identification of marked recurrence with exact-level
membership and the comparison with gauges at the scales $1/Q_n$.
Theorem~\ref{thm:finite-state-quant} and
Corollary~\ref{cor:finite-state-nondet} have paper proofs.

\subsection{Frequency spectra for weighted finite-state constraints}
Let $G$ be a strongly connected finite graph with distinct outgoing
labels in $\{-k,\ldots,k\}$ at each vertex, and at least one extreme
label $\pm k$. For $p\in[0,1]$, let $S_{k,p}(G)$ consist of points of
the exact level $S_k$ whose eventual defect word is read in $G$ and
whose frequency of extreme labels tends to $p$. Exact-level membership
requires infinitely many extreme labels even when $p=0$.
Set $r_n=1/Q_n$, and for a positive Hausdorff gauge set
\[
 \alpha(g)=\lim_{n\to\infty}\frac{-\log g(r_n)}{n}
\]
when this limit exists, allowing $+\infty$.
Put $\chi(e)=1$ on an extreme-labelled edge and $0$ elsewhere, and set
\[
 A_t(v,u)=\sum_{e:v\to u}\exp(t\chi(e)),\qquad
 P(t)=\log\rho(A_t).
\]
Here $\rho$ is the Perron root. Use $p_t=P'(t)$ and
$h_t=P(t)-tp_t$.

\begin{theorem}[Frequency gauges from finite-state pressure]\label{thm:pressure-frequency}
For every $t\in\R$, the set $S_{k,p_t}(G)$ is nonempty.
If $\alpha(g)>h_t$, it is $\mathcal H^g$-null.
If $p_t\ne p_0$ and $\alpha(g)\le h_t$, its intersection with every
nonempty open arc is non-$\sigma$-finite for $\mathcal H^g$.
If $p_t=p_0$, then $h_t=P(0)$. For this frequency the full three-way
gauge law of Theorem~\ref{thm:finite-state-tail} holds, with
$\lambda_*=\exp(P(0))$ when $\lambda_*>1$.
When $\rho(A_0)=1$, the level at $p_0$ is countable.
\end{theorem}

\begin{proof}
Let $w_t>0$ be a right Perron vector, and from a vertex $v$ choose each
edge $e:v\to u$ with probability
\[
 K_t(e)=\exp(t\chi(e)-P(t))w_t(u)/w_t(v).
\]
This finite irreducible Markov chain has stationary law $\pi_t$ on
vertices. Differentiating the simple Perron eigenvalue with left and
right eigenvectors gives
$p_t=\sum_v\pi_t(v)\sum_{e:\,v\to u}K_t(e)\chi(e)$.
The conditional edge entropy is $h_t$ by telescoping the logarithms
of $w_t$. In particular $p_t>0$.

We record two counts for words of length $m$ beginning at a fixed
vertex. If $J$ counts extreme labels along a word, its probability is
$\exp(tJ-mP(t))w_t(v_m)/w_t(v_0)$. Thus the number of words with
$|J/m-p_t|\le\epsilon$ is at most
$C_t\exp(m(h_t+|t|\epsilon))$. To obtain many words of narrower
width, let $F$ solve the finite linear equation
$(I-K_t)F(v)=\mathbb E_t[\chi(e)\mid v]-p_t$, where $K_t$
acts on vertex functions by summing over outgoing edges.
The stationary mean of the right-hand side is zero, so
the simple eigenvalue $1$ gives solvability. The increments
$\chi(e_i)-p_t+F(v_{i+1})-F(v_i)$ are bounded martingale differences.
Consequently $\mathbb E_t[(J-mp_t)^2]=O_t(m)$, uniformly in
the initial vertex. Chebyshev's inequality gives probability
at least $1/2$ for $|J-mp_t|\le m^{2/3}$ at large $m$.
For the strong law, apply the same second-moment bound at lengths
$j^2$, use Borel--Cantelli, and interpolate between these lengths
using the bounded increments of $J$. In particular every fixed
starting vertex has almost surely the required frequency.

For the null conclusion, fix an eventual starting index and the
two starting integers of the recurrence code. A prefix of length
$m$ is contained in an interval of diameter $r_{N+m}$.
Choose $\epsilon$ small enough that
$h_t+|t|\epsilon<\alpha(g)$. The first count makes the cover cost
tend to zero, and a countable union over starts and onset indices
proves nullity.

The stationary relative entropy of $K_t$ with respect to $K_0$
equals $P(0)-h_t$. It is nonnegative, and equality implies equality
of all edge probabilities, hence $p_t=p_0$. If $p_t=p_0$, the
relative entropy in the reverse direction equals $h_t-P(0)$,
so $h_t=P(0)$ and one can use $K_0$.
For $\lambda_*=\rho(A_0)>1$, this distribution puts full mass on
the words of frequency $p_0$. Its cylinder estimate from
Theorem~\ref{thm:finite-state-quant} gives positive measure for
the reference gauge $g_{\lambda_*}(r_n)=\lambda_*^{-n}$.
The sublevel covering argument and adjacent-scale comparison in
Theorem~\ref{thm:finite-state-tail} apply to this subset. In the
positive finite normalized-cost regime, local infinitude follows
by the same invariance argument: the tail frequency is unchanged
by the dense arithmetic subgroup, and the level is Borel and
Lebesgue-null. The case $\lambda_*=1$ is covered by the
countability assertion there.

For $p_t\ne p_0$, use the block argument below with $p=p_t$,
$h=h_t$ and $h_0=P(0)>h_t$. To supply its loops, the narrow
window count just proved gives at least
$\exp(mh_t-O_t(m^{2/3}))$ words with $J=mp_t+O(m^{2/3})$.
Indeed, every such cylinder has probability at most
$C_t\exp(-mh_t+|t|m^{2/3})$. Choose an end vertex supporting
at least a fixed fraction of them and append a fixed return path.
Likewise the words at $t=0$ give return blocks at rate $h_0$.
Distinct words remain distinct after this operation because
outgoing labels are distinct.
\end{proof}

\paragraph{Block argument at the critical rate.}
Suppose two collections of loops based at a common vertex are
available at each large length parameter $m$, with a common
length $m+O(1)$ within each collection. Suppose the first has
at least $\exp(mh-O(m^{2/3}))$ distinct words of marked count
$mp+O(m^{2/3})$, and the second has at least
$\exp(mh_0-O(m^{2/3}))$ words of marked count
$mp_0+O(m^{2/3})$, where $h_0>h$ and $p_0>0$.
For a gauge with $\alpha(g)\le h$, define
\[
 u_n=\max(0,-\log g(r_n)-hn)=o(n),\qquad
 \epsilon_j=\max\!\left(j^{-1/8},
   \sqrt{\sup_{n\ge\max(1,\lfloor j^2/10\rfloor)}u_n/n}\right).
\]
The sequence $\epsilon_j$ decreases to zero. In block $j$, choose
a first type loop of parameter $j$, followed by a second type
loop of parameter $\lceil j\epsilon_j\rceil$, independently and
uniformly from their respective collections. The entire word has
frequency $p$, since the second type has total length $o(j^2)$
by stage $j$, and the errors in the first type total $O(j^{5/3})$.
An incomplete block has length $O(j)$ and does not change the
limit. Each sufficiently late second type loop has an extreme
label, so the exact-level condition holds for $p=0$ as well.

Start coding at a fixed large recurrence index $N$ with initial
integers zero and let $\mu$ be the image of the block probability
on the real-line chart. A prefix ending inside block $j$ has
length $m\asymp j^2$. Its probability is at most the product
of inverse word counts of all completed blocks. Taking logs,
the second type contributes an excess of
$(h_0-h)\sum_{i<j}\lceil i\epsilon_i\rceil\gg j^2\epsilon_j$
over rate $h$. The total errors $O(j^{5/3})$ and the possible
incomplete block cost $O(j)$ are $o(j^2\epsilon_j)$.
Thus for some $c>0$ the mass of any such prefix is at most
$\exp(-hm-cj^2\epsilon_j)$ for all sufficiently large $j$.

For a set of diameter $r_{n+1}\le r<r_n$, use the bound of
three cylinders of depth $m=n-N$ from
Theorem~\ref{thm:finite-state-quant}. At this depth $n+1\ge j^2/10$
and $n+1=O(j^2)$. Monotonicity gives
\[
 \frac{\mu(U)}{g(\operatorname{diam}U)}
 \ \le\ 3\exp\bigl(h(N+1)-cj^2\epsilon_j+u_{n+1}\bigr)
 \longrightarrow 0 ,
\]
since $u_{n+1}\le(n+1)\epsilon_j^2$. Singletons have zero mass
by the cylinder bound. The convergence is uniform as diameter
goes to zero. Hence every set of finite $\mathcal H^g$ measure
has zero $\mu$-outer measure, and a countable union of such
sets cannot cover the coded set. Taking $N$ as large as needed
makes the coded set lie in an arbitrarily small chart.
The dense group from Theorem~\ref{thm:finite-state-tail}
translates a copy into every nonempty open arc, preserving
eventual words. This proves local non-$\sigma$-finiteness.

\begin{corollary}[All extreme-label frequencies on the full alphabet]
\label{cor:full-alphabet-frequency}
For the full alphabet $\{-k,\ldots,k\}$, let $S_{k,p}$ denote
the corresponding exact-level frequency set. Then the threshold is
\[
 h(p)=-p\log p-(1-p)\log(1-p)
      +p\log 2+(1-p)\log(2k-1)
\]
with $0\log0=0$. If $\alpha(g)>h(p)$, the set is null.
If $p\ne 2/(2k+1)$ and $\alpha(g)\le h(p)$, it is locally
non-$\sigma$-finite. At $p=2/(2k+1)$ the arbitrary-gauge
law of Theorem~\ref{thm:pressure-frequency} applies. In particular,
the frequency-one exact level is locally non-$\sigma$-finite
for $g_2(r_n)=2^{-n}$, while the subfamily with eventually
only extreme defects has $\sigma$-finite $\mathcal H^{g_2}$.
\end{corollary}

\begin{proof}
The full alphabet is the one-vertex graph with
$P(t)=\log((2k-1)+2e^t)$. Its tilted probabilities give
the displayed $h(p)$ for $0<p<1$. For $p=0,1$, the number
of words with $J$ extreme labels out of $m$ is
$\binom mJ 2^J(2k-1)^{m-J}$. The binomial bound
$\binom mJ\le \exp(m[-r\log r-(1-r)\log(1-r)])$, $r=J/m$,
and continuity give the same null assertion at the endpoints.
For the lower assertion take as first type all loops of length $m$
with only unmarked labels at $p=0$, or only marked labels at $p=1$.
The second type comes from the maximal-entropy distribution.
The block argument applies even when $h(0)=0$.
The eventually-marked subfamily has the binary gauge law already
established by Theorem~\ref{thm:trichotomy}(ii).
\end{proof}

\begin{remark}[A constrained pressure calculation]
Let $k=1$ and forbid consecutive extreme labels. The weighted matrix is
$\left(\begin{smallmatrix}1&2e^t\\1&0\end{smallmatrix}\right)$.
The Perron root obeys $\lambda^2-\lambda-2e^t=0$. At $t=0$ one gets
$p_0=1/3$. At $t=\log3$ one gets $\lambda=3$, $p_t=2/5$ and
$h_t=(3/5)\log3$. The frequency-$2/5$ subset is locally
non-$\sigma$-finite for $g(r_n)=3^{-3n/5}$ and null for gauges
with a strictly greater exponential rate.
\end{remark}

The entropy formula is related to the classical digit-frequency
dimension calculation reviewed in Olsen~\cite[\S1.4]{OlsenFrequency}.
Elekes and Keleti~\cite[\S2]{ElekesKeletiFrequency} show a
null-or-non-$\sigma$-finite phenomenon for some base-expansion
frequency classes and translation-invariant measures. Here the scale
is the arithmetic sequence $1/Q_n$. The block estimate determines
the critical-rate side of the transition for finite-state exact
defect levels. These frequency results have paper proofs.


\subsection{Edge frequencies, endpoint rates and reducible graphs}

\begin{theorem}[Prescribed edge frequencies]\label{thm:edge-frequency-circulation}
Let $G$ be a strongly connected graph with distinct outgoing labels at
each vertex and at least one extreme edge. For a probability vector
$f=(f_e)_{e\in E(G)}$, define $S_k(G,f)$ by requiring an eventual
path with limiting edge frequencies $f$, as well as infinitely many
extreme labels. The set is nonempty exactly when $f$ is a circulation,
that is,
\[
 f_v:=\sum_{e:\,v\to u}f_e=\sum_{e:\,u\to v}f_e
 \quad(v\in V(G)).
\]
For such $f$ put
\[
 h(f)=-\sum_{e:\,v\to u}f_e\log(f_e/f_v),
\]
with zero terms omitted. If $\alpha(g)>h(f)$ the set is null.
If $f$ differs from the stationary edge law of the Perron chain
$K_0$, then it is locally non-$\sigma$-finite for
$\alpha(g)\le h(f)$. At the Perron edge law it inherits the
arbitrary-gauge classification of Theorem~\ref{thm:pressure-frequency}.
\end{theorem}

\begin{proof}
An empirical edge count has imbalance at most one at each vertex,
which proves necessity. For the upper bound, fix an initial vertex
and an empirical type $a_e/m$, where the nonnegative integers $a_e$
sum to $m$. For this type use the stochastic edge weights
$a_e/\sum_{e:\,v\to u}a_e$ out of each visited vertex $v$.
Every path of the type has the same positive probability
$\exp(-m h(a/m))$, so there are at most $\exp(m h(a/m))$ such paths.
There are at most $(m+1)^{|E|}$ types. Continuity of
$h(f)=-\sum_e f_e\log f_e+\sum_v f_v\log f_v$ and the cylinder
covers give nullity for $\alpha(g)>h(f)$, as in
Theorem~\ref{thm:pressure-frequency}.

For existence and the lower bounds, the positive support of a
circulation splits into strongly connected components with no
positive edge between them. Indeed a source in the component
condensation has zero outgoing mass by balance; remove it and
repeat. Within each support component of mass $s_i>0$, take the
irreducible chain with conditional probabilities $f_e/f_v$ and
stationary edge law $f_e/s_i$. For a length parameter $m$, sample
words of length $\lfloor s_i m\rfloor$ in that component.
Apply the martingale estimate in Theorem~\ref{thm:pressure-frequency}
to each of the finitely many edge indicators. With probability
bounded away from zero, all their counts are within $O(m^{2/3})$
of their means. A word of this kind has probability at most
$\exp(-\lfloor s_i m\rfloor h_i+O(m^{2/3}))$, where $h_i$
is the conditional entropy in that component. This follows by
writing its log probability as the sum of edge counts times
the logarithms of the positive transition probabilities.
Choose one end vertex with a fixed fraction of the good words
and join the component blocks with bounded connectors in $G$,
starting and ending at a chosen common vertex. For each parameter
$m$ the connectors are fixed across all word choices. The
result is a loop collection of length $m+O(1)$ with coordinate
counts $mf_e+O(m^{2/3})$ and cardinality
$\exp(mh(f)-O(m^{2/3}))$, since $h(f)=\sum_i s_i h_i$.

Let $h_0=P(0)$. The relative entropy of the edge transition
probabilities on the positive support with respect to $K_0$,
averaged with weights $f_v$, equals $h_0-h(f)$: balance cancels
the logarithms of Perron eigenvector coordinates. Equality
forces the two edge laws out of every vertex with $f_v>0$ to
agree. This support is then forward closed in $G$, hence all
of $G$ by strong connectivity, and stationarity forces
$f$ to be the Perron edge law.

For strict inequality $h(f)<h_0$, use those loop collections
as the first type in the critical-rate block argument after
Theorem~\ref{thm:pressure-frequency}, with the second type
coming from $K_0$. Frequency convergence there holds coordinate
by coordinate; the second type guarantees infinitely many
extreme labels even if the first type has none. This also
constructs a path at any non-Perron $f$, independently of the
gauge. At the Perron law each edge frequency converges almost
surely by the same martingale estimate. The Perron measure
cylinder bound and the proof for $p_0$ apply to this full-measure
frequency class. This proves the remaining gauge law and
existence.
\end{proof}

\begin{theorem}[Scalar endpoint entropy from tight cycles]
\label{thm:frequency-endpoints}
Assume $G$ is as in Theorem~\ref{thm:edge-frequency-circulation}.
Let $p_-,p_+$ be the minimum and maximum proportions of extreme
edges on its directed simple cycles. The set of attainable
extreme-label frequencies in $S_k(G)$ is $[p_-,p_+]$.
If $p_-<p_+$, then each interior $p$ is attained by a unique
tilt parameter $t$ with $P'(t)=p$, and has critical rate $h_t$.

For $p=p_+$ define the edge weight $w(e)=\chi(e)-p_+$, and for
$p=p_-$ define $w(e)=p_- -\chi(e)$. In either case put
\[
 F(v)=\max_{\pi:\,\mathrm{starts}(\pi)=v}\sum_{e\in\pi} w(e),
 \qquad
 G_p=\{\,e:v\to u\mid F(v)=w(e)+F(u)\,\}.
\]
The maximum includes the empty path. Let $\lambda_p$ be the maximum
Perron root of a strongly connected component of $G_p$ containing
a directed cycle. The sharp endpoint threshold is $\log\lambda_p$.
When $p_-<p_+$, the endpoint class is null for
$\alpha(g)>\log\lambda_p$ and locally non-$\sigma$-finite for
$\alpha(g)\le \log\lambda_p$. If all cycles have the same
proportion, that proportion has the original gauge law.
\end{theorem}

\begin{proof}
Removing directed cycles from a path leaves a simple path with at
most $|V|-1$ edges. Thus the count of extreme labels lies between
$mp_- - O(1)$ and $mp_++O(1)$ for a path of length $m$.
A cycle realizing an endpoint gives a periodic path with that
mean. For the pressure, summing $\exp(tJ)$ over words is
comparable to $\exp(mP(t))$ by a positive Perron vector.
The path bounds and a repeated extremal cycle give
$tp_+\le P(t)\le tp_++P(0)$ for $t\ge0$ and
$tp_-\le P(t)\le tp_-+P(0)$ for $t\le0$. In particular the
limits of the nondecreasing derivative $P'$ at the two
infinities are $p_-$ and $p_+$. Convexity here follows by
applying H\"older's inequality to the word sums and taking
their normalized logs. If $P'(t)=P'(s)$ at distinct parameters,
the two stationary relative entropies of their edge laws
sum to zero by telescoping. The laws agree, and comparing
their log probabilities along any cycle shows that all cycle
proportions equal $(P(t)-P(s))/(t-s)$. This proves the strict
interior assertion.

For an endpoint, every directed cycle has nonpositive $w$-weight,
so the maximum defining $F$ is obtained on simple paths.
The slacks $F(v)-w(e)-F(u)$ are nonnegative; a cycle has zero
total slack exactly when all its edges are tight. At least one
extremal cycle exists, so $\lambda_p\ge1$. If a path has limiting
proportion $p$, its empirical edge types have limit points that
are circulations supported on tight edges: for a circulation,
the sum of slacks equals minus the average of $w$, which is
zero at $p$.

For any circulation on $G_p$, its entropy is at most
$\log\lambda_p$. Decompose the support into strongly
connected components as in Theorem~\ref{thm:edge-frequency-circulation}
and compare each conditional law to the Perron law of the
corresponding positive-support component. Its Perron root is
at most that of the containing strongly connected component of
$G_p$, by monotonicity of spectral radius for nonnegative matrices. Conversely a Perron circulation in a component
with root $\lambda_p$ attains $\log\lambda_p$ and has marked
proportion $p$ by vanishing slack. Continuity and compactness
therefore bound $h(a/m)$ for all sufficiently long words
of proportion tending to $p$ by any number strictly larger
than $\log\lambda_p$, uniformly whenever the proportion lies
in a sufficiently small neighborhood of $p$. More explicitly,
a failure of this assertion gives empirical types along
$m\to\infty$ with proportions tending to $p$ and with entropy
bounded away above $\log\lambda_p$; any convergent subsequence
contradicts the preceding circulation bound. The type count
proves nullity.

When $p_-<p_+$ the extremal circulation cannot be the
Perron edge law of $G$: that law is positive on every edge,
whereas at least one slack is positive because a nonextremal
cycle has positive total slack. Theorem~\ref{thm:edge-frequency-circulation}
now gives the lower assertion and existence, including when
$p=0$ since its construction still inserts infinitely many
extreme labels. If all cycles have the same proportion, the
path bounds give that proportion on every path. It is positive
because there is an extreme edge in a cycle. The original
gauge law applies.
\end{proof}

\begin{remark}[Endpoint calculation with a transient tight edge]
For the graph forbidding consecutive extreme labels at $k=1$,
the feasible interval is $[0,1/2]$. At $p=0$, the tight graph has
the unmarked loop and the unmarked return to it; the loop alone
is recurrent. At $p=1/2$, only alternating paths are tight:
two marked transitions go out and one unmarked transition comes
back. The recurrent roots are $1$ and $\sqrt2$, giving endpoint
rates $0$ and $(\log 2)/2$. The finite telescoping identity for the
slacks is checked in \texttt{OrbitAndFlow.lean}.
\end{remark}

\begin{corollary}[Variational rate]\label{cor:variational-frequency-rate}
At every feasible $p$ in a strongly connected marked component, the
critical rate equals
\[
 H_G(p)=\max\{h(f): f\text{ is a probability circulation},\
                      \sum_e f_e\chi(e)=p\}.
\]
\end{corollary}
\begin{proof}
For an interior frequency, the stationary relative entropy of
$f_e/f_v$ with respect to $K_t$ gives
$h(f)\le P(t)-tp$, with equality at the stationary edge law of
$K_t$. At the endpoints use the tight-graph maximizer in
Theorem~\ref{thm:frequency-endpoints}. For constant proportions
the unrestricted Perron law maximizes the entropy.
\end{proof}

\begin{corollary}[Scalar frequencies in reducible labelled graphs]
\label{cor:reducible-frequency-rates}
For a general right-resolving finite graph, retain the strongly
connected components having an internal extreme edge and for
which $p$ lies in the interval from
Theorem~\ref{thm:frequency-endpoints}. If none remain, the
exact level of frequency $p$ is empty. Otherwise put
$H_*=\max_C H_C(p)$ over those components.
The level is null for $\alpha(g)>H_*$ and locally
non-$\sigma$-finite for $\alpha(g)<H_*$. At $\alpha(g)=H_*$,
if a maximizing component has $p\ne p_0(C)$, the level is
locally non-$\sigma$-finite. If every maximizing component
has $p=p_0(C)$, use the arbitrary-gauge law at base
$\exp(H_*)$ for $H_*>0$; for $H_*=0$ the level is countable.
For an arbitrary label presentation the scalar statement uses
the subset graph from Corollary~\ref{cor:finite-state-nondet}.
\end{corollary}
\begin{proof}
Every infinite path eventually remains inside one strongly
connected component. Moving the starting point by finitely many
positions preserves frequency and infinite recurrence of an
extreme label. The scalar level is consequently the finite
union of the component levels. At criticality all components
of smaller entropy are null; among the maximizing components
one applies their critical gauge laws. Countable unions of
$\sigma$-finite sets are $\sigma$-finite, and a single locally
non-$\sigma$-finite subset suffices for the opposite conclusion.
For arbitrary label presentations the deterministic subset
construction preserves the full eventual label language and
therefore the scalar label frequencies.
\end{proof}

The flow description is related to the rotation sets of finite
graphs \cite{ZiemianRotation}: edge counts determine balanced empirical vectors.
The related zero--infinity result for generic points in Gibbs
symbolic metrics \cite{GenericZeroInfinity} concerns a different
metric setting. Here the recurrence-based scales and the
high-entropy insertion estimate give the stated result for
arithmetic gauges, including endpoint frequencies and reducible
labelled graphs. The new results in this subsection are paper proofs.


\subsection{Frequency levels with a bound on discrepancy}

Fix a right-resolving strongly connected marked graph $G$ and
let $p=c/d$ be in $(0,1)$, written in lowest terms with $d>0$.
For the first $m$ edges of a path let $J_m$ count its marked edges.
Given an integer $H\ge0$, require that some eventual tail of the
defect code has a path for which
\[
 |z+dJ_m-cm|\le H\quad(m\ge0)
\]
for some integer initial height $z\in[-H,H]$. Denote the resulting
subset of $S_k$ by $S_{p,H}(G)$. Allowing the initial height to vary
makes this condition stable under taking a later tail.

\begin{theorem}[Finite strips and their entropy limit]\label{thm:frequency-strip}
Form a graph $G_{p,H}$ on $(v,z)$ with $v\in V(G)$ and
$z\in[-H,H]\cap\Z$. Each original edge $e:v\to u$ induces the edge
\[
 (v,z)\longrightarrow(u,z+d\chi(e)-c)
\]
with the original label if the new height lies in the same interval.
If there is no directed cycle, $S_{p,H}(G)$ is empty. Otherwise let
$\lambda_H$ be the largest Perron root of a component containing
a cycle. Then $S_{p,H}(G)$ is countable if $\lambda_H=1$.
For $\lambda_H>1$ it has the arbitrary-gauge law of
Theorem~\ref{thm:finite-state-tail}, with base $\lambda_H$.
If $p$ is strictly between the smallest and largest cycle means of $G$,
then $\lambda_H$ is eventually defined and
\[
 \lambda_H\uparrow \exp(H_G(p))\qquad(H\to\infty),
\]
where $H_G(p)$ is the variational frequency rate from
Corollary~\ref{cor:variational-frequency-rate}.
\end{theorem}
\begin{proof}
The sequence of visited heights is exactly the accumulated
integer discrepancy plus $z$. The graph is right-resolving.
On any directed cycle of positive length in the strip,
$dJ_m=cm$, so the cycle has a marked edge. Thus the strip's
components containing a cycle are exactly its marked components,
and Theorem~\ref{thm:finite-state-tail} gives the measure statements.
Monotonicity follows by inclusion of strip matrices in the larger
matrices. To bound $\lambda_H$ above, use the tilted path counts
from Theorem~\ref{thm:pressure-frequency}: all strip paths of length
$m$ have $|dJ_m-cm|\le2H$, so the number of label words from each
initial vertex is at most a constant depending on $H$ times
$\exp(mH_G(p))$. This bounds the Perron root.

For the lower bound, fix a base vertex $v$ of $G$ and use the
tilt at the interior frequency $p$. The bounded-variance estimate
and the path probability formula in
Theorem~\ref{thm:pressure-frequency} give at least
$\exp(nH_G(p)-O(n^{2/3}))$ length-$n$ paths from $v$ with
$|dJ_n-cn|=O(n^{2/3})$.
Partition these paths by their endpoint and their integer displacement.
There are $O(n^{2/3}+1)$ classes. Keep a largest class and append
the same bounded-length path back to $v$. This gives a family of
$M_n$ distinct equal-length loops, with
$\log M_n\ge nH_G(p)-O(n^{2/3})$ and common displacement $s_n=O(n^{2/3})$.

There exist fixed based loops of displacements $a>0$ and $-b<0$.
To construct them, go from $v$ to cycles on either side of $p$,
repeat the cycle enough times, and return to $v$.
Write $g_0=\gcd(a,b)$. Concatenate $g_0$ arbitrary loops from the
family and then append copies of the two fixed loops to cancel
the displacement $g_0s_n$. There are nonnegative integers $u,w$ with
$ua-wb=-g_0s_n$ and $u+w=O(|s_n|+1)$: start with a B\'ezout
solution and add a common multiple of $(b/g_0,a/g_0)$ to
make both entries nonnegative. The correction word is fixed
for this $n$. The resulting $M_n^{g_0}$ distinct loops all have
zero displacement and common length $L_n=g_0n+O(n^{2/3})$.
Each fits in the strip with initial height zero and width
$H=dL_n$, since each edge changes the height by at most $d$.
Concatenating the loops in that strip bounds its Perron root
below by $M_n^{g_0/L_n}$. Letting $n$ tend to infinity and using
monotonicity proves the claimed limit.
\end{proof}

\begin{corollary}[The full alphabet at frequency one half]\label{cor:half-strip}
For the full defect alphabet $\{-k,\ldots,k\}$ and $p=1/2$, set
$a=2$ and $b=2k-1$. The width-zero strip is empty and for $H\ge1$
\[
 \lambda_H=2\sqrt{ab}\cos\frac{\pi}{2H+2}.
\]
The loss in logarithmic rate relative to the unrestricted frequency
layer is $\pi^2/(8(H+1)^2)+O((H+1)^{-4})$.
\end{corollary}
\begin{proof}
A marked label moves the height up by one with multiplicity $a$;
an unmarked label moves it down by one with multiplicity $b$.
Conjugation by the diagonal matrix with consecutive ratios
$\sqrt{b/a}$ reduces the adjacency matrix to the symmetric path
matrix with off-diagonal entries $\sqrt{ab}$. There are $2H+1$
vertices. The sine vector with consecutive phases $\pi/(2H+2)$
is positive on those vertices and vanishes at the two exterior
indices; applying the sine addition formula gives the eigenvalue.
The unrestricted base is $2\sqrt{ab}$ by
Theorem~\ref{thm:frequency-strip}. Expanding $-\log\cos$ gives
the stated loss.
\end{proof}

\begin{corollary}[A speed restriction with unchanged exponential rate]
\label{cor:sublinear-speed}
Under the interior hypotheses of Theorem~\ref{thm:frequency-strip},
let $R_m\to\infty$ with $R_m/m\to0$. Require
$|J_m-mp|\le R_m$ eventually along a defect tail.
The resulting level is null when $\alpha(g)>H_G(p)$ and locally
non-$\sigma$-finite when $\alpha(g)<H_G(p)$.
For the full alphabet at $p=2/(2k+1)$, if also $R_m=o(\sqrt m)$,
the speed-restricted level is null for the reference gauge
$h_{2k+1}$, although the unrestricted frequency level is $\sigma$-finite
and has infinite $h_{2k+1}$ measure on every nonempty open arc.
\end{corollary}
\begin{proof}
The speed condition implies the frequency $p$ and contains
every fixed strip level, whose unshifted discrepancy is at most
$2H/d$. The upper bound therefore follows from the unrestricted
level. For $\alpha(g)<H_G(p)$ choose a strip with
$\log\lambda_H>\alpha(g)$. This forces
$\liminf\lambda_H^n g(r_n)=+\infty$ and gives the lower conclusion
by the arbitrary-gauge law.

Put $\lambda=2k+1$ for the full alphabet.
For any fixed starting index, initial integer pair and large $m$,
the eligible prefixes of length $m$ are at most $\lambda^m$ times
the probability that a $\operatorname{Bin}(m,2/\lambda)$ variable
lies within $R_m$ of its mean. By Stirling's bounds the maximum
point probability is $O_k(m^{-1/2})$, so the probability is
$O_k((1+R_m)/\sqrt m)=o(1)$.
Each prefix at start $N$ is covered at scale $r_{N+m}$ and has
$h_\lambda$ cost at most $\lambda^{-(N+m)}$.
Thus the total cost tends to zero. Take the countable union over
initial integer pairs, starting times and thresholds for the
eventual speed condition. At this frequency the unrestricted level
has the Perron frequency and obeys the finite-$\rho_\lambda$ critical
law of Theorem~\ref{thm:pressure-frequency}.
\end{proof}
The displacement identity for the strip lift is audited
in \texttt{ResearchAdvance.lean}. The Perron and measure arguments
are paper proofs. The coefficient checks for the path matrix and
the binary-alphabet calculation are in
\texttt{verify\_research\_advance.py}.


\begin{theorem}[A sharp second-order boundary at the reference gauge]
\label{thm:frequency-upper-class}
For the full alphabet let $\lambda=2k+1$, $p=2/\lambda$ and
$\sigma^2=p(1-p)$. Fix $b\in\R$. At large $m$ put
\[
 R^{(b)}_m=\sigma\sqrt{m\bigl(2\log\log m+
                                      b\log\log\log m\bigr)}.
\]
Require $|J_m-mp|\le R^{(b)}_m$ for all sufficiently large $m$
along an eventual defect tail. The resulting level is
$\mathcal H^{h_\lambda}$-null for $b\le3$.
For $b>3$ it is $\sigma$-finite and has infinite
$\mathcal H^{h_\lambda}$ measure in every nonempty open arc.
The logarithmic rate is $\log\lambda$ for every $b$.
\end{theorem}
\begin{proof}
For uniformly chosen labels, the centered mark indicators are
independent, bounded, mean-zero variables with variance $\sigma^2$.
Set $\psi_m=R^{(b)}_m/(\sigma\sqrt m)$.
Feller's upper/lower-class test \cite[Theorem~2]{Feller1943}
applies on both signs: the ratio of the increment bound to
$\sigma\sqrt m$ is $O(m^{-1/2})=O(\psi_m^{-3})$, and
$\psi_m$ is eventually increasing. Its series is, up to a factor
tending to a positive constant,
\[
 \sum_m\frac{\psi_m}{m}\exp(-\psi_m^2/2)
 \ \asymp\
 \sum_m\frac{1}{m\log m\,(\log\log m)^{(b-1)/2}},
\]
which converges exactly for $b>3$. A finite initial segment has no
effect on the test.

We transfer the Bernoulli classification to the arithmetic metric.
For a fixed starting time $N$ and integer pair in the recurrence
code, a word of length $m$ has product probability $\lambda^{-m}$
and its image is covered at diameter $r_{N+m}$ for
$h_\lambda$-cost at most $\lambda^{-(N+m)}$.
A Bernoulli-null Borel set of words is covered by open sets of
arbitrarily small probability. Decompose each open set into
disjoint prefix cylinders, refining to any prescribed minimum depth.
The total cost of their images is at most $\lambda^{-N}$ times the
probability of the open set. Thus its image is $h_\lambda$-null.
For $b\le3$ the eventual bound is Bernoulli-null by the divergent
case of the test. A countable union over starts and integer pairs
gives the first assertion.

For $b>3$ use the set of points whose defect word satisfies
\[
 \limsup_{m\to\infty}
 \frac{(J_m-mp)^2/(m\sigma^2)-2\log\log m}{\log\log\log m}=3 .
\]
This set has full Bernoulli probability for a fixed start, by
applying the test at rational parameters approaching $3$ from
above and below. The fixed-start mass-distribution estimate gives
positive $h_\lambda$ measure for its image. Dropping a finite
prefix changes the discrepancy by a constant and shifts $m$ by
a fixed integer. Under the displayed growth bound this changes
the normalized expression by $o(1)$, so the subset is stable under
later tails. It is Borel and invariant under the dense translation
group $L$, by the eventual equality of defect tails used in
Theorem~\ref{thm:finite-state-tail}. The same finite-measure
invariance argument gives infinite measure in every nonempty
open arc. This subset obeys the eventual bound at parameter
$b>3$. The unrestricted finite-rate law supplies $\sigma$-finiteness.
Finally $R^{(b)}_m\to\infty$ and $R^{(b)}_m/m\to0$, so
Corollary~\ref{cor:sublinear-speed} gives the rate statement.
\end{proof}


\subsection{Upper and lower classes for finite-state tails}

Let $G$ be a strongly connected labelled graph as in
Theorem~\ref{thm:finite-state-tail}, with a marked edge and Perron
root $\lambda>1$. For an edge $e:v\to u$, let
$K(e)=w_u/(\lambda w_v)$ for a positive right Perron vector $w$.
Let $\nu$ be the stationary vertex law, $\chi(e)$ the mark indicator,
and $p=\sum_v\nu_v\sum_{e:v\to u}K(e)\chi(e)>0$.
Put $a(e)=\chi(e)-p$ and write $Z_m=J_m-mp$.
All the statements below use $h_\lambda(r_n)=\lambda^{-n}$.

\begin{proposition}[Variance and bounded discrepancy]
\label{prop:markov-discrepancy-variance}
Fix a vertex $v_*$ and let $(T,Y)$ be the length and the sum of
$a(e)$ along the first return excursion under $K$.
Then $\mathbb EY=0$, and
\[
 \sigma^2=\frac{\mathbb EY^2}{\mathbb ET}.
\]
This quantity is independent of $v_*$, equals the pressure curvature
$P''(0)$, and has the finite formula
\[
 \sigma^2=\sum_v\nu_v\sum_{e:v\to u}
                 K(e)\bigl(a(e)+F(u)-F(v)\bigr)^2,
 \qquad (I-K)F(v)=\sum_{e:v\to u}K(e)a(e),
\]
where $K$ on vertex functions sums over outgoing edges.
The variance is zero if and only if every directed cycle has mark
mean $p$. In that case $Z_m$ is uniformly bounded over all finite
paths, and every diverging tolerance eventually holds on
$S_k(G)$.
\end{proposition}
\begin{proof}
From each vertex there is a path to $v_*$ of bounded length
and probability uniformly bounded away from zero. Consequently
the first return has an exponential tail, and $|Y|\le T$.
The expected time spent at each vertex before the return,
divided by $\mathbb ET$, is stationary: expected arrivals and
departures agree because the walk starts and ends at $v_*$.
Uniqueness identifies this occupation law as $\nu$, and summing
conditional expected edge rewards gives $\mathbb EY=0$.

The right side of the Poisson equation has stationary mean zero,
so the finite linear equation is solvable, with $F$ determined
up to a constant. The adjusted edge increments
$a(e)+F(u)-F(v)$ are bounded martingale differences.
Their sum during the return excursion is $Y$ by telescoping.
Stopping first at $T\wedge n$ gives the second moment as the
expected sum of conditional increment variances. Exponential
tails justify passing to $T$, and the occupation formula gives
the displayed identity.

To identify the pressure curvature, let $\mathsf M_t(v,u)$
sum $K(e)\exp(t a(e))$ over $e:v\to u$.
Its Perron root is $\exp(P(t)-P(0)-tp)$.
Normalizing the right eigenvector to equal $\mathbf 1$ at
$t=0$, its derivative satisfies the Poisson equation for $F$,
since $P'(0)=p$. Twice differentiating the eigenvalue
equation and averaging with $\nu$ gives
$P''(0)=\mathbb E_\nu[a(e)^2+2a(e)F(u)]$.
By stationarity,
$\mathbb E_\nu[(F(u)-F(v))^2]=2\mathbb E_\nu[F(v)(I-K)F(v)]$.
Expanding the square in the variance formula and applying the
Poisson equation gives the same expression.
The variance terms are nonnegative and every edge has positive
weight. Hence $\sigma^2=0$ exactly when
$a(e)=F(v)-F(u)$ on all edges, which implies the cycle condition
and the uniform bound. Conversely, when every cycle sum is zero,
the sum from $v_*$ to a vertex is independent of path by appending
a fixed return path. This gives a potential with zero adjusted
increments (with the sign chosen as in the formula).
\end{proof}

\begin{theorem}[Sharp higher-order bounds at the Perron frequency]
\label{thm:markov-upper-class}
Assume $\sigma^2>0$. Put $L_i(m)=\log^{\circ i}(m)$ at sufficiently
large $m$, and for a fixed integer $\ell\ge3$ define
\[
 A_3(m)=2L_2(m),\qquad
 A_\ell(m)=2L_2(m)+3L_3(m)+2\sum_{i=4}^{\ell-1}L_i(m)
       \quad(\ell\ge4).
\]
For $b\in\R$, impose along an eventual defect tail the bound
\[
 |Z_m|\le \sigma\sqrt{m\bigl(A_\ell(m)+bL_\ell(m)\bigr)}
                 \qquad\text{for all sufficiently large }m.
\]
At gauge $h_\lambda$ this set is null for $b\le b_\ell$ and
$\sigma$-finite with infinite measure in every nonempty open
arc for $b>b_\ell$, where $b_3=3$ and $b_\ell=2$ for $\ell\ge4$.
The equality cases are included in the null phase. The conclusion
holds for periodic as well as aperiodic $G$.
\end{theorem}
\begin{proof}
Start the Perron chain at $v_*$ and decompose it at returns.
The pairs $(T_i,Y_i)$ are independent and identically distributed.
Write $\mu=\mathbb ET$, $\beta^2=\mathbb EY^2>0$,
$S_n=\sum_{i\le n}Y_i$ and $t_n=\sum_{i\le n}T_i$.
Here $\beta^2=\mu\sigma^2$ and $Z_{t_n}=S_n$.
Exponential tails, the exponential Markov inequality and
Borel--Cantelli give almost surely
$t_n=\mu n+O(n^{2/3})$ and $T_{n+1}=O(\log(n+2))$.
For instance the centered lengths have moment generating
function bounded by $\exp(Cs^2)$ near zero, so the two
tail probabilities at $n^{2/3}$ are summable.

We apply Feller's test \cite[Theorem~2]{Feller1943} to independent
bounded variables obtained from $Y_i$ by truncation at
$c\log(i+2)$ and centering. Choose $c$ so large that the removed
tail probabilities and removed first and second moments are
summable. The new partial sums differ from $S_n$ by $O(1)$
almost surely, and their variance sum is
$s_n^2=n\beta^2+O(1)$. Their increment bound divided by $s_n$
is $O(\log n/\sqrt n)$. For either sign and every fixed real $q$,
take the eventual nondecreasing sequence
\[
 \psi_n^2=A_\ell(\mu n)+bL_\ell(\mu n)+q .
\]
The increment bound is $O(\psi_n^{-3})$, as required by
Feller's test. Its series is equivalent for convergence to
$\sum_n n^{-1}\psi_n e^{-\psi_n^2/2}$. For $\ell=3$ this is
comparable to
$\sum_n [n\log n\,L_2(n)^{(b-1)/2}]^{-1}$.
For $\ell\ge4$ it is comparable to
\[
 \sum_n \frac{1}{n L_1(n)\cdots L_{\ell-2}(n)
                                L_{\ell-1}(n)^{b/2}}.
\]
By substitution in the integral test these converge exactly
for $b>b_\ell$. The same conclusion holds for either sign and
is independent of $q$.

We detail the transfer of the equality cases. Adding a fixed
positive constant under the squared normalized boundary gives
a margin of order $\sqrt{n/L_2(n)}$ for the sums.
Replacing $\mu n$ by $t_n$ in the function
$\sigma\sqrt{t(A_\ell(t)+bL_\ell(t)+q)}$ changes the boundary by
$O(n^{1/6}\sqrt{L_2(n)})$, by the mean value theorem.
Replacing $\beta\sqrt n$ by $s_n$ introduces a smaller error.
Between adjacent returns the change in $Z_m$ is $O(\log n)$,
and the boundary changes by $o(\sqrt{n/L_2(n)})$ as well.
All these errors are smaller than the fixed-constant margin.
Use $q=-1$ in the convergent part to obtain the eventual bound
at all times with $q=0$. Use $q=1$ in the divergent part to get
infinitely many violations at returns with $q=0$.
Any other initial vertex hits $v_*$ almost surely; the finite
initial offset and time shift are absorbed by the same margin.

We transfer the probability statement to the arithmetic level.
For a fixed starting time $N$, vertex and starting integer pair,
the probability of a length-$m$ cylinder is
$w_u/(\lambda^m w_v)$, where $u$ is the endpoint.
This is bounded above and below by fixed multiples of
$\lambda^{-m}$. Its image has a cover of diameter $r_{N+m}$
with cost $\lambda^{-(N+m)}$ for the reference gauge.
A probability-zero Borel set is covered by open sets of arbitrarily
small probability. Partition each open set into disjoint cylinders
at an arbitrarily large minimum depth. The image cost is bounded
by a fixed multiple of the probability. Taking the countable union
over starts and initial pairs proves nullity for $b\le b_\ell$.

In the positive phase consider the subset with
\[
 \limsup_{m\to\infty}
 \frac{Z_m^2/(m\sigma^2)-A_\ell(m)}{L_\ell(m)}=b_\ell .
\]
The two probability phases, applied at countably many parameters
approaching $b_\ell$, give probability one for this equality.
Using a fixed start and the mass-distribution bound from
Theorem~\ref{thm:finite-state-quant}, its image has positive
$h_\lambda$ measure. The equality is stable under dropping a
finite prefix: it implies $Z_m=O(\sqrt{mL_2(m)})$, so a fixed
time shift and an additive discrepancy change alter the
normalized expression by $o(1)$. Thus the corresponding set
of arithmetic points is Borel and invariant under the dense group
$L$ of tail-preserving translations. It is Lebesgue-null.
The finite-measure invariance argument of
Theorem~\ref{thm:finite-state-tail} then gives infinite measure
in every nonempty open arc. For $b>b_\ell$ this subset lies in
the speed-restricted set, whose $\sigma$-finiteness follows
from that of $S_k(G)$.
\end{proof}

\begin{remark}[Two dependent examples]
Take a graph with a loop at $0$ and edges $0\to1$ and $1\to0$;
mark only $0\to1$. Its Perron root is $\varphi$.
Set $q=\varphi^{-2}$. Return excursions from $0$ have
lengths $1,2$ with probabilities $1-q,q$, and mark counts
$0,1$, respectively. Hence
$p=q/(1+q)$ and $\sigma^2=q(1-q)/(1+q)^3$.
Both the second-order and all the stated higher-order
boundaries apply to this dependent label process.
For a periodic zero-variance example, use two marked parallel
edges from $0$ to $1$ and two unmarked parallel edges from $1$
to $0$, with distinct labels at each vertex and $k\ge2$.
Then $\lambda=2$, $p=1/2$ and every discrepancy
has magnitude at most $1/2$.
\end{remark}

Potential telescoping and the two-excursion variance identity
are checked in \texttt{TraceMarkov.lean}. The certificate
\texttt{verify\_trace\_markov.py} checks the two example variances
with exact arithmetic. The regeneration, probability test, and
Hausdorff measure argument are paper proofs.


\subsection{Discrepancy constraints away from the Perron frequency}

\begin{lemma}[Uniform tilted fluctuation estimates]
\label{lem:tilted-clt}
In the notation of Theorem~\ref{thm:pressure-frequency},
take $p=P'(t)$ and use the law $K_t$ on paths from a prescribed
vertex. Suppose the graph has directed cycles with distinct
mark means. For $Z_j=J_j-jp$, the law of $Z_j/\sqrt j$ tends
to $N(0,P''(t))$, uniformly over the initial vertex, with
$P''(t)>0$. There is also a uniform maximal bound
\[
 \Pr_v\!\left(\max_{j\le L}|Z_j|>s\right)\le C_t L/s^2
 \qquad(s\ge C_t,\ L\ge1).
\]
\end{lemma}
\begin{proof}
The variance and Poisson calculation of
Proposition~\ref{prop:markov-discrepancy-variance} applies with
$K_t$ in place of $K$: the centered moment matrix at parameter $z$
has Perron root $\exp(P(t+z)-P(t)-pz)$.
The formula for its second derivative gives the same sum of squares
under the stationary law of $K_t$. Its vanishing would force all
cycle means to be $p$. The martingale maximal inequality for
the bounded Poisson-corrected increments gives the stated bound,
absorbing the bounded potential in $s$.

For the central limit assertion write this moment matrix as $M_z$.
It has a positive right Perron vector $w_{t+z}/w_t$, with
coordinates tending to $1$ as $z\to0$, using continuously
normalized Perron vectors. Comparison with this vector above
and below bounds $(M_z^j\mathbf 1)(v)$ by the $j$th power of the
Perron root times factors tending to $1$ uniformly in $v$.
Taking $z=u/\sqrt j$ and Taylor expanding $P$ yields
$\exp(u^2P''(t)/2)$ for the limiting moment generating function.
The moment-generating-function convergence theorem
\cite{CurtissMGF} gives the normal limit. This comparison needs
no aperiodicity assumption.
\end{proof}

\begin{theorem}[Cubic-root discrepancy threshold]
\label{thm:cubic-root-threshold}
Let $G$ be strongly connected and right resolving, with a marked
edge, and suppose its cycle mark means are not all equal.
Fix $p=P'(t)\ne P'(0)$ for finite $t$ and set
$h=P(t)-tp$. At this frequency let $g(r_n)=\exp(-hn)$,
with monotone interpolation between the scales.
Let $D_{\alpha,C}$ be the part of the exact level read in $G$
at frequency $p$ for which, along some eventual tail,
$|J_m-mp|\le C m^\alpha$ for all sufficiently large $m$,
where $C>0$ and $0<\alpha<1$.
Then:
\begin{enumerate}
\item If $\alpha<1/3$, $D_{\alpha,C}$ is $\mathcal H^g$-null
for every fixed $C$.
\item If $\alpha>1/3$, $D_{\alpha,C}$ is locally
non-$\sigma$-finite for $\mathcal H^g$ for every fixed $C$.
\item At exponent $1/3$ the elementary argument below gives nullity
for small $C$ and local non-$\sigma$-finiteness for large $C$.
Theorem~\ref{thm:sharp-cubic-amplitude} gives the exact
critical amplitude.
\end{enumerate}
The constants in the last assertion may depend on $G$ and $p$.
\end{theorem}
\begin{proof}
At this tilt, a length-$m$ cylinder has probability
\begin{equation}
 \exp(-hm+tZ_m)\frac{w_t(v_m)}{w_t(v_0)} .
 \label{eq:cubic-tilted-weight}
\end{equation}
Also $h>0$: the graph has a branching vertex, since a strongly
connected graph with only one outgoing edge per vertex is one
cycle. At a branching vertex the tilted edge law has positive
entropy, and its stationary weight is positive.

\emph{Upper bound.}
For a fixed onset of the discrepancy condition and large $m$,
look between times $\lfloor m/2\rfloor$ and $m$.
Any admissible word stays in $[-Cm^\alpha,Cm^\alpha]$ there.
Divide the interval into complete blocks of length
$L=\lceil K C^2m^{2\alpha}\rceil$. Choose $K$ large enough that,
uniformly over each block's initial vertex and large $m$, the
probability that the centered increment over the block has absolute
value at most $2Cm^\alpha$ is bounded by a fixed $\theta<1$.
This follows from Lemma~\ref{lem:tilted-clt}, since the endpoint
threshold divided by $\sqrt L$ is at most $2/\sqrt K$.
Admissibility implies that bound on every complete block increment.
The Markov property gives probability at most
$\exp(-\kappa C^{-2}m^{1-2\alpha})$ for some $\kappa>0$,
when $\alpha<1/2$ and $m$ is sufficiently large for the fixed
$C$. The constant $\kappa$ can be chosen independently of $C$.
By \eqref{eq:cubic-tilted-weight}, each eligible word has
probability bounded below by a fixed multiple of
$\exp(-hm-|t|Cm^\alpha)$. Right resolution identifies words
with paths at a prescribed start. The $g$-cost of covering their
images at diameter $r_{N+m}$, for a fixed recurrence start $N$
and pair of initial integers, is consequently at most
\[
 O_N(1)\exp\!\left(|t|Cm^\alpha
                    -\kappa C^{-2}m^{1-2\alpha}\right).
\]
This tends to zero for $\alpha<1/3$, and also for $\alpha=1/3$
when $C$ is sufficiently small. Take the countable union over
starts, starting data and onsets.

\emph{Lower bound.}
It suffices to treat $\alpha<1/2$,
since for any larger $\alpha$ the construction at a smaller
exponent above $1/3$ gives a subset. Write
$V_m=-\operatorname{sgn}(t)Z_m$ and $R(n)=Cn^\alpha$.
Use a large virtual time $n_0$ and the shifted coordinate
$\widetilde V_{n_0+m}=V_m+R(n_0)/2$.
Starting at $n=n_0$, take consecutive blocks whose length is
$L_n=\lfloor\delta R(n)^2\rfloor$, for a small fixed
$\delta>0$. At time $n$ require $\widetilde V_n$ to lie in
the inner interval
$[3R(n)/8,5R(n)/8]$. The requirement at the end of the block
is membership in the corresponding inner interval at $n+L_n$;
also require the absolute centered increment throughout the
block to be at most $R(n)/8$.

There is a lower bound $a_*>0$ for the probability of this
block event, uniform over initial vertices and inner positions
at large $n$ and independent of $C$. In fact, the two terminal
increment intervals $[\sqrt{L_n},2\sqrt{L_n}]$ and
$[-2\sqrt{L_n},-\sqrt{L_n}]$ have probabilities bounded below
by the normal limit in Lemma~\ref{lem:tilted-clt}.
Choose $\delta$ so $2\sqrt{L_n}\le R(n)/16$.
As $\alpha<1/2$, the change $R(n+L_n)-R(n)$ is $o(\sqrt{L_n})$.
If an inner starting position is above the center $R(n)/2$,
the negative interval lands inside the next inner interval;
if it is below the center, the positive interval does.
The maximal bound in Lemma~\ref{lem:tilted-clt} makes the
probability of violating the within-block increment condition
at most $O(\delta)$. A smaller $\delta$ makes this less than
half the common terminal probability bound.
Thus the law $K_t$ conditioned successively on these block
events is well defined from any specified starting vertex.

Every path in its support satisfies
$R(n)/4\le\widetilde V_j\le3R(n)/4$ inside that block.
Since $L_n=o(n)$, after undoing the initial shift it follows
that all large actual times $m$ satisfy
$Cm^\alpha/6\le V_m\le 7Cm^\alpha/8$.
The number of blocks met up to time $m$ is at most
$O(1)+O(C^{-2}m^{1-2\alpha})$, with the second constant
independent of $C$. To see this, each completed block
contributes a fixed lower bound, proportional to $\delta C^2$,
to the integral of $u^{-2\alpha}$ across it.
Conditioning multiplies cylinder probabilities by at most
$a_*^{-1}$ per completed block and once more for an incomplete
block. Hence the new cylinder probabilities, by
\eqref{eq:cubic-tilted-weight}, are at most
\[
 O(1)\exp\!\left(-hm-\frac{|t|C}{6}m^\alpha+
                       B C^{-2}m^{1-2\alpha}\right),
\]
with $B$ independent of $C$ for fixed $\alpha$.
This is $o(e^{-hm})$ uniformly in the words if $\alpha>1/3$,
and also at $\alpha=1/3$ for sufficiently large $C$.

Map the path law into the recurrence code at a sufficiently
late fixed start. This law is supported on the required
frequency, which is positive. A small set of diameter
$r_{n+1}\le r<r_n$ meets at most three prefixes of depth
$n-N$, as in Theorem~\ref{thm:finite-state-quant}.
Thus its mass divided by $g(r)$ tends to zero uniformly,
so every set of finite $\mathcal H^g$ measure has zero
outer mass. The image cannot be $\sigma$-finite.
It can be placed in arbitrarily small charts by taking the
recurrence start late. The dense arithmetic translation
group moves a copy into every nonempty open arc. Dropping
a finite word prefix changes discrepancy by a fixed
constant and shifts its time by a fixed amount; the strict
factor $7/8$ leaves room for both changes. Local
non-$\sigma$-finiteness follows.
\end{proof}

\begin{corollary}[A one-sided null condition]
\label{cor:off-perron-one-sided}
Under the hypotheses of Theorem~\ref{thm:cubic-root-threshold},
the part of the frequency level satisfying $t(J_m-mp)\ge0$
eventually along a tail is $\mathcal H^g$-null, with no upper
bound imposed on the discrepancy.
\end{corollary}
\begin{proof}
Under the tilted law take integers $n_{j+1}=n_j^3$ with
$n_1\ge2$. Regardless of the past,
$|Z_{n_j}|\le n_j=o(\sqrt{n_{j+1}-n_j})$.
By Lemma~\ref{lem:tilted-clt}, the conditional probability
that $tZ_{n_{j+1}}<0$ is uniformly bounded away from zero
for large $j$. Iterated conditioning shows that eventual
nonnegativity has probability zero. For a fixed onset,
the probabilities of surviving with $tZ_i\ge0$ up to time
$m$ therefore tend to zero. Each eligible length-$m$
cylinder at a prescribed vertex has probability at least
a constant times $\exp(-hm)$ by \eqref{eq:cubic-tilted-weight}.
The corresponding cover at diameter $r_{N+m}$ has
$g$-cost bounded by a fixed multiple of the survival probability.
Taking the countable union over starts and onsets proves nullity.
\end{proof}


For example, in the full alphabet for $k=2$, there are two
marked and three unmarked labels. At $p=1/2$ one has
$t=\log(3/2)$, $P''(t)=1/4$ and $e^h=2\sqrt6$.
The theorem separates the eventually prescribed square-root
tolerance from, for example, the $m^{1/4}$ tolerance at
the same reference gauge $g(r_n)=(2\sqrt6)^{-n}$.
The next subsection computes the sharp amplitude at exponent
$1/3$ by evaluating the confinement rate. The elementary exponent and amplitude
comparisons are checked in \texttt{DensityCube.lean}.
The tilted probability and Hausdorff measure arguments are paper proofs.



\subsection{The exact amplitude at the cubic-root scale}

In the hypotheses of Theorem~\ref{thm:cubic-root-threshold}, put
$\beta=|t|$, $\sigma^2=P''(t)>0$ and
$V_m=-\operatorname{sgn}(t)(J_m-mp)$.
Let $K=\pi^2\sigma^2/2$. A tilted length-$m$ cylinder has probability
comparable, uniformly in the endpoints, to $\exp(-hm-\beta V_m)$.

\begin{lemma}[Confinement rate in a widening interval]
\label{lem:sharp-moving-strip}
Let $a<b$ be real and let
$E_{a,b}(n;s)$ be the event that
$a j^{1/3}\le V_j\le b j^{1/3}$ for $s\le j\le n$.
For any fixed $s$ and any $\epsilon>0$,
\[
 \Pr_t(E_{a,b}(n;s))
 \ll_{a,b,s,\epsilon}
 \exp\left[-\left(\frac{3K}{(b-a)^2}-\epsilon\right)n^{1/3}\right].
\]
There is also a path measure, from any prescribed starting
vertex, supported on paths with
\[
 a m^{1/3}-O(1)\le V_m\le b m^{1/3}+O(1)
 \qquad(m\ \text{sufficiently large}),
\]
whose probability of any length-$m$ cylinder is at most
\[
 O_\epsilon(1)\exp\left[
   \left(\frac{3K}{(b-a)^2}+\epsilon\right)m^{1/3}\right]
                  \Pr_t(\text{cylinder}).
\]
The additive constants in the support condition are uniform over
the constructed paths.
\end{lemma}
\begin{proof}
Under the tilted law, the Poisson-corrected increments used in
Lemma~\ref{lem:tilted-clt} are bounded martingale differences.
Their conditional variances $s(v)=\mathbb E_t[\xi^2\mid v]$ are
functions of the current vertex alone, with stationary mean
$\pi_t(s)=\sigma^2$ by the variance formula. Two standard
consequences of Perron--Frobenius theory for finite irreducible
chains turn this into an almost-sure strong law from every
starting vertex, with no aperiodicity assumption. First, the
tilted kernel $K_t$ has a unique stationary law $\pi_t$, and
the eigenvalue $1$ of $K_t$ is simple (in the periodic case the
remaining unit-modulus eigenvalues are the nontrivial
$\gamma$-th roots of unity), so the Poisson equation used in
Lemma~\ref{lem:tilted-clt} is solvable on the zero-mean
subspace. Second, if $K_t$ has period $\gamma\ge2$, the vertex
set splits into $\gamma$ cyclic classes, and on each class the
$\gamma$-step chain $K_t^{\gamma}$ is irreducible and
aperiodic; the strong law for these $\gamma$ skeletons, summed
over the classes, gives the Ces\`aro-average strong law
$\frac1m\sum_{i<m}s(X_i)\to\sigma^2$ almost surely from every
starting vertex. The ergodic theorem therefore gives convergence
of the rescaled
predictable quadratic variation to $\sigma^2 u$ uniformly
on bounded $u$-intervals. The rescaled jumps tend uniformly
to zero. The martingale functional central limit theorem
\cite[Theorem~2.1(ii)]{WhittFCLT} thus gives convergence of
$V_{\lfloor R^2 u\rfloor}/R$ to variance-$\sigma^2$ Brownian
motion, on every fixed time interval; the bounded Poisson
potential vanishes upon rescaling. The convergence is uniform
over the finite set of initial vertices.

For Brownian motion killed on leaving a fixed interval of
width $W$, the Dirichlet heat kernel has sine eigenfunctions
and eigenvalues $e^{-Kj^2 A/W^2}$ at time $A$, $j\ge1$.
The sine expansion bounds the maximum survival probability
by $C_W e^{-KA/W^2}$ for large $A$.
For starting points in a fixed compact interior interval
and endpoint in a fixed interior interval of positive length,
it bounds the minimum such survival probability below by
$c_W e^{-KA/W^2}$, since the first sine function is
strictly positive there and the remaining terms decay faster.

Fix a large $A$ and split times into blocks starting at $n_i$,
of lengths $\lfloor A n_i^{2/3}\rfloor$, from some sufficiently
large $n_0$. Then $n_{i+1}^{1/3}-n_i^{1/3}\to A/3$,
so the number of completed blocks by $n$ is
$(3/A+o(1))n^{1/3}$. On any one block the ratio of the
moving scale $j^{1/3}$ to $n_i^{1/3}$ tends uniformly to one.
For the upper bound, widen $[a,b]$ by $\eta$ on each side:
confinement to the moving interval implies survival in this
wider static interval after division by $n_i^{1/3}$.
The path functionals $\sup_{0\le u\le A}x(u)$ and
$\inf_{0\le u\le A}x(u)$ are continuous on $C[0,A]$, so the
confinement event for the widened closed interval is a closed
subset of $C[0,A]$. By the Portmanteau theorem applied to the
functional convergence, the conditional probability of this
block event is bounded above by twice the Brownian upper bound
for all sufficiently late blocks. This holds uniformly over
the allowed starting positions: the block probability is, by
the Markov property, a function of the finitely many starting
vertices and of the rescaled starting position alone. Were the
uniformity to fail, a subsequence of blocks with rescaled
starting positions in the fixed widened compact interval would
have probabilities bounded away from the Brownian bound;
compactness gives a further subsequence whose rescaled starts
converge, and functional convergence together with continuity
of the bounded deterministic shift (Slutsky's theorem) bounds
the probabilities along that subsequence by the same Brownian
bound, a contradiction. The closed-path event has the same
Brownian probability as
killed survival: a Brownian path that first reaches the
boundary before the final time crosses to the exterior
immediately with conditional probability one, by the
reflection principle, so the closed and open confinement
events differ only by a null set, and a hit
exactly at the final time has probability zero since the
endpoint has a continuous distribution.
Discard the final incomplete block. Multiply the upper bounds
and take first $\eta$ sufficiently small and then $A$
sufficiently large, giving the claimed upper rate.

For the lower bound, use the contracted interval
$[a+\eta,b-\eta]$ and a compact interval $J$ of positive length
strictly in its interior. Require a start in $J$, survival in
the open contracted interval, and an endpoint in a smaller
interval $J'$ strictly inside $J$.
Functional convergence gives a uniform lower probability of
at least half the corresponding Brownian lower bound.
Since $J'$ lies strictly inside $J$, rescaling the endpoint
by $(n_{i+1}/n_i)^{1/3}$ puts the next starting position in $J$.
For large $n_i$ the contracted block interval lies inside
the moving interval. A virtual initial time $n_0$ and an
additive initial shift place the first position inside $J$.
Choose each block under its tilted law conditioned on the
block event. The ratio of the new cylinder probability to
its original tilted probability is bounded above by the
product of inverse block probabilities, including at most
one extra factor for a partial block. With $\eta$ small
and then $A$ large, the block count above gives the stated
rate. Undoing the initial shift and the finite time offset
changes the support bounds by bounded additive errors.
\end{proof}

\begin{theorem}[Sharp cubic-root amplitude and endpoint localization]
\label{thm:sharp-cubic-amplitude}
Under the hypotheses of Theorem~\ref{thm:cubic-root-threshold},
put
\[
 C_*=\left(\frac{81\pi^2P''(t)}{8|t|}\right)^{1/3}.
\]
At the same reference Hausdorff gauge:
\begin{enumerate}
\item If $0<C<C_*$, even the one-sided restriction
$V_m\le C m^{1/3}$ eventually along a tail gives a null set.
\item If $C>C_*$, the symmetric restriction
$|J_m-mp|\le C m^{1/3}$ eventually gives a locally
non-$\sigma$-finite set.
\item At $C=C_*$, the portion of the restriction with
\[
 \liminf_{m\to\infty} V_m/m^{1/3}\ne C_*/3
 \quad\text{or}\quad
 \limsup_{m\to\infty} V_m/m^{1/3}\ne C_*
\]
is null. The assertion applies also to the one-sided
restriction at $C=C_*$.
\end{enumerate}
\end{theorem}
\begin{proof}
For the upper bound fix the recurrence start, initial integers,
graph vertex and the onset of the bound. If
$V_m\le d m^{1/3}$ at a time $m$, the sum of the costs of
eligible cylinders at that time is at most a constant times
$e^{\beta d m^{1/3}}$ times their tilted probability.
An eventual lower bound $V_j\ge a j^{1/3}$, together with
the upper bound $C j^{1/3}$, improves this estimate to
\begin{equation}
 O(1)\exp\left[
    \left(\beta d-\frac{3K}{(C-a)^2}+\epsilon\right)m^{1/3}
 \right]
 \label{eq:sharp-cubic-limsup-cost}
\end{equation}
by Lemma~\ref{lem:sharp-moving-strip}, for fixed $a<C$.
A negative exponent is summable in $m$ and gives nullity
for infinitely many such cylinders in the limit.

If $v=\liminf V_m/m^{1/3}<0$, infinitely many visits below
a fixed negative threshold already have summable cost,
using the probability bound $1$ alone.
If $0\le v<C$ and $\beta v(C-v)^2<3K$, choose rational
$a<v<d$ so that $\beta d<3K/(C-a)^2$; this is possible because
$\beta d(C-a)^2$ is continuous in $(a,d)$ and tends to
$\beta v(C-v)^2<3K$ as $(a,d)\to(v,v)$ within $a<v<d$.
Since $v=\liminf V_j/j^{1/3}>a$, the definition of the limit
inferior gives $V_j\ge a\,j^{1/3}$ eventually, while $d>v$
forces $V_m\le d\,m^{1/3}$ infinitely often. Infinitely many
visits below $d$ then give a null set by
\eqref{eq:sharp-cubic-limsup-cost}.
The remaining case $v=C$ (which the one-sided hypothesis does
not exclude) follows by the same reasoning: pick rationals
$a<C<d$; then $\liminf V_j/j^{1/3}=C>a$ gives
$V_j\ge a\,j^{1/3}$ eventually and $d>C$ forces infinitely many
visits below $d\,m^{1/3}$, and since $3K/(C-a)^2\to\infty$ as
$a\uparrow C$ while $\beta d\to\beta C$, the rationals can be
chosen with $\beta d<3K/(C-a)^2$, making the exponent in
\eqref{eq:sharp-cubic-limsup-cost} strictly negative and
yielding the same null conclusion.
There are only countably many choices of $a,d$ and starting data.
The optimization is
\[
 \frac{4C^3}{27}-v(C-v)^2
    =(v-C/3)^2(4C/3-v) .
\]
On $[0,C]$ the maximum is $4C^3/27$, attained only at $C/3$.
Since $\beta 4C_*^3/27=3K$, all cases below $C_*$ are null.
At $C_*$ the same proof excludes every liminf other than $C_*/3$,
up to a null set. Any smaller limsup would imply an eventual
upper bound with a rational positive amplitude below $C_*$,
giving the second endpoint restriction.

For the lower bound pick $C_*<b<C$ and $a=b/3>0$. Then
$\beta a>3K/(b-a)^2$: here $b-a=2b/3$, so the inequality reads
$\beta b^3/3>27K/4$, that is, $\beta b^3>81K/4=\beta C_*^3$,
which holds because $b>C_*$. Choose $\epsilon>0$ smaller than
this gap in Lemma~\ref{lem:sharp-moving-strip}.
Since $V_m=-\operatorname{sgn}(t)(J_m-mp)$, the symmetric
restriction $|J_m-mp|\le C m^{1/3}$ of part~(2) is exactly the
condition $|V_m|\le C m^{1/3}$. Along the supported paths the
lower support bound with $a=b/3>0$ forces
$0<V_m\le b m^{1/3}+O(1)$ eventually, so
$|V_m|=V_m\le b m^{1/3}+O(1)\le C m^{1/3}$ eventually, and the
strict margin $(C-b)\,m^{1/3}-O(1)$ stays positive for all
large $m$. This margin absorbs the bounded additive errors
produced by the initial shift and by any finite prefix change:
the resulting measure is supported in the required symmetric
bound. Its cylinder probabilities
are $o(e^{-hm})$ uniformly, because the tilted cylinder
weight contains the factor $\exp(-\beta V_m)$ and
$V_m\ge a m^{1/3}-O(1)$.
The at-most-three-prefix covering estimate from
Theorem~\ref{thm:finite-state-quant} shows that every set of
finite $\mathcal H^g$ measure has zero outer mass for this
measure. The whole supported set is therefore
non-$\sigma$-finite. The chart can be placed inside any
nonempty open arc by the dense arithmetic translations,
as in Theorem~\ref{thm:cubic-root-threshold}; the strict
amplitude margin survives these translations,
so it gives local non-$\sigma$-finiteness.
\end{proof}

The constant has the same optimization as the branching
random-walk barrier in \cite[Theorem~1.2]{JaffuelBarrier}
after rescaling the displacement by $\beta$. The conclusion
here concerns reference-gauge Hausdorff measure on deterministic
arithmetic codes of finite-state paths.
At the exact amplitude the theorem localizes the remaining
measure question to the two specified extremal limits.
\texttt{SharpWindow.lean} kernel-checks exactly three algebraic
lemmas and nothing beyond them: the factorization identity
$4c^3/27-v(c-v)^2=(v-c/3)^2(4c/3-v)$ (\texttt{cubicGap}), the
strict maximality of $v(c-v)^2$ on $[0,c]$, attained only at
$c/3$ (\texttt{strictPeak}), and the location-independent
error bound for residue-class floor counts
(\texttt{residueCountScaledError}). All probabilistic and
measure arguments in this subsection are paper proofs.


\section{General polynomial-growth denominator recurrences}\label{sec:poly-rec}

The quantization, exact gauge, three-way positive-level splitting, and
subcritical dimension phase transition extend from the Euler recurrence
\eqref{eq:rec} to arbitrary polynomial-growth three-term recurrences. Throughout
this subsection, fix a degree $r\in\Z_{\ge1}$, a leading coefficient
$c\in\Z_{\ge1}$, and an integer shift $d\in\Z$. Set $A_j:=cj^r+d$, choose an
integer $J_0\ge1$ such that $A_{J_0}=cJ_0^r+d\ge2$ (so $A_j$ is strictly
increasing and $A_j\ge2$ for all $j\ge J_0$), and let $(Q_j)_{j\ge0}$ be any
sequence of positive integers satisfying
\begin{equation}\label{eq:poly-rec}
Q_{j+1}=(cj^r+d)Q_j+Q_{j-1}\qquad\text{for all }j\ge J_0.
\end{equation}
For $x\in\R$ and $\alpha\ge0$, define the weighted approximation exponents
\[
\Lam^{(r,c,d)}(x):=\limsup_{j\to\infty}j^r\|Q_jx\|\in[0,\infty],\qquad
\Lam_\alpha^{(r,c,d)}(x):=\limsup_{j\to\infty}j^\alpha\|Q_jx\|\in[0,\infty],
\]
so that $\Lam^{(r,c,d)}=\Lam_r^{(r,c,d)}$. Fix the Borel rounding rule
$M_j(x):=\lfloor Q_jx+1/2\rfloor\in\Z$ and signed remainder
$u_j(x):=Q_jx-M_j(x)\in[-1/2,1/2)$ (so $|u_j(x)|=\|Q_jx\|$), and for $j\ge J_0$
define the integer recurrence defect (digit)
\begin{equation}\label{eq:poly-defect-id}
d_j(x):=M_{j+1}(x)-A_jM_j(x)-M_{j-1}(x)=A_ju_j(x)+u_{j-1}(x)-u_{j+1}(x)\in\Z.
\end{equation}

\begin{lemma}[Basic growth and non-asymptotic factorial bound, Layer~A in Lean]\label{lem:poly-growth}
Put $J_1:=\max(J_0+1,\,2|d|+1)$. Then:
\begin{enumerate}[label=(\roman*),leftmargin=2em]
\item For all $j\ge J_0$, $Q_{j+1}>A_jQ_j\ge2Q_j$.
\item For all $j\ge J_0+1$, $0<Q_{j-1}<\frac12Q_j<Q_j$, and hence
\[
cj^r+d<\frac{Q_{j+1}}{Q_j}<cj^r+d+1\le(c+|d|+1)j^r.
\]
\item For all $j\ge J_1$, $A_j\ge\frac{c}{2}j^r$, and for every $m\ge J_1$,
\[
Q_m\ge C_{J_1}\left(\frac{c}{2^{r+1}}\right)^m(m!)^r,\qquad
C_{J_1}:=Q_{J_1}\left(\frac{c}{2}\right)^{-J_1}\bigl((J_1-1)!\bigr)^{-r}>0.
\]
\item As $n\to\infty$, $\log Q_n=rn\log n+(\log c-r)n+O(\log n)$.
\end{enumerate}
\end{lemma}

\begin{proof}
Since $Q_{j-1}\ge1$ and $A_j\ge2$ for $j\ge J_0$, \eqref{eq:poly-rec} yields
$Q_{j+1}>A_jQ_j\ge2Q_j$, proving~(i). If $j\ge J_0+1$, then $j-1\ge J_0$, so~(i)
gives $Q_j>2Q_{j-1}$, i.e.\ $0<Q_{j-1}/Q_j<1/2<1$; adding $A_j$ gives
$cj^r+d<Q_{j+1}/Q_j<cj^r+d+1\le(c+|d|+1)j^r$ (using $j^r\ge1$), proving~(ii). For
$j\ge J_1$, we have $|d|<j/2\le\frac12 cj^r$ and therefore $A_j\ge\frac{c}{2}j^r$.
Multiplying $Q_{j+1}>A_jQ_j\ge\frac{c}{2}j^r Q_j$ over $j=J_1,\dots,m-1$ gives
\[
Q_m\ge Q_{J_1}\left(\frac{c}{2}\right)^{m-J_1}
\left(\frac{(m-1)!}{(J_1-1)!}\right)^r
=C_{J_1}\left(\frac{c}{2}\right)^m\frac{(m!)^r}{m^r}
\ge C_{J_1}\left(\frac{c}{2^{r+1}}\right)^m(m!)^r,
\]
since $m\le2^m$ for all $m\ge1$ implies $m^r\le(2^r)^m$, proving~(iii). Finally,
by~(ii), $\log(Q_{j+1}/Q_j)=r\log j+\log c+O(j^{-r})$ for $j\ge J_1$; summing
over $j=J_1,\dots,n-1$ and using $\sum_{j=1}^{n-1}\log j=n\log n-n+O(\log n)$
together with $\sum_{j=J_1}^{n-1}j^{-r}=O(\log n)$ (for $r=1$) or $O(1)$ (for
$r\ge2$) gives~(iv).
\end{proof}

When no pre-existing convergent sequence of $\ee$ is present, the role of the
group $L=\Z+\Z\frac{\ee+1}{2}$ is played by an intrinsic rank-$2$ lattice built
directly from the homogeneous recurrence $X_{j+1}=A_jX_j+X_{j-1}$ ($j\ge J_0$).
Let $(U_j)_{j\ge J_0-1}$ and $(V_j)_{j\ge J_0-1}$ be the two fundamental integer
solutions specified by
\[
U_{J_0-1}=1,\quad U_{J_0}=0;\qquad V_{J_0-1}=0,\quad V_{J_0}=1.
\]

\begin{lemma}[Intrinsic zero-level lattice and Casorati rigidity, Layer~A in Lean]\label{lem:poly-lattice}
\begin{enumerate}[label=(\roman*),leftmargin=2em]
\item For all $j\ge J_0-1$, $Q_j=Q_{J_0-1}U_j+Q_{J_0}V_j$ and
$U_jV_{j+1}-U_{j+1}V_j=(-1)^{j-J_0+1}$.
\item The limits $\alpha_U:=\lim_{j\to\infty}U_j/Q_j$ and
$\alpha_V:=\lim_{j\to\infty}V_j/Q_j$ exist, satisfy
$Q_{J_0-1}\alpha_U+Q_{J_0}\alpha_V=1$, and obey
\[
\left|\alpha_U-\frac{U_j}{Q_j}\right|<\frac{Q_{J_0}}{Q_jQ_{j+1}},\qquad
\left|\alpha_V-\frac{V_j}{Q_j}\right|<\frac{Q_{J_0-1}}{Q_jQ_{j+1}}
\qquad\text{for all }j\ge J_0.
\]
\item Set $g:=\gcd(Q_{J_0-1},Q_{J_0})\ge1$, write $Q_{J_0-1}=gq_0$ and
$Q_{J_0}=gq_1$ with $\gcd(q_0,q_1)=1$, pick integers $p_0,p_1\in\Z$ with
$p_0q_1-p_1q_0=1$, and put $\tilde Q_j:=Q_j/g=q_0U_j+q_1V_j$ and
$\tilde P_j:=p_0U_j+p_1V_j$. Then
\[
\theta:=\lim_{j\to\infty}\frac{\tilde P_j}{\tilde Q_j}=g(p_0\alpha_U+p_1\alpha_V)\in\R\setminus\Q,
\]
and the intrinsic rank-$2$ lattice
\[
L^{(r,c,d)}:=\Z\alpha_U+\Z\alpha_V=\frac{1}{g}\bigl(\Z+\Z\theta\bigr)
\]
is independent of the choice of $(p_0,p_1)$, contains $\Z$, and projects to a
dense countable subgroup $\bar L^{(r,c,d)}:=L^{(r,c,d)}/\Z$ of $\R/\Z$.
\item For $x\in\R$, $\Lam^{(r,c,d)}(x)=0$ if and only if $u_j(x)\to0$ and
eventually $d_j(x)=0$, if and only if $x\in L^{(r,c,d)}$. Moreover, every
$x\in L^{(r,c,d)}$ satisfies $\|Q_jx\|=O(1/Q_{j+1})$, whence
$\Lam_\alpha^{(r,c,d)}(x)=0$ for all $\alpha\ge0$.
\end{enumerate}
\end{lemma}

\begin{proof}
Part~(i) holds at $j=J_0-1$ by inspection and propagates to all $j\ge J_0-1$
because $(U_j),(V_j),(Q_j)$ satisfy the same recurrence \eqref{eq:poly-rec} whose
transfer matrices $\bigl(\begin{smallmatrix}A_j&1\\1&0\end{smallmatrix}\bigr)$
have determinant $-1$. Expanding $Q_j=Q_{J_0-1}U_j+Q_{J_0}V_j$ in the adjacent
ratio differences yields, for $j\ge J_0$,
\[
\frac{U_{j+1}}{Q_{j+1}}-\frac{U_j}{Q_j}
=\frac{(-1)^{j-J_0}Q_{J_0}}{Q_jQ_{j+1}},\qquad
\frac{V_{j+1}}{Q_{j+1}}-\frac{V_j}{Q_j}
=\frac{(-1)^{j-J_0+1}Q_{J_0-1}}{Q_jQ_{j+1}}.
\]
By Lemma~\ref{lem:poly-growth}(i), $Q_{j+2}>2Q_{j+1}>Q_j$ for all $j\ge J_0$,
so $1/(Q_jQ_{j+1})$ decreases strictly to $0$; the alternating series test
gives convergence and bounds the tail strictly by the magnitude of the first
term, proving~(ii) (and $Q_{J_0-1}\alpha_U+Q_{J_0}\alpha_V=1$ follows by taking
$j\to\infty$ in $Q_{J_0-1}(U_j/Q_j)+Q_{J_0}(V_j/Q_j)=1$).

For~(iii), $\bigl(\begin{smallmatrix}p_0&p_1\\q_0&q_1\end{smallmatrix}\bigr)\in\mathrm{SL}_2(\Z)$
gives $\tilde P_j\tilde Q_{j+1}-\tilde P_{j+1}\tilde Q_j=(-1)^{j-J_0+1}$, so
$\theta=\lim_{j\to\infty}\tilde P_j/\tilde Q_j=\lim_{j\to\infty}g(p_0U_j+p_1V_j)/Q_j=g(p_0\alpha_U+p_1\alpha_V)$
satisfies $0<|\theta-\tilde P_j/\tilde Q_j|<(\tilde Q_j\tilde Q_{j+1})^{-1}$ for
all $j\ge J_0$. If $\theta=p/q\in\Q$ ($q\ge1$), then
$0<|q\tilde P_j-p\tilde Q_j|<q/\tilde Q_{j+1}<1$ for large $j$, contradicting
$|q\tilde P_j-p\tilde Q_j|\in\Z_{\ge1}$; hence $\theta\notin\Q$. Since
$\bigl(\begin{smallmatrix}p_0&p_1\\q_0&q_1\end{smallmatrix}\bigr)\in\mathrm{SL}_2(\Z)$
is unimodular and $q_0\alpha_U+q_1\alpha_V=1/g$, we have
$\Z\alpha_U+\Z\alpha_V=\Z(q_0\alpha_U+q_1\alpha_V)+\Z(p_0\alpha_U+p_1\alpha_V)
=\Z\frac{1}{g}+\Z\frac{\theta}{g}=\frac{1}{g}(\Z+\Z\theta)$, which is thus
independent of $(p_0,p_1)$ and contains $1=g\cdot\frac{1}{g}$. Because
$\theta/g\notin\Q$, Kronecker's theorem shows that $\bar L^{(r,c,d)}$ is dense
in $\R/\Z$.

For~(iv), if $\Lam^{(r,c,d)}(x)=0$, then $j^r|u_j(x)|\to0$, so
$|d_j(x)|\le cj^r|u_j(x)|+|d||u_j(x)|+|u_{j-1}(x)|+|u_{j+1}(x)|\to0$; since
$d_j(x)\in\Z$, eventually $d_j(x)=0$. Conversely, if $u_j(x)\to0$ and $d_j(x)=0$
for all $j\ge N\ge J_0+1$, then by~(i)
$\det\bigl(\begin{smallmatrix}U_{N-1}&V_{N-1}\\U_N&V_N\end{smallmatrix}\bigr)=\pm1$,
so $(M_{N-1},M_N)=a(U_{N-1},U_N)+b(V_{N-1},V_N)$ for unique $a,b\in\Z$, and by
induction $M_j(x)=aU_j+bV_j$ for all $j\ge N-1$; dividing by $Q_j$ and using
$u_j(x)\to0$ yields $x=a\alpha_U+b\alpha_V\in L^{(r,c,d)}$. Finally, if
$x=a\alpha_U+b\alpha_V\in L^{(r,c,d)}$, then~(ii) gives
$\|Q_jx\|\le|Q_jx-(aU_j+bV_j)|\le(|a|Q_{J_0}+|b|Q_{J_0-1})/Q_{j+1}$ for
$j\ge J_0$, and Lemma~\ref{lem:poly-growth}(iii) yields $j^\alpha/Q_{j+1}\to0$
for every $\alpha\ge0$.
\end{proof}

\begin{theorem}[General $(cj^r+d)$ recurrences: complete spectrum discretization,
Layer~A in Lean]\label{thm:general-rec}
For any $r,c\in\Z_{\ge1}$, $d\in\Z$, and $(Q_j)_{j\ge0}$ satisfying
\eqref{eq:poly-rec}:
\begin{enumerate}[label=(\roman*),leftmargin=2em]
\item Whenever $u_j(x)\to0$ (in particular, whenever $\Lam^{(r,c,d)}(x)<\infty$),
the digits $d_j(x)$ satisfy
\[
\lim_{j\to\infty}\bigl(d_j(x)-cj^ru_j(x)\bigr)=0,\qquad
\lim_{j\to\infty}\bigl(|d_j(x)|-cj^r\|Q_jx\|\bigr)=0.
\]
\item The full spectrum of $\Lam^{(r,c,d)}$ is completely discrete with step $1/c$:
\[
\Lam^{(r,c,d)}(\R)=\frac{1}{c}\Z_{\ge0}\cup\{\infty\}
=\left\{0,\frac{1}{c},\frac{2}{c},\frac{3}{c},\dots\right\}\cup\{\infty\},
\]
so every open interval $\bigl(\frac{k}{c},\frac{k+1}{c}\bigr)$ ($k\in\Z_{\ge0}$)
is an empty spectral gap, and for every $t\in\bigl[\frac{k}{c},\frac{k+1}{c}\bigr)$
the sublevel set $E_t^{(r,c,d)}:=\{x\in\R:\Lam^{(r,c,d)}(x)\le t\}$ equals
$E_{k/c}^{(r,c,d)}=\{x\in\R:u_j(x)\to0\text{ and }\limsup_{j\to\infty}|d_j(x)|\le k\}$
(with $E_t^{(r,c,d)}=L^{(r,c,d)}$ for $t<1/c$).
\end{enumerate}
(Lean: for all $r\ge1$, proved at Layer~A in
\texttt{EpiFormal/\allowbreak TrichotomyAndPolyExtension.lean} via
\texttt{PolyRecData.\allowbreak poly\_\allowbreak signed\_\allowbreak defect\_\allowbreak approx},
\texttt{poly\_\allowbreak abs\_\allowbreak defect\_\allowbreak approx},
\texttt{poly\_\allowbreak spectral\_\allowbreak gap\_\allowbreak open\_\allowbreak interval},
\texttt{poly\_\allowbreak levelSet\_\allowbreak eq\_\allowbreak on\_\allowbreak gap},
\texttt{poly\_\allowbreak levelSet\_\allowbreak below\_\allowbreak eq\_\allowbreak lattice}, and
\texttt{poly\_\allowbreak range\_\allowbreak apprExponent}; the $r=1$
specialization also appears in
\texttt{EpiFormal/\allowbreak Full\allowbreak Spectrum\allowbreak Discrete.lean}.)
\end{theorem}

\begin{proof}
Rearranging \eqref{eq:poly-defect-id} gives
$d_j(x)-cj^ru_j(x)=du_j(x)+u_{j-1}(x)-u_{j+1}(x)\to0$ whenever $u_j(x)\to0$
(which holds whenever $\Lam^{(r,c,d)}(x)<\infty$), and the reverse triangle
inequality $\bigl||d_j(x)|-cj^r|u_j(x)|\bigr|\le|d_j(x)-cj^ru_j(x)|$ proves~(i).
Since $|d_j(x)|\in\Z_{\ge0}$, whenever $\Lam^{(r,c,d)}(x)<\infty$ we obtain
$c\Lam^{(r,c,d)}(x)=\limsup_{j\to\infty}|d_j(x)|\in\Z_{\ge0}$, excluding all
open intervals $(k/c,(k+1)/c)$ and forcing $E_t^{(r,c,d)}=E_{k/c}^{(r,c,d)}$ on
$[k/c,(k+1)/c)$. Conversely, $0\in L^{(r,c,d)}$ realizes $0$,
Theorem~\ref{thm:poly-phase-transition}(i) below realizes $\infty$ on a set of
full Lebesgue measure, and for each $k\ge1$ the constant-digit construction
$d_j\equiv k$ in Theorem~\ref{thm:poly-gauge-trichotomy} below produces a point
$x_k^+$ with $|u_j(x_k^+)|<4k/A_j=O(j^{-r})\to0$ and $d_j(x_k^+)=k$ for large
$j$, giving $\Lam^{(r,c,d)}(x_k^+)=k/c$ by~(i). The $t<1/c$ identity follows from
Lemma~\ref{lem:poly-lattice}(iv).
\end{proof}

For each integer $k\ge1$, define the continuous nondecreasing critical gauge
$h_k^{(r,c,d)}:[0,\infty)\to[0,\infty)$ by $h_k^{(r,c,d)}(0)=0$,
$h_k^{(r,c,d)}(1/Q_n)=(2k+1)^{-n}$ for all $n\ge J_0$, linear interpolation on
each $[1/Q_{n+1},1/Q_n]$ ($n\ge J_0$), and constant value $(2k+1)^{-J_0}$ for
$t\ge1/Q_{J_0}$ (well-defined because $Q_n$ is strictly increasing for
$n\ge J_0$ by Lemma~\ref{lem:poly-growth}(i)). Similarly define the universal
binary gauge $g_2^{(r,c,d)}(1/Q_n):=2^{-n}$ ($n\ge J_0$).

\begin{theorem}[Critical gauges and positive-level trichotomy for general $(cj^r+d)$, part~(iii) Layer~A in Lean]\label{thm:poly-gauge-trichotomy}
Fix an integer $k\ge1$. At the positive spectral level $k/c$:
\begin{enumerate}[label=(\roman*),leftmargin=2em]
\item \textbf{Limsup level (alphabet size $2k+1$, gauge base $2k+1$):}
$x\in E_{k/c}^{(r,c,d)}$ if and only if $u_j(x)\to0$ and eventually
$d_j(x)\in\{-k,\dots,k\}$. On $\R/\Z$, the exact level
$S_k^{(r,c,d)}:=\{x\in\R/\Z:\Lam^{(r,c,d)}(x)=k/c\}$ has infinite but
$\sigma$-finite $\mathcal H^{h_k^{(r,c,d)}}$ measure (with infinite measure in
every nonempty open arc), whereas $\{\Lam^{(r,c,d)}<k/c\}$ is
$\mathcal H^{h_k^{(r,c,d)}}$-null.
\item \textbf{Unsigned exact limit set (universal binary alphabet $\{\pm k\}$, gauge base $2$):}
$E_{\mathrm{lim}}^{(r,c,d)}(k/c):=\{x\in\R:\lim_{j\to\infty}j^r\|Q_jx\|=k/c\}$
consists precisely of those $x$ with $u_j(x)\to0$ and eventually
$d_j(x)\in\{-k,+k\}$ (two branches per step, independently of $k\ge1$). On
$\R/\Z$, $E_{\mathrm{lim}}^{(r,c,d)}(k/c)$ has infinite, $\sigma$-finite
$\mathcal H^{g_2^{(r,c,d)}}$ measure (infinite on every nonempty open arc),
while $\mathcal H^{h_k^{(r,c,d)}}(E_{\mathrm{lim}}^{(r,c,d)}(k/c))=0$. For
$t\notin\frac1c\Z_{\ge0}$, $E_{\mathrm{lim}}^{(r,c,d)}(t)=\emptyset$.
\item \textbf{Signed exact limit sets (single tail branch, coset collapse):}
$E_{\mathrm{lim}}^{(r,c,d),\pm}(k/c):=\{x\in\R:\lim_{j\to\infty}j^ru_j(x)=\pm k/c\}$
consists precisely of those $x$ with $u_j(x)\to0$ and eventually $d_j(x)=\pm k$.
With $x_k^+$ built from $d_j\equiv+k$ ($j\ge N$) and $x_k^-:=-x_k^+$, each signed
exact limit set collapses to a single countable coset of the intrinsic lattice
$L^{(r,c,d)}$ of Lemma~\ref{lem:poly-lattice}:
\[
E_{\mathrm{lim}}^{(r,c,d),+}(k/c)=x_k^+ + L^{(r,c,d)},\qquad
E_{\mathrm{lim}}^{(r,c,d),-}(k/c)=x_k^- + L^{(r,c,d)},
\]
and is therefore null for every non-atomic Hausdorff gauge.
\end{enumerate}
\end{theorem}

\begin{proof}
The digit characterizations follow directly from Theorem~\ref{thm:general-rec}.
For the upper bounds and $\sigma$-finiteness on $[0,1]$ (hence on $\R/\Z$ since
the quotient map is $1$-Lipschitz), partition $E_{k/c}^{(r,c,d)}\cap[0,1]$ into
countably many Borel classes $\mathcal C(N,m_{N-1},m_N)$ ($N\ge J_0+1$,
$m_{N-1},m_N\in\Z$) requiring $M_{N-1}=m_{N-1}$, $M_N=m_N$,
$|u_j|\le(k+1)/(cj^r)$, and $d_j\in\Sigma$ for all $j\ge N$, where
$\Sigma=\{-k,\dots,k\}$ ($|\Sigma|=2k+1$) for $E_{k/c}^{(r,c,d)}$,
$\Sigma=\{-(k-1),\dots,k-1\}$ ($|\Sigma|=2k-1$) for $\{\Lam^{(r,c,d)}<k/c\}$,
and $\Sigma=\{-k,+k\}$ ($|\Sigma|=2$) for $E_{\mathrm{lim}}^{(r,c,d)}(k/c)$.
At level $n>N$, the recurrence $M_{j+1}=A_jM_j+M_{j-1}+d_j$ yields at most
$|\Sigma|^{n-N}$ values of $M_n$, each confining $x$ to an interval of diameter
at most $1/Q_n$ (using $|u_n|\le1/2$). Since $n>J_0$, covering with these
intervals gives $\mathcal H_\delta^{h_k^{(r,c,d)}}$ cost at most
$(2k+1)^{-N}\bigl(|\Sigma|/(2k+1)\bigr)^{n-N}$ (bounded by $(2k+1)^{-N}<\infty$
when $|\Sigma|=2k+1$, and $\to0$ when $|\Sigma|<2k+1$), and when $\Sigma=\{-k,+k\}$
gives $\mathcal H_\delta^{g_2^{(r,c,d)}}$ cost at most $2^{n-N}2^{-n}=2^{-N}<\infty$.

For the lower bounds, choose an integer $N\ge J_0+1$ such that
$A_N=cN^r+d>32k$. Since $A_j\ge A_N>32k\ge32>6$ and $Q_N\ge1$, Lemma~\ref{lem:poly-growth}(i)
gives $Q_{N+1}>32kQ_N\ge32k$ and $Q_{j+1}>6Q_j$ for all $j\ge N$. Set
$M_{N-1}=M_N=0$ and, for $\Sigma\in\{\{-k,\dots,k\},\{-k,+k\}\}$ and
$\omega=(d_N,d_{N+1},\dots)\in\Sigma^{\{j\ge N\}}$, define $M_{j+1}=A_jM_j+M_{j-1}+d_j$
($j\ge N$). Then $W_j:=M_{j+1}Q_j-M_jQ_{j+1}$ satisfies $W_{N-1}=0$ and
$W_j=d_jQ_j-W_{j-1}$ ($j\ge N$). Since $j\ge N\ge J_0+1$ implies $Q_j>2Q_{j-1}$,
induction gives $|W_j|\le2kQ_j$ for all $j\ge N-1$, whence
$|M_{j+1}/Q_{j+1}-M_j/Q_j|\le2k/Q_{j+1}$. Using $Q_{m+1+s}\ge2^sQ_{m+1}$,
$x(\omega):=\lim_{j\to\infty}M_j/Q_j$ satisfies $|x(\omega)-M_m/Q_m|\le4k/Q_{m+1}$
for all $m\ge N$ (in particular $|x(\omega)|\le4k/Q_{N+1}$ since $M_N=0$).
Multiplying by $Q_m$ gives $|Q_mx(\omega)-M_m|\le4kQ_m/Q_{m+1}\le4k/A_m\le4k/A_N<1/8<1/2$
for all $m\ge N$, and at $m=N-1$ we also have $Q_{N-1}|x(\omega)|\le4kQ_{N-1}/Q_{N+1}<1/16<1/2$
(since $Q_{N+1}>32kQ_N>64kQ_{N-1}$). Hence $M_m(x(\omega))=M_m$ for all $m\ge N-1$,
no half-integer ties occur, and $d_j(x(\omega))=d_j$ for all $j\ge N$. Moreover,
$K_\Sigma:=x(\Sigma^{\{j\ge N\}})\subseteq[-4k/Q_{N+1},4k/Q_{N+1}]$, an interval
of length $<1/4<1/2$ embedding isometrically into $\R/\Z$. Pushing forward the
uniform Bernoulli product measure on $\Sigma^{\{j\ge N\}}$ yields $\mu$ on $K_\Sigma$.
For any interval $U$ of length $1/Q_{n+1}\le h<1/Q_n$ ($n>N$) and any $x,y\in K_\Sigma\cap U$,
we have $|M_j(x)-M_j(y)|<\frac14+\frac{Q_j}{Q_n}\le\frac14+\frac16<1$ for $N\le j<n$
(so $M_j(x)=M_j(y)$) and $|M_n(x)-M_n(y)|<\frac14+Q_nh<\frac54<2$ (at most $3$ values).
Thus $K_\Sigma\cap U$ meets at most $3$ cylinders of depth $n-N$, so
$\mu(U)\le3|\Sigma|^{-(n-N)}=3|\Sigma|^{N+1}|\Sigma|^{-(n+1)}\le3|\Sigma|^{N+1}g_{|\Sigma|}(h)$
(as $n+1>J_0$), where $g_{2k+1}:=h_k^{(r,c,d)}$ and $g_2:=g_2^{(r,c,d)}$. The
mass distribution principle yields $\mathcal H^{g_{|\Sigma|}}(K_\Sigma)\ge(3|\Sigma|^{N+1})^{-1}>0$.
Taking $\Sigma=\{-k,\dots,k\}$ and subtracting the $\mathcal H^{h_k^{(r,c,d)}}$-null
set $\{\Lam^{(r,c,d)}<k/c\}$ gives $\mathcal H^{h_k^{(r,c,d)}}(S_k^{(r,c,d)})>0$;
taking $\Sigma=\{-k,+k\}$ gives $K_\Sigma\subseteq E_{\mathrm{lim}}^{(r,c,d)}(k/c)$
and $\mathcal H^{g_2^{(r,c,d)}}(E_{\mathrm{lim}}^{(r,c,d)}(k/c))>0$.

By Lemma~\ref{lem:poly-lattice}(iv), $j^r\|Q_jy\|\to0$ for every $y\in L^{(r,c,d)}$,
so both $S_k^{(r,c,d)}$ and $E_{\mathrm{lim}}^{(r,c,d)}(k/c)$ are Borel,
Lebesgue-null (by Theorem~\ref{thm:poly-phase-transition}(i), or directly since
$\operatorname{Leb}_{\R/\Z}(\{\|Q_jx\|\le m/j^r\})=2m/j^r\to0$), and invariant
under translation by the dense subgroup $\bar L^{(r,c,d)}\subset\R/\Z$ of
Lemma~\ref{lem:poly-lattice}(iii). The exact Haar-uniqueness and compactness
covering arguments of Theorem~\ref{thm:gauge} and Corollary~\ref{cor:exact-gauge}(i)
then force both $\mathcal H^{h_k^{(r,c,d)}}(S_k^{(r,c,d)}\cap I)=\infty$ and
$\mathcal H^{g_2^{(r,c,d)}}(E_{\mathrm{lim}}^{(r,c,d)}(k/c)\cap I)=\infty$ on
every nonempty open arc $I\subseteq\R/\Z$.

Finally, for~(iii), take $x_k^+$ from the constant sequence $d_j\equiv+k$ ($j\ge N$).
Since $|u_j(x_k^+)|<1/8<1/2$ for all $j\ge N-1$, $x_k^-=-x_k^+$ satisfies
$M_j(x_k^-)=-M_j(x_k^+)$ and $u_j(x_k^-)=-u_j(x_k^+)$ for $j\ge N-1$, so
$d_j(x_k^-)=-k$ for $j\ge N$ and $x_k^\pm\in E_{\mathrm{lim}}^{(r,c,d),\pm}(k/c)$.
If $x,y\in E_{\mathrm{lim}}^{(r,c,d),\pm}(k/c)$, then
$j^r\|Q_j(x-y)\|\le j^r|u_j(x)-u_j(y)|\to0$, so $x-y\in L^{(r,c,d)}$ by
Lemma~\ref{lem:poly-lattice}(iv). Conversely, if $z\in L^{(r,c,d)}$, then
$u_j(z)=o(j^{-r})\to0$ while $u_j(x_k^\pm)=O(j^{-r})\to0$, so eventually
$|u_j(x_k^\pm)+u_j(z)|<1/2$, forcing $u_j(x_k^\pm+z)=u_j(x_k^\pm)+u_j(z)$ and
$j^ru_j(x_k^\pm+z)\to\pm k/c$. Thus $E_{\mathrm{lim}}^{(r,c,d),\pm}(k/c)=x_k^\pm+L^{(r,c,d)}$.
\end{proof}

\begin{corollary}[Exact critical gauge for $(cj^r+d)$]\label{cor:poly-exact-gauge}
Let $k\ge1$ and let $g:[0,\infty)\to[0,\infty)$ be any nondecreasing gauge with $g(0)=0$.
\begin{enumerate}[label=(\roman*),leftmargin=2em]
\item If $g(1/Q_n)/h_k^{(r,c,d)}(1/Q_n)\to0$ as $n\to\infty$, then
$\mathcal H^g(\{\Lam^{(r,c,d)}\le k/c\})=0$.
\item If $h_k^{(r,c,d)}(t)/g(t)\to0$ as $t\to0^+$, then $S_k^{(r,c,d)}$ is not
$\sigma$-finite for $\mathcal H^g$ (in particular $\mathcal H^g(S_k^{(r,c,d)})=\infty$).
\end{enumerate}
The exact same dichotomy holds for $E_{\mathrm{lim}}^{(r,c,d)}(k/c)$ with
$g_2^{(r,c,d)}$ in place of $h_k^{(r,c,d)}$.
\end{corollary}

\begin{proof}
This follows by the identical covering comparison and $\sigma$-finiteness
contradiction of Corollary~\ref{cor:exact-gauge}(ii)--(iii), using the upper
covers from Theorem~\ref{thm:poly-gauge-trichotomy}.
\end{proof}

\begin{remark}[Interpolation-free logarithmic scale and the $1/r$ factor]\label{rem:poly-gauge-scale}
For $1/Q_{n+1}\le t\le 1/Q_n$ ($n\ge J_0$), $-\log h_k^{(r,c,d)}(t)$ lies between
$n\log(2k+1)$ and $(n+1)\log(2k+1)$. By Lemma~\ref{lem:poly-growth}(iv),
$\log(1/t)=rn\log n+O(n)$, so $\log\log(1/t)=\log n+O(\log\log n)\sim\log n$,
yielding $\frac{\log(1/t)}{\log\log(1/t)}\sim rn$ and therefore
\[
-\log h_k^{(r,c,d)}(t)\sim\frac{\log(2k+1)}{r}\cdot\frac{\log(1/t)}{\log\log(1/t)},\qquad
-\log g_2^{(r,c,d)}(t)\sim\frac{\log2}{r}\cdot\frac{\log(1/t)}{\log\log(1/t)}
\qquad(t\downarrow0).
\]
Consequently, any gauge $g$ whose logarithmic rate against
$\frac{\log(1/t)}{\log\log(1/t)}$ exceeds $\frac{\log(2k+1)}{r}$ gives
$\mathcal H^g(\{\Lam^{(r,c,d)}\le k/c\})=0$, while any strictly smaller rate
yields non-$\sigma$-finite $\mathcal H^g$ measure on $S_k^{(r,c,d)}$.
\end{remark}

\begin{theorem}[Critical dimension zero, subcritical phase transition $1-\alpha/r$,
and supercritical collapse, Layer~A in Lean]\label{thm:poly-phase-transition}
For any $r,c\in\Z_{\ge1}$, $d\in\Z$, and $(Q_j)_{j\ge0}$ satisfying \eqref{eq:poly-rec}:
\begin{enumerate}[label=(\roman*),leftmargin=2em]
\item \textbf{Critical weight ($\alpha=r$):} For every finite $t\in\R$,
$\dimH\{x\in\R:\Lam^{(r,c,d)}(x)\le t\}=0$. Thus
$F^{(r,c,d)}:=\{x\in\R:\Lam^{(r,c,d)}(x)<\infty\}$ has Hausdorff dimension $0$
and Lebesgue measure $0$, while its complement
$F_\infty^{(r,c,d)}:=\{x\in\R:\Lam^{(r,c,d)}(x)=\infty\}$ has full Lebesgue
measure and $\dimH(F_\infty^{(r,c,d)}\cap[a,b])=1$ for every $a<b$.
\item \textbf{Subcritical weights ($0\le\alpha<r$):} For every $0\le\alpha<r$ and
every $t\ge0$ (with $0\le t\le1/2$ when $\alpha=0$),
\[
\dimH\bigl\{x\in[0,1]:j^\alpha\|Q_jx\|\longrightarrow t\bigr\}
=\dimH\bigl\{x\in\R:\Lam_\alpha^{(r,c,d)}(x)\le t\bigr\}
=1-\frac{\alpha}{r}.
\]
(For $\alpha=0$ and $t>1/2$, the prescribed-limit set is empty while
$\{\Lam_0^{(r,c,d)}\le t\}=\R$ has dimension $1=1-0/r$.)
\item \textbf{Supercritical weights ($\alpha>r$):} For every finite $t\ge0$,
$\{x\in\R:\Lam_\alpha^{(r,c,d)}(x)\le t\}=L^{(r,c,d)}$ (countable, Hausdorff
dimension $0$), while for $t<0$ the sublevel set is empty.
\end{enumerate}
\end{theorem}

\begin{proof}
\emph{Step 1: Upper bound $\dimH\{x\in\R:\Lam_\alpha^{(r,c,d)}(x)\le t\}\le1-\alpha/r$
for all $0\le\alpha\le r$.} Fix $0\le\alpha\le r$ and $t\in\R$, and put
$t_0:=\max(t,0)+1>0$. By $\Z$-periodicity and countable stability, it suffices to
show $\dimH C_{J,\alpha}\le1-\alpha/r$ for each $J\ge J_1$ (from
Lemma~\ref{lem:poly-growth}), where
$C_{J,\alpha}:=[0,1]\cap\bigcap_{j=J}^\infty\{x:\|Q_jx\|\le t_0j^{-\alpha}\}$.
Put $r_{j,\alpha}:=t_0/(j^\alpha Q_j)$ and
$I_{j,\alpha}(k):=[k/Q_j-r_{j,\alpha},k/Q_j+r_{j,\alpha}]$. Any $x\in C_{J,\alpha}$
satisfies $x\in I_{j,\alpha}(k_j(x))$ for $k_j(x)=\mathrm{round}(Q_jx)\in\{0,\dots,Q_j\}$,
with $|k_j(x)/Q_j-k_{j+1}(x)/Q_{j+1}|\le r_{j,\alpha}+r_{j+1,\alpha}$. Multiplying
by $2Q_{j+1}$ and using Lemma~\ref{lem:poly-growth}(ii) together with
$1/(j+1)^\alpha\le1\le j^{r-\alpha}$ ($0\le\alpha\le r$, $j\ge1$), each parent
$k$ yields child centers $k'$ in an interval of length
\[
2(r_{j,\alpha}+r_{j+1,\alpha})Q_{j+1}
=2t_0\left(\frac{Q_{j+1}}{j^\alpha Q_j}+\frac{1}{(j+1)^\alpha}\right)
\le 2t_0(c+|d|+2)j^{r-\alpha},
\]
so each parent has at most $Kj^{r-\alpha}$ children, where $K:=2t_0(c+|d|+2)+1$.
Hence the surviving center set $S_{n,\alpha}\subseteq\{0,\dots,Q_{J+n}\}$ at
level $m:=J+n\ge J_1$ satisfies
$|S_{n,\alpha}|\le(Q_J+1)K^n(m!)^{r-\alpha}$, while
Lemma~\ref{lem:poly-growth}(iii) and $m^\alpha\ge1$ give
$\operatorname{diam}I_{m,\alpha}(k)=2t_0/(m^\alpha Q_m)\le2t_0/(C_{J_1}(c/2^{r+1})^m(m!)^r)\to0$.
For any $s>1-\alpha/r$, set $\beta:=rs-(r-\alpha)=r(s-(1-\alpha/r))>0$ and
$K_2:=K(2^{r+1}/c)^s>0$. Using $m!\ge(m/\ee)^m$, the $s$-dimensional covering
cost satisfies
\begin{align*}
\sum_{k\in S_{n,\alpha}}\bigl(\operatorname{diam}I_{m,\alpha}(k)\bigr)^s
&\le(Q_J+1)K^{-J}\left(\frac{2t_0}{C_{J_1}}\right)^s K_2^m(m!)^{-\beta}\\
&\le(Q_J+1)K^{-J}\left(\frac{2t_0}{C_{J_1}}\right)^s
\left(\frac{\ee K_2^{1/\beta}}{m}\right)^{\beta m}\longrightarrow0
\end{align*}
as $m\to\infty$. Thus $\mathcal H^s(C_{J,\alpha})=0$ for every $s>1-\alpha/r$,
giving $\dimH\{x\in\R:\Lam_\alpha^{(r,c,d)}(x)\le t\}\le1-\alpha/r$. Taking
$\alpha=r$ proves~(i) (with $\operatorname{Leb}(F^{(r,c,d)})=0$ and
$\dimH(F_\infty^{(r,c,d)}\cap[a,b])=1$ exactly as in Corollary~\ref{cor:infinite-level}).

\emph{Step 2: Lower bound for the prescribed-limit set.}
Fix $0\le\alpha<r$ and an admissible $t$. Given $0<s<1-\alpha/r$, put
\[
  \varepsilon=\frac{r-\alpha-rs}{3},\qquad
  \beta=\alpha+\varepsilon,\qquad \gamma=r-\varepsilon .
\]
Then $0<\beta<\gamma<r$ and $\gamma-\beta-rs=\varepsilon>0$.
For all sufficiently large $j$, Lemma~\ref{lem:poly-growth}(iii) gives
$Q_{j+1}/Q_j\ge 4j^\gamma$. Set
\[
  w_j=j^{-\beta},\qquad R_j=\frac{w_j}{Q_j},\qquad
  \theta_j=\max\{w_j,\min\{tj^{-\alpha},1/2-w_j\}\}.
\]
Take $N$ large enough that, for $j\ge N$, $w_j\le1/8$ and
$tj^{-\alpha}\in[0,1/2]$. This includes $\alpha=0$, $t=0$ or $1/2$.
The interval $I_j(k)$ of radius $R_j$ centered at $(k+\theta_j)/Q_j$
satisfies $|j^\alpha\|Q_jx\|-t|\le2j^{\alpha-\beta}$ for $x\in I_j(k)$.
Indeed, $w_j\le\theta_j\le1/2-w_j$ and
$|\theta_j-tj^{-\alpha}|\le w_j$. Moreover, $I_N(0)\subset[0,1]$.

For each parent $I_j(k)$ the real interval of admissible child indices
has length $2Q_{j+1}(R_j-R_{j+1})$. Since $w_{j+1}\le w_j$ and
$Q_{j+1}/Q_j\ge 4j^\gamma\ge4$, this length is at least
$6j^{\gamma-\beta}$. A closed real interval of length at least $b$
contains at least $b$ integers whenever $b$ is a positive integer,
by starting with the ceiling of its left endpoint. Retain exactly
$b_j=\lfloor 6j^{\gamma-\beta}\rfloor$ children per parent, starting from
$I_N(0)$. At level $n\ge N$ the intervals are disjoint, and their number
$M_n=\prod_{j=N}^{n-1} b_j$ satisfies
\[
 M_n\ge \left(\frac{(n-1)!}{(N-1)!}\right)^{\gamma-\beta},
\]
because $j^{\gamma-\beta}\ge1$. Lemma~\ref{lem:poly-growth}(iv) bounds
$Q_{n+1}\le C B^n ((n-1)!)^r$ with suitable $C,B>0$.
The positive exponent $\gamma-\beta-rs$ therefore implies
$2Q_{n+1}^s/M_n\to0$ by factorial domination of exponentials.

The compact intersection $K$ is contained in the prescribed-limit set.
Assigning each child an equal fraction of its parent's mass gives a
probability measure $\mu$ supported on $K$ with mass $1/M_n$ per level-$n$
interval. For diameters $Q_{n+1}^{-1}<h\le Q_n^{-1}$,
the centers of level-$n$ intervals meeting a set of diameter $h$ lie
in a span at most $h+2R_n<2Q_n^{-1}$. At most two centers occur because
they lie on a grid of spacing $Q_n^{-1}$. For sufficiently small $h$,
the set has outer $\mu$-measure at most $2/M_n\le Q_{n+1}^{-s}\le h^s$.
Singletons have zero mass since $M_n\to\infty$. Thus
$\mathcal H^s(K)\ge1$. Letting $s$ increase to $1-\alpha/r$ gives the
desired lower bound. The prescribed-limit set is contained in the
corresponding sublevel, so Step~1 yields both equalities in (ii).
When $\alpha=0$ and $t>1/2$, the distance to $\Z$ is at most $1/2$ at
every index, which gives the two endpoint claims.

\emph{Step 3: Supercritical collapse ($\alpha>r$).} If $t<0$, the sublevel set
is empty because $\Lam_\alpha^{(r,c,d)}\ge0$. Fix $t\ge0$. If
$\Lam_\alpha^{(r,c,d)}(x)\le t$, then $\|Q_jx\|=O(j^{-\alpha})=o(j^{-r})$, so
$\Lam^{(r,c,d)}(x)=0$ and $x\in L^{(r,c,d)}$ by Lemma~\ref{lem:poly-lattice}(iv);
conversely, every $x\in L^{(r,c,d)}$ satisfies $\Lam_\alpha^{(r,c,d)}(x)=0\le t$
by Lemma~\ref{lem:poly-lattice}(iv), proving~(iii).
\end{proof}

\paragraph{Lean status for Section~\ref{sec:poly-rec}.}
In \texttt{EpiFormal/\allowbreak TrichotomyAndPolyExtension.lean}, thirty-seven
audited declarations under \texttt{PolyRecData} (base axioms only:
\texttt{propext}, \texttt{Classical.choice}, \texttt{Quot.sound}; zero
\texttt{sorryAx}) kernel-check the core general-$(r,c,d)$ theory for
\emph{all} degrees $r\ge1$, leading coefficients $c\ge1$, and shifts $d\in\Z$:
\begin{itemize}[leftmargin=1.8em]
\item Lemma~\ref{lem:poly-growth}(i)--(iii) and decay of $m^\alpha/Q_m$
(\texttt{poly\_\allowbreak Q\_\allowbreak ge\_\allowbreak factorial},
\texttt{poly\_\allowbreak pow\_\allowbreak div\_\allowbreak Q\_\allowbreak tendsto\_\allowbreak zero});
\item Lemma~\ref{lem:poly-lattice}(i)--(iv) including the Casorati identity,
fundamental solution error bounds, irrationality $\theta\notin\Q$, and
zero-level lattice equivalence
(\texttt{poly\_\allowbreak casorati},
\texttt{alphaU\_\allowbreak error\_\allowbreak bound},
\texttt{alphaV\_\allowbreak error\_\allowbreak bound},
\texttt{poly\_\allowbreak alpha\_\allowbreak linear\_\allowbreak relation},
\texttt{theta\_\allowbreak irrational},
\texttt{polyApprExponent\_\allowbreak eq\_\allowbreak zero\_\allowbreak tfae},
\texttt{polyApprExponent\_\allowbreak eq\_\allowbreak zero\_\allowbreak iff\_\allowbreak mem\_\allowbreak latticePoly});
\item Theorem~\ref{thm:general-rec}(i)--(ii), establishing the full discrete
spectrum $\Lam^{(r,c,d)}(\R)=\frac{1}{c}\Z_{\ge0}\cup\{\infty\}$ for all
$r\ge1$
(\texttt{poly\_\allowbreak signed\_\allowbreak defect\_\allowbreak approx},
\texttt{poly\_\allowbreak abs\_\allowbreak defect\_\allowbreak approx},
\texttt{polyApprExponent\_\allowbreak xPlusPoly},
\texttt{poly\_\allowbreak spectral\_\allowbreak gap\_\allowbreak open\_\allowbreak interval},
\texttt{poly\_\allowbreak levelSet\_\allowbreak eq\_\allowbreak on\_\allowbreak gap},
\texttt{poly\_\allowbreak levelSet\_\allowbreak below\_\allowbreak eq\_\allowbreak lattice},
\texttt{poly\_\allowbreak range\_\allowbreak apprExponent});
\item Theorem~\ref{thm:poly-gauge-trichotomy}(iii), establishing the signed
exact-limit coset collapse $E_{\mathrm{lim}}^{(r,c,d),\pm}(k/c)=x_k^\pm+L^{(r,c,d)}$
(\texttt{xPlusPoly\_\allowbreak mem\_\allowbreak signedExactLimitSet},
\texttt{neg\_\allowbreak xPlusPoly\_\allowbreak mem\_\allowbreak signedExactLimitSet},
\texttt{poly\_\allowbreak trichotomy\_\allowbreak part\_\allowbreak iii\_\allowbreak complete});
\item Theorem~\ref{thm:poly-phase-transition}(i) (critical weight $\alpha=r$:
dimension $0$ for all finite level sets, Lebesgue nullity of
$F^{(r,c,d)}$, full Lebesgue measure of $F_\infty^{(r,c,d)}$, and
$\dimH(F_\infty^{(r,c,d)}\cap[a,b])=1$:
\texttt{poly\_\allowbreak dimH\_\allowbreak levelSet\_\allowbreak finite},
\texttt{poly\_\allowbreak dimH\_\allowbreak finiteLevel\_\allowbreak zero},
\texttt{poly\_\allowbreak volume\_\allowbreak finiteLevel\_\allowbreak zero},
\texttt{poly\_\allowbreak ae\_\allowbreak apprExponent\_\allowbreak eq\_\allowbreak top},
\texttt{poly\_\allowbreak volume\_\allowbreak infiniteLevel\_\allowbreak Icc},
\texttt{poly\_\allowbreak dimH\_\allowbreak infiniteLevel\_\allowbreak inter\_\allowbreak Icc\_\allowbreak eq\_\allowbreak one},
\texttt{poly\_\allowbreak dimH\_\allowbreak infiniteLevel\_\allowbreak eq\_\allowbreak one})
and Theorem~\ref{thm:poly-phase-transition}(iii) (supercritical collapse
$\alpha>r$:
\texttt{poly\_\allowbreak supercritical\_\allowbreak levelSet\_\allowbreak eq\_\allowbreak lattice},
\texttt{poly\_\allowbreak supercritical\_\allowbreak dimH\_\allowbreak zero},
\texttt{poly\_\allowbreak supercritical\_\allowbreak empty\_\allowbreak neg}).
\end{itemize}
The module \texttt{PolyGaugeAndPhaseTransition.lean} formalizes the
general critical gauges, the logarithmic scale, and the metric measure
statements in Theorem~\ref{thm:poly-gauge-trichotomy}(i)--(ii) and
Corollary~\ref{cor:poly-exact-gauge}. It provides strict-sublevel nullity,
local infinitude and $\sigma$-finiteness at the critical gauges on
$\R$ and $\R/\Z$, as well as the fast, slow and $\liminf$ comparisons.
In \texttt{PolyWeightedUpper.lean},
\texttt{dimH\_\allowbreak polyWeightedExponent\_\allowbreak le} verifies
the upper bound of Theorem~\ref{thm:poly-phase-transition}(ii) for every
$P:\texttt{PolyRecData}$ and $\alpha\le r$. It uses recurrence defects to count successive pairs of rounding
integers. The general lower bound in \texttt{PolyWeightedLower.lean}
uses the same nested intervals as Step~2. The proof obtains the uniform
mass bound by counting finite extensions of a prefix and applying
compactness to finite covers. The exact dimensions of the prescribed-limit
set and the sublevel are verified by
\texttt{dimH\_\allowbreak polyExactLimit\_\allowbreak eq} and
\texttt{dimH\_\allowbreak polyWeightedSublevel\_\allowbreak eq}.

\section{Continued-fraction window constraints and digit residues}\label{sec:windows-residues}

\subsection{Dimensions of continued-fraction window constraints}

Put $H_j=\lfloor Q_j\sqrt j\rfloor$. For irrational $x\in(0,1)$ let
$k(j)$ be its last convergent index with denominator at most $H_j$
(the integer part has index zero) and $r_j=k(j+1)-k(j)$.

\begin{theorem}[Hermite window constraints]\label{thm:windows}
For each prescribed sequence $\epsilon_j\in\{0,1\}$, the set with
$k(j)\equiv\epsilon_j\pmod2$ eventually has Hausdorff dimension one.
The set with eventually bounded $r_j$ has dimension $1/2$.
For every bounded sequence of positive integers $(m_j)$, the set
with $r_j=m_j$ eventually has dimension $1/2$. At a sufficiently
large starting index, either initial parity of $k(j)$ can be imposed.
\end{theorem}

\begin{proof}
Since $Q_{j+1}=(4j+2)Q_j+Q_{j-1}$ gives $Q_{j+1}/Q_j=4j+O(1)$, the window
thresholds $H_j=\lfloor Q_j\sqrt j\rfloor$ satisfy $H_{j+1}/H_j=4j+O(1)\asymp j$
and $\log H_j=j\log j+O(j)$. We work with fundamental cylinders of finite
regular continued fractions: for positive integers $a_1,\dots,a_r\ge1$, let
$I(a_1,\dots,a_r)$ be the closure of the set of irrationals in $(0,1)$ whose
first $r$ partial quotients are $a_1,\dots,a_r$. If $p_{r-1}/q_{r-1}$ and
$p_r/q_r$ are the last two convergents of $[0;a_1,\dots,a_r]$ and we write
$u:=q_r$, $v:=q_{r-1}$, then $I(a_1,\dots,a_r)$ is the closed interval with
endpoints $\frac{p_r}{u}$ and $\frac{p_r+p_{r-1}}{u+v}$, of length
\[
\bigl|I(a_1,\dots,a_r)\bigr|=\frac{1}{u(u+v)}\in\left[\frac{1}{2u^2},\frac{1}{u^2}\right].
\]
At stage $j$ we stop each constructed cylinder at the first convergent index
$\tau_j=k(j)+1$ whose denominator $u=q_{\tau_j}$ crosses $H_j$, maintaining the
two-sided invariant
\[
\frac43 H_j\le u\le \frac52 H_j.
\]
In the construction prescribing parity, the stopped cylinder will also have
last partial quotient $a_{\tau_j}=4$, so that $v=q_{\tau_j-1}$ satisfies
$u=4v+q_{\tau_j-2}>4v$, i.e.\ $v<u/4$. To initialize either construction at any
sufficiently large starting stage $j_0$, choose an integer $d$ with
$H_{j_0}/3\le d\le H_{j_0}/2$: the rational $1/d$ has the two terminal continued
fraction representations $[0;d]$ (length $1$, terminal denominator $d$) and
$[0;d-1,1]$ (length $2$, terminal denominator $(d-1)\cdot1+1=d$). Choosing the
representation whose terminal index has parity $\epsilon_{j_0}$ and appending
the digit $4$ gives a stopped cylinder at index $\tau_{j_0}=k(j_0)+1$ with
$k(j_0)\equiv\epsilon_{j_0}\pmod2$ and denominator
$u\in\{4d+1,4d+(d-1)\}\subset[4H_{j_0}/3,5H_{j_0}/2]$.

\emph{Full-dimensional parity construction.} Fix stage $j\ge j_0$ and a parent
cylinder stopped at $\tau_j=k(j)+1$ with last two denominators $u=q_{\tau_j}$
and $v=q_{\tau_j-1}\le u/4$. Write $T:=H_{j+1}$ and $R:=T/u\asymp j$. We choose
coprime positive integers $(m,n)$ in the rectangle
\[
\mathcal R_R:=\bigl[0.34R,0.36R\bigr]\times\bigl[0.20R,0.22R\bigr]\cap\Z^2.
\]
By M\"obius inversion $\sum_{d\mid\gcd(m,n)}\mu(d)=\mathbf 1_{\gcd(m,n)=1}$, the
number of coprime integer pairs in $\mathcal R_R$ is
\begin{align*}
\#\{(m,n)\in\mathcal R_R:\gcd(m,n)=1\}
&=\sum_{d=1}^{\lfloor0.36R\rfloor}\mu(d)
\left(\left\lfloor\frac{0.36R}{d}\right\rfloor-\left\lceil\frac{0.34R}{d}\right\rceil+1\right)\\
&\qquad\times\left(\left\lfloor\frac{0.22R}{d}\right\rfloor-\left\lceil\frac{0.20R}{d}\right\rceil+1\right)\\
&=\sum_{d=1}^{\lfloor0.36R\rfloor}\mu(d)
\left(\frac{0.02R}{d}+O(1)\right)\left(\frac{0.02R}{d}+O(1)\right)\\
&=0.0004R^2\sum_{d=1}^{\lfloor0.36R\rfloor}\frac{\mu(d)}{d^2}
+O\left(R\sum_{d=1}^{\lfloor0.36R\rfloor}\frac1d\right)\\
&=\frac{0.0004}{\zeta(2)}R^2+O(R\log R),
\end{align*}
since $\sum_{d>\lfloor0.36R\rfloor}d^{-2}=O(R^{-1})$. Thus there are $\gg R^2\asymp j^2$
coprime pairs $(m,n)\in\mathcal R_R$ uniformly for all large $j$.

For each coprime pair $(m,n)\in\mathcal R_R$, we have
$1<\frac{0.34}{0.22}\le\frac{m}{n}\le\frac{0.36}{0.20}=1.8$, so $m/n>1$ has a
standard finite regular continued-fraction expansion $[b_1;b_2,\dots,b_s]$ with
$b_1=\lfloor m/n\rfloor=1$, $s\ge2$, and $b_s\ge2$, as well as the paired
expansion $[b_1;b_2,\dots,b_s-1,1]$ of length $s+1$. By the standard matrix
identity for continuants, if $\begin{pmatrix}m & m'\\ n & n'\end{pmatrix}
=\begin{pmatrix}b_1 & 1\\ 1 & 0\end{pmatrix}\cdots\begin{pmatrix}b_s & 1\\ 1 & 0\end{pmatrix}$,
then replacing $\begin{pmatrix}b_s & 1\\ 1 & 0\end{pmatrix}$ by
$\begin{pmatrix}b_s-1 & 1\\ 1 & 0\end{pmatrix}\begin{pmatrix}1 & 1\\ 1 & 0\end{pmatrix}
=\begin{pmatrix}b_s & b_s-1\\ 1 & 1\end{pmatrix}$ leaves the first column
$\binom{m}{n}$ unchanged. Consequently, when we append either representation of
$m/n$ to the parent cylinder (whose denominator matrix has final row $(u,v)$),
the resulting terminal denominator is
\[
t=mu+nv,
\]
independently of which of the two terminal representations of $m/n$ is chosen
(and every intermediate denominator along the appended tail is strictly less
than $t$, while the penultimate denominator $t'=m'u+n'v$ satisfies $1\le t'<t$).
Using $0.34R\le m\le0.36R$, $0.20R\le n\le0.22R$, $u>0$, $0<v\le u/4$, and
$Ru=T$, we obtain the two-sided inequality
\[
\frac13 T<0.34T=0.34Ru<mu+nv\le\left(0.36+\frac{0.22}{4}\right)Ru=0.415T<\frac12 T.
\]
Thus $t\in(T/3,T/2)$, so the appended tail of $m/n$ has not yet crossed
$T=H_{j+1}$. Choose the representation of $m/n$ (of length $s$ or $s+1$) that
makes the terminal convergent index have parity $\epsilon_{j+1}\pmod2$, and
then append the digit $4$. The new terminal denominator is
\[
u_{\mathrm{new}}=4t+t',
\]
where $0<t'<t$ implies
\[
\frac43 T<4t<u_{\mathrm{new}}=4t+t'<5t<\frac52 T,
\]
and the new penultimate denominator $v_{\mathrm{new}}=t$ satisfies
$v_{\mathrm{new}}<u_{\mathrm{new}}/4$. Thus the first crossing of $T=H_{j+1}$
occurs precisely at the appended digit $4$ (index $\tau_{j+1}=k(j+1)+1$), so the
index immediately preceding the $4$ is indeed $k(j+1)$ and satisfies
$k(j+1)\equiv\epsilon_{j+1}\pmod2$, while $(u_{\mathrm{new}},v_{\mathrm{new}})$
restores the stage-$(j+1)$ invariant. Moreover, distinct coprime pairs
$(m,n)\in\mathcal R_R$ produce stopped child cylinders with disjoint interiors:
each child word is stopped at its first denominator crossing of $T$, so no child
word is a prefix of another, and deleting the parent prefix and the final digit
$4$ recovers the chosen representation of $m/n$ and hence the unique coprime
pair $(m,n)$.

\emph{Prescribed bounded increment construction.} Given a bounded sequence of
positive integers $(m_j)$ ($1\le m_j\le M_*$), start from a stage-$j$ cylinder
stopped at $\tau_j=k(j)+1$ with denominator $4H_j/3\le u\le5H_j/2$ and previous
denominator $0<v<u$. Append $m_j-1$ copies of the digit $1$ (appending no digit
if $m_j=1$), and let $u_*,v_*$ be the resulting last two denominators. By the
Fibonacci continuant relation, $u\le u_*\le F_{m_j+1}u\le F_{M_*+1}u$, so for all
sufficiently large $j$ we have $u_*<T/2$ ($T=H_{j+1}$), meaning none of the
$m_j-1$ ones crosses $T$. Next choose a final digit $a\in\Z_{\ge1}$ such that
the next denominator $u_{\mathrm{new}}=au_*+v_*$ satisfies
\[
\frac43 T\le au_*+v_*\le\frac52 T.
\]
Since $0<v_*<u_*$, this condition holds for all integers $a$ in an interval of
length $(7/6)T/u_*\asymp j$, giving $\gg T/u_*\gg j$ eligible choices of $a$.
For each such $a$, the first crossing of $T=H_{j+1}$ occurs at this final digit
$a$, so $\tau_{j+1}=\tau_j+m_j$, i.e.\ $r_j=k(j+1)-k(j)=\tau_{j+1}-\tau_j=m_j$,
and the new stopped denominator $u_{\mathrm{new}}\in[4T/3,5T/2]$ restores the
inductive hypothesis. Note that $k(j+1)=k(j)+m_j$ is determined thereafter by
the initial parity $k(j_0)\pmod2$ and the cumulative sums of $m_i$; in
particular, taking $m_j\equiv3\pmod2$ alternates the parity of $k(j)$ and
preserves the condition that $k(j)+j$ is odd.

\emph{Mass distribution lower bounds.} In both constructions, assign equal
conditional probability to the retained children of each level-$j$ stopped
cylinder. Let $b=2$ for the full-dimensional parity construction (where each
parent has $\gg R^2\gg j^2$ children) and $b=1$ for the prescribed-count
construction (where each parent has $\gg j$ children). By stage $j$, the
measure $\mu(I_j)$ of any retained stopped cylinder $I_j$ satisfies
\[
\mu(I_j)\le\prod_{i=j_0}^{j-1}\bigl(C_1 i^{-b}\bigr)=\exp\bigl(-bj\log j+O(j)\bigr).
\]
Every retained level-$j$ cylinder $I_j$ has last denominator
$u\in[4H_j/3,5H_j/2]$, so its length $|I_j|=1/(u(u+v))$ satisfies
$(25/2)^{-1}H_j^{-2}\le|I_j|\le(9/16)H_j^{-2}$, and the interiors of distinct
level-$j$ cylinders are disjoint. Consequently, any interval $U\subset(0,1)$ of
length $h\in[H_{j+1}^{-2},H_j^{-2}]$ can intersect the interiors of at most
$\lceil h/((25/2)^{-1}H_j^{-2})\rceil+2\le 15$ level-$j$ cylinders (and its
endpoints meet at most $2$ more). Since
$\log(h^{-1})\le2\log H_{j+1}=2j\log j+O(j)$, for every $0<s<b/2$ we have
\[
\mu(U)\le 17\exp\bigl(-bj\log j+O(j)\bigr)
=O\bigl(\exp(-2s j\log j- (b-2s)j\log j+O(j))\bigr)=O(h^s).
\]
By the mass distribution principle, the Hausdorff dimension is at least $b/2$
 for every $s<b/2$, yielding $\dimH\ge1$ in the first construction (hence $=1$)
and $\dimH\ge1/2$ in the second.

\emph{Upper bound $1/2$ for eventually bounded $r_j$.} If $r_j(x)$ is
eventually bounded, then there exist integers $C\ge1$ and $N_0\ge1$ such that
$k(j)\le Cj$ for all $j\ge N_0$. For any $n\ge2CN_0$, setting
$j=\lfloor n/(2C)\rfloor\ge N_0$ gives $k(j)\le Cj\le n/2<n$, so the $n$th
convergent denominator $d_n(x)=q_n(x)$ has already passed $H_j$:
\[
\frac{\log d_n(x)}{n}\ge\frac{\log H_{\lfloor n/(2C)\rfloor}}{n}
=\frac{\frac{n}{2C}\log\frac{n}{2C}+O(n)}{n}\longrightarrow\infty
\qquad(n\to\infty).
\]
Thus for every $M>0$, every point $x$ with eventually bounded $r_j(x)$ satisfies
$d_n(x)\ge\ee^{Mn}$ for all sufficiently large $n$. Fix any $s>1/2$ and choose
$t\in(1/2,s)$. Since $d_n(a_1,\dots,a_n)\ge\prod_{i=1}^n a_i$, summing
$d_n^{-2t}$ over all $(a_1,\dots,a_n)\in\Z_{\ge1}^n$ gives
\[
\sum_{a_1,\dots,a_n\ge1} d_n(a_1,\dots,a_n)^{-2t}
\le\sum_{a_1,\dots,a_n\ge1}\prod_{i=1}^n a_i^{-2t}=\zeta(2t)^n<\infty.
\]
For the subfamily of level-$n$ cylinders $I(a_1,\dots,a_n)$ with
$d_n\ge\ee^{Mn}$, each cylinder has diameter
$|I(a_1,\dots,a_n)|\le d_n^{-2}\le\ee^{-2Mn}$, and their $s$-dimensional
covering cost is bounded by
\[
\sum_{d_n\ge\ee^{Mn}}\bigl|I(a_1,\dots,a_n)\bigr|^s
\le\sum_{d_n\ge\ee^{Mn}} d_n^{-2s}
\le\ee^{-2M(s-t)n}\sum_{a_1,\dots,a_n\ge1} d_n^{-2t}
\le\bigl(\zeta(2t)\ee^{-2M(s-t)}\bigr)^n.
\]
Choosing $M>\frac{\log\zeta(2t)}{2(s-t)}$ makes the base
$\zeta(2t)\ee^{-2M(s-t)}<1$, so the $s$-cost tends to $0$ as $n\to\infty$.
Taking the countable union over starting indices $N_1$ of the tail sets
$\bigcap_{n\ge N_1}\{d_n\ge\ee^{Mn}\}$ shows that the set of points with
eventually bounded $r_j$ has $\mathcal H^s$-measure zero for all $s>1/2$, hence
Hausdorff dimension at most $1/2$. Together with the lower bound $1/2$, both
sets have exact Hausdorff dimension $1/2$.
\end{proof}

For related spectra of continued-fraction growth see Fan, Liao, Wang
and Wu~\cite{FLLW}. The lower constructions here prescribe the
Hermite windows themselves. Removing the countable set $\Q+\Q\ee$
leaves both dimensions unchanged. These statements concern sets of
variable $x$ and leave the window pattern of the single point $\pi$
undetermined (translation by its integer part preserves denominators).

\subsection{Short continued-fraction windows and divergent window counts}

For $x\in(0,1)\setminus\Q$, retain $H_j,k(j),r_j$ from the window
constraints above.

\begin{theorem}[Short windows at almost every point]\label{thm:short-ae}
There is a constant $c>0$ such that for Lebesgue-almost every $x\in(0,1)$,
\[
 \#\{1\le j\le J:r_j(x)\le1\}\ge c\log J
\]
for all sufficiently large $J$ (with a starting index depending on $x$).
\end{theorem}

\begin{proof}
For irrational $x=[0;a_1,a_2,\dots]\in(0,1)\setminus\Q$, let
$\tau_j(x):=k(j)+1$ be the first convergent index whose denominator $q_{\tau_j}$
strictly exceeds $H_j$. Stop each continued-fraction cylinder at index $\tau_j$
and denote by $\mathcal F_j$ the $\sigma$-field generated by the countable
partition of $(0,1)\setminus\Q$ into stopped cylinders
$I(a_1,\dots,a_{\tau_j})$. Because the stopping condition is determined by the
first crossing of $H_j$, the stopping words at level $j$ are prefix-free and cover
$(0,1)\setminus\Q$. Since $H_j<H_{j+1}$, we have $\tau_j\le\tau_{j+1}$, so
$\mathcal F_j\subseteq\mathcal F_{j+1}$, and the short-window event
\[
E_j:=\bigl\{x\in(0,1)\setminus\Q:r_j(x)=\tau_{j+1}(x)-\tau_j(x)\le1\bigr\}
\]
belongs to $\mathcal F_{j+1}$.

Fix a stopped cylinder $I=I(a_1,\dots,a_r)$ in $\mathcal F_j$ (so $r=\tau_j$),
and write $p/q=p_r/q_r$ and $p'/q'=p_{r-1}/q_{r-1}$ for its last two
convergents, with $q>H_j$ and $0\le q'\le q$. Any irrational $x\in I$ is
uniquely parameterized by its remaining tail $t=[0;a_{r+1},a_{r+2},\dots]\in(0,1)\setminus\Q$
via the M\"obius transformation
\[
x=\Phi_I(t):=\frac{p+p't}{q+q't},\qquad t\in(0,1).
\]
Since $|pq'-p'q|=1$, the Jacobian derivative of $\Phi_I$ is
$|\Phi'_I(t)|=(q+q't)^{-2}$, while the total Lebesgue measure of $I$ is
$\operatorname{Leb}(I)=\int_0^1(q+q't)^{-2}\,dt=\frac{1}{q(q+q')}$. Dividing
$|\Phi'_I(t)|$ by $\operatorname{Leb}(I)$ shows that, conditionally on $I$, the
tail $t\in(0,1)$ has the explicit probability density
\[
p_I(t)=\frac{|\Phi'_I(t)|}{\operatorname{Leb}(I)}=\frac{q(q+q')}{(q+q't)^2},
\qquad 0<t<1.
\]
We now bound $\mathbb P(E_j\mid I)$. There are two cases:
\begin{itemize}[leftmargin=1.8em]
\item If $q>H_{j+1}$, then the denominator $q_{\tau_j}=q$ has already crossed
$H_{j+1}$, so $\tau_{j+1}=\tau_j$ and $r_j=0\le1$ identically on $I$; hence
$\mathbb P(E_j\mid I)=1$.
\item If $H_j<q\le H_{j+1}$, put $A:=\lfloor H_{j+1}/q\rfloor+1\ge2$, so that
$A-1\le H_{j+1}/q<A$. Whenever the next partial quotient $a=a_{r+1}$ satisfies
$a\ge A$, the next denominator is $q_{r+1}=aq+q'\ge Aq>H_{j+1}$, which forces
$\tau_{j+1}=r+1=\tau_j+1$ and hence $r_j=1$ ($E_j$ occurs). In terms of the tail
$t=[0;a,a_{r+2},\dots]\in(0,1)$, the condition $a=\lfloor 1/t\rfloor\ge A$ is
equivalent to $0<t<1/A$ (up to the rational endpoint $1/A$). Integrating the
conditional density $p_I(t)$ over $(0,1/A)$ gives
\[
\mathbb P(E_j\mid I)\ge\mathbb P(a_{r+1}\ge A\mid I)
=\int_0^{1/A}\frac{q(q+q')}{(q+q't)^2}\,dt
=\left[-\frac{q(q+q')}{q'(q+q't)}\right]_{t=0}^{t=1/A}
=\frac{q+q'}{Aq+q'}.
\]
Because $A\ge1$ and $q'\ge0$, we have $A(q+q')=Aq+Aq'\ge Aq+q'$, so
$\frac{q+q'}{Aq+q'}\ge\frac1A$. Moreover, since $q>H_j$ and $H_{j+1}/q\ge1$, we
have $A=\lfloor H_{j+1}/q\rfloor+1\le\frac{H_{j+1}}{q}+1\le\frac{2H_{j+1}}{q}<\frac{2H_{j+1}}{H_j}$,
which yields
\[
\mathbb P(E_j\mid I)\ge\frac{q+q'}{Aq+q'}\ge\frac1A\ge\frac{H_j}{2H_{j+1}}.
\]
\end{itemize}
Since $H_j/(2H_{j+1})<1$, the bound $\mathbb P(E_j\mid I)\ge H_j/(2H_{j+1})$
holds in the first case $q>H_{j+1}$ as well, and therefore holds uniformly on
every atom $I$ of $\mathcal F_j$:
\[
\mathbb P(E_j\mid\mathcal F_j)\ge\frac{H_j}{2H_{j+1}}.
\]
By \eqref{eq:rec} and $H_j=\lfloor Q_j\sqrt j\rfloor$, we have
$H_{j+1}/H_j=4j+O(1)$, so there exist an index $j_0\ge1$ and a constant
$c_0>0$ such that
\[
\mathbb P(E_j\mid\mathcal F_j)\ge\frac{c_0}{j}\qquad\text{almost surely for all }j\ge j_0.
\]

Write $Z_J:=\sum_{j=1}^J\mathbf 1_{E_j}$. Since $E_1,\dots,E_{j-1}$ are
$\mathcal F_j$-measurable, $Z_{j-1}$ is $\mathcal F_j$-measurable, while on binary
indicators $\mathbf 1_{E_j}\in\{0,1\}$ we have the exact identity
$\ee^{-\mathbf 1_{E_j}}=1-(1-\ee^{-1})\mathbf 1_{E_j}$. Using $1-u\le\ee^{-u}$
for all $u\in\R$, the one-step conditional expectation for $j\ge j_0$ satisfies
\[
\mathbb E\bigl[\ee^{-\mathbf 1_{E_j}}\bigm|\mathcal F_j\bigr]
=1-(1-\ee^{-1})\mathbb P(E_j\mid\mathcal F_j)
\le 1-\frac{(1-\ee^{-1})c_0}{j}
\le\exp\left(-\frac{b}{j}\right),\qquad b:=(1-\ee^{-1})c_0>0.
\]
Iterating the tower property $\mathbb E[\ee^{-Z_j}]=\mathbb E\bigl[\ee^{-Z_{j-1}}\mathbb E[\ee^{-\mathbf 1_{E_j}}\mid\mathcal F_j]\bigr]$
from $j=j_0$ to $J$ and using $\sum_{j=j_0}^J\frac1j\ge\log J-\log j_0$ gives
\[
\mathbb E\bigl[\ee^{-Z_J}\bigr]
\le\exp\left(-b\sum_{j=j_0}^J\frac1j\right)
\le C J^{-b},\qquad C:=\ee^{b\log j_0}.
\]
By Markov's inequality applied to the nonnegative random variable $\ee^{-Z_J}$,
\[
\mathbb P\left(Z_J<\frac{b}{2}\log J\right)
=\mathbb P\left(\ee^{-Z_J}>J^{-b/2}\right)
\le J^{b/2}\mathbb E\bigl[\ee^{-Z_J}\bigr]
\le C J^{-b/2}.
\]
Evaluating along the dyadic subsequence $J_m:=2^m$ ($m\ge1$), the probabilities
$\mathbb P(Z_{2^m}<(b/2)\log(2^m))\le C\cdot(2^{-b/2})^m$ form a convergent
geometric series $\sum_{m=1}^\infty C(2^{-b/2})^m<\infty$ since $b>0$. By the
first Borel--Cantelli lemma, for Lebesgue-almost every $x\in(0,1)$ there exists
$m_0(x)\ge2$ such that $Z_{2^m}(x)\ge\frac{b}{2}\log(2^m)$ for all
$m\ge m_0(x)$. Finally, for any integer $J\ge 2^{m_0(x)}$, choose $m\ge m_0(x)$
with $2^m\le J<2^{m+1}$; since $J\mapsto Z_J(x)$ is nondecreasing and
$m\ge\frac{m+1}{2}>\frac{\log_2 J}{2}$ for $m\ge1$, we obtain
\[
Z_J(x)\ge Z_{2^m}(x)\ge\frac{b}{2}\log(2^m)>\frac{b}{4}\log(2^{m+1})>\frac{b}{4}\log J,
\]
proving the theorem with $c:=b/4>0$.
\end{proof}

The argument applies to any increasing integer thresholds tending to
infinity with successive ratios $O(j)$. The estimate concerns
Lebesgue measure on the variable $x$ and supplies a frequency bound
for the short-window criterion of the Phase-1 project.

\begin{theorem}[Full dimension for divergent window counts]\label{thm:no-short-dim}
Let $\varphi=(1+\sqrt5)/2$ and prescribe arbitrary
$\epsilon_j\in\{0,1\}$. For any $0<\eta<1/\log\varphi$, the set of
$x\in(0,1)$ satisfying both $r_j(x)\ge\eta\log j$ and
$k(j)\equiv\epsilon_j\pmod2$ eventually has Hausdorff dimension at
least $1-\eta\log\varphi$. The set satisfying the prescribed
eventual parity and $r_j(x)\to\infty$ has Hausdorff dimension one
and Lebesgue measure zero.
\end{theorem}

\begin{proof}
We modify the stopped-cylinder construction of Theorem~\ref{thm:windows},
initialized at a large stage $j_0$ with the prescribed parity
$k(j_0)\equiv\epsilon_{j_0}\pmod2$. At stage $j\ge j_0$, a retained cylinder ends
at index $\tau_j=k(j)+1$ with last partial quotient $a_{\tau_j}=4$ and last two
denominators $u=q_{\tau_j}$ and $v=q_{\tau_j-1}$, where
\[
\frac43 H_j\le u\le\frac52 H_j,\qquad 0<v\le\frac{u}{4}.
\]
First append $\ell_j:=\lceil\eta\log j\rceil$ partial quotients equal to $1$.
Write $u':=q_{\tau_j+\ell_j}$ and $v':=q_{\tau_j+\ell_j-1}$ for the resulting
last two denominators. Since multiplication by $\begin{pmatrix}1 & 1\\ 1 & 0\end{pmatrix}^{\ell_j}
=\begin{pmatrix}F_{\ell_j+1} & F_{\ell_j}\\ F_{\ell_j} & F_{\ell_j-1}\end{pmatrix}$
governs the denominator recurrence during a block of $\ell_j$ ones (where
$F_0=0,F_1=1,F_{i+1}=F_i+F_{i-1}$ are the Fibonacci numbers), we have
\[
\begin{pmatrix}u'\\ v'\end{pmatrix}
=\begin{pmatrix}F_{\ell_j+1} & F_{\ell_j}\\ F_{\ell_j} & F_{\ell_j-1}\end{pmatrix}
\begin{pmatrix}u\\ v\end{pmatrix},
\]
so $0<v'\le u'$ and, using $v\le u$ and $F_{\ell_j+1}+F_{\ell_j}=F_{\ell_j+2}\le\varphi^{\ell_j+1}$,
\[
u\le u'=F_{\ell_j+1}u+F_{\ell_j}v\le F_{\ell_j+2}u\le\varphi^{\ell_j+1}u
\le\varphi^2 j^{\eta\log\varphi}u.
\]
With $T:=H_{j+1}$ and $H_{j+1}/H_j=4j+O(1)$, the remaining scale ratio
$R:=T/u'$ satisfies
\[
R=\frac{T}{u'}\ge\frac{H_{j+1}}{\varphi^2 j^{\eta\log\varphi}\cdot(5H_j/2)}
\ge c_1 j^{1-\eta\log\varphi}
\]
for a constant $c_1>0$. Since $0<\eta<1/\log\varphi$, the exponent
$\gamma:=1-\eta\log\varphi>0$ is strictly positive, so $R\to\infty$ uniformly as
$j\to\infty$; in particular, $u'=T/R<T$ for all large $j$, so the prelude of
$\ell_j$ ones stays strictly below $T$.

Next, choose coprime positive integers $(m,n)$ in the rectangle
\[
\mathcal R'_R:=\bigl[0.34R,0.35R\bigr]\times\bigl[0.12R,0.13R\bigr]\cap\Z^2.
\]
By the same M\"obius inversion calculation as in Theorem~\ref{thm:windows}, the
number of coprime pairs $(m,n)\in\mathcal R'_R$ is
$\frac{0.01\times0.01}{\zeta(2)}R^2+O(R\log R)\gg R^2\ge c_2 j^{2\gamma}$. For
every $(m,n)\in\mathcal R'_R$, the quotient $m/n$ satisfies
\[
2<\frac{34}{13}=\frac{0.34R}{0.13R}\le\frac{m}{n}\le\frac{0.35R}{0.12R}=\frac{35}{12}<3.
\]
Thus $m/n\in(2,3)$ has first partial quotient $\lfloor m/n\rfloor=2$ and two
finite regular continued-fraction representations of opposite lengths modulo
two (both of length at least $2$). Append the digits of $m/n$ using the
representation whose terminal convergent index has parity
$\epsilon_{j+1}\pmod2$. As shown in Theorem~\ref{thm:windows}, both
representations produce the same terminal denominator
\[
t=mu'+nv'
\]
(and all intermediate denominators along the tail of $m/n$ are at most $t$).
Using $0.34R\le m\le0.35R$, $0.12R\le n\le0.13R$, $0<v'\le u'$, and $Ru'=T$, we
obtain
\[
\frac13 T<0.34T=0.34Ru'<t=mu'+nv'\le(0.35+0.13)Ru'=0.48T<\frac12 T.
\]
Thus $t\in(T/3,T/2)$ stays below $T=H_{j+1}$. Now append the final digit $4$.
The new denominator is $u_{\mathrm{new}}=4t+t'$ with $0<t'<t$, so
\[
\frac43 T<4t<u_{\mathrm{new}}=4t+t'<5t<\frac52 T
\]
and $v_{\mathrm{new}}=t<u_{\mathrm{new}}/4$. Hence the first crossing of
$T=H_{j+1}$ occurs at this final digit $4$ (index $\tau_{j+1}=k(j+1)+1$), so the
index immediately before the $4$ is $k(j+1)$ with
$k(j+1)\equiv\epsilon_{j+1}\pmod2$, and the window increment satisfies
\[
r_j=\tau_{j+1}-\tau_j=\ell_j+\operatorname{len}(m/n)+1\ge\ell_j+3>\eta\log j.
\]
Distinct coprime pairs $(m,n)\in\mathcal R'_R$ produce stopped child cylinders
with disjoint interiors because each child stops at its first crossing of $T$,
and deleting the parent prefix, the fixed prelude of $\ell_j$ ones, and the
terminal digit $4$ uniquely recovers the chosen representation of $m/n$.

Since each stage-$j$ parent has at least $c_2 j^{2\gamma}$ children, assigning
equal conditional masses gives each level-$j$ stopped cylinder a measure at most
$\prod_{i=j_0}^{j-1}(c_2^{-1}i^{-2\gamma})=\exp(-2\gamma j\log j+O(j))$, while
its length is comparable to $H_j^{-2}$. As in Theorem~\ref{thm:windows}, any
interval of length $h\in[H_{j+1}^{-2},H_j^{-2}]$ meets at most $17$ level-$j$
cylinders, and since $\log H_{j+1}=j\log j+O(j)$, its measure is $O(h^s)$ for
every $0<s<\gamma$. The mass distribution principle yields Hausdorff dimension
at least $\gamma=1-\eta\log\varphi$. Finally, every point in the constructed
set for a given $\eta>0$ satisfies $r_j(x)\ge\eta\log j\to\infty$ and
$k(j)\equiv\epsilon_j\pmod2$ eventually, so the set with $r_j(x)\to\infty$ and
prescribed eventual parity contains these sets for all $\eta\in(0,1/\log\varphi)$;
letting $\eta\downarrow0$ shows that it has Hausdorff dimension~$1$, while
Theorem~\ref{thm:short-ae} implies that it has Lebesgue measure zero (since
almost every $x$ has $r_j(x)\le1$ infinitely often).
\end{proof}

In particular, the set with eventually $r_j\ge3$ has Lebesgue
measure zero and full Hausdorff dimension~$1$. The same is true if one
also imposes any prescribed parity of the terminal convergent
indices.

\subsection{Eventual periodicity of digit residues}

For irrational $x=[0;a_1,a_2,\ldots]$ in $(0,1)$ let $E$ consist of
the points for which $(a_i\bmod4)$ is eventually periodic. Let $F$
consist of points with every digit in $\mathcal A=\{1,5,9,\ldots\}$.
Continued fractions with a fixed arithmetic-progression alphabet are
studied by Chousionis, Leykekhman and Urba\'nski~\cite{CLU}. Here
we compare different sequences of allowed residues, including sequences
that change from one position to the next.

\begin{theorem}[Digit residues modulo four]\label{thm:residues}
$\dimH E=\dimH F$, with $357/512<\dimH F<359/512$.
\end{theorem}

\begin{proof}
\emph{Step 1: Uniform Lipschitz equivalence and $\dimH E=\dimH F$.}
Fix any sequence of residue shifts $\mathbf c=(c_1,c_2,\dots)\in\{0,1,2,3\}^{\N}$,
and define the digit-shift map $\Psi_{\mathbf c}:F\to(0,1)\setminus\Q$ by
\[
\Psi_{\mathbf c}\bigl([0;a_1,a_2,\dots]\bigr):=[0;a_1+c_1,a_2+c_2,\dots]
\qquad(a_i\in\mathcal A=\{1,5,9,\dots\}).
\]
We prove that $\Psi_{\mathbf c}$ is Lipschitz on $F$ with Lipschitz constant
$40/3$, \emph{uniformly} over all shift sequences $\mathbf c$. Let
$x=[0;a_1,a_2,\dots]$ and $y=[0;b_1,b_2,\dots]$ be two distinct points of $F$,
and let $n\ge0$ be the length of their longest common prefix (so $a_i=b_i$ for
$1\le i\le n$ and $a_{n+1}\ne b_{n+1}$). By symmetry we may assume
$a:=a_{n+1}<b:=b_{n+1}$. Since $a,b\in\mathcal A=\{1,5,9,\dots\}$ are distinct
positive integers congruent to $1\bmod4$, we have $a\ge1$ and $b-a\ge4$ (in
particular $b\ge5$). Write the source tails starting at position $n+1$ as
\[
x_t:=[0;a,x_1]=\frac{1}{a+x_1},\qquad
y_t:=[0;b,y_1]=\frac{1}{b+y_1}\qquad\bigl(x_1,y_1\in(0,1)\bigr),
\]
and the shifted image tails as
\begin{align*}
x'_t&:=[0;a+c,x'_1]=\frac{1}{a+c+x'_1},\\
y'_t&:=[0;b+c,y'_1]=\frac{1}{b+c+y'_1}\qquad\bigl(c:=c_{n+1}\in\{0,1,2,3\},\ x'_1,y'_1\in(0,1)\bigr).
\end{align*}
Since $0<x_1,y_1<1$ and $a+1<b$, the source tail distance is bounded below by
\[
|x_t-y_t|=\frac{1}{a+x_1}-\frac{1}{b+y_1}
>\frac{1}{a+1}-\frac{1}{b}=\frac{b-a-1}{(a+1)b}>0.
\]
Similarly, since $0<x'_1,y'_1<1$ and $c\ge0$, the shifted tail distance is
bounded above by
\[
|x'_t-y'_t|=\frac{1}{a+c+x'_1}-\frac{1}{b+c+y'_1}
<\frac{1}{a+c}-\frac{1}{b+c+1}=\frac{b-a+1}{(a+c)(b+c+1)}
\le\frac{b-a+1}{a(b+1)}.
\]
Combining these two inequalities bounds the tail distance ratio by
\[
\frac{|x'_t-y'_t|}{|x_t-y_t|}
\le\frac{b-a+1}{b-a-1}\cdot\frac{(a+1)b}{a(b+1)}.
\]
Put $d:=b-a\ge4$. Then $\frac{d+1}{d-1}=1+\frac{2}{d-1}$ is strictly decreasing
in $d\ge4$, so $\frac{b-a+1}{b-a-1}\le\frac{4+1}{4-1}=\frac53$. Moreover,
since $b\ge a+4$ and $a\ge1$, we have
\[
\frac{(a+1)b}{a(b+1)}
=\frac{a+1}{a}\cdot\frac{b}{b+1}
<2\cdot1=2
\quad\left(\text{and in fact }\frac{b-a+1}{b-a-1}\frac{(a+1)b}{a(b+1)}\le\frac53\cdot\frac{2\cdot5}{1\cdot6}=\frac{50}{18}<\frac{10}{3}\right),
\]
so in all cases
\[
|x'_t-y'_t|\le\frac{10}{3}|x_t-y_t|.
\]
Now restore the common prefix of length $n$. For the source points,
$x=M(x_t)$ and $y=M(y_t)$, where $M(t)=\frac{p_n+p_{n-1}t}{q_n+q_{n-1}t}$ is the
M\"obius map of the common prefix $(a_1,\dots,a_n)$, with last denominator
$q:=q_n\ge1$ and previous denominator $0\le q_{n-1}\le q$. For all $t\in[0,1]$,
$q\le q+q_{n-1}t\le 2q$, so the derivative $|M'(t)|=(q+q_{n-1}t)^{-2}$ satisfies
\[
\frac{1}{4q^2}\le |M'(t)|\le\frac{1}{q^2}\qquad\text{for all }t\in[0,1].
\]
By the mean value theorem, $|x-y|=|M(x_t)-M(y_t)|\ge\frac{1}{4q^2}|x_t-y_t|$.
For the shifted points, $\Psi_{\mathbf c}(x)=M'(x'_t)$ and
$\Psi_{\mathbf c}(y)=M'(y'_t)$, where $M'$ is the M\"obius map of the shifted
prefix $(a_1+c_1,\dots,a_n+c_n)$ with last denominator $q_{\mathrm{new}}$.
Because the continuant $q_n(d_1,\dots,d_n)$ is strictly increasing in each
partial quotient $d_i\ge1$ and $a_i+c_i\ge a_i$, we have
$q_{\mathrm{new}}\ge q$, and the mean value theorem gives
\[
\bigl|\Psi_{\mathbf c}(x)-\Psi_{\mathbf c}(y)\bigr|
=\bigl|M'(x'_t)-M'(y'_t)\bigr|
\le\frac{1}{q_{\mathrm{new}}^2}|x'_t-y'_t|
\le\frac{1}{q^2}|x'_t-y'_t|.
\]
Combining the prefix and tail bounds yields the uniform Lipschitz inequality
\[
\bigl|\Psi_{\mathbf c}(x)-\Psi_{\mathbf c}(y)\bigr|
\le\frac{1}{q^2}\cdot\frac{10}{3}|x_t-y_t|
\le 4\cdot\frac{10}{3}|x-y|=\frac{40}{3}|x-y|\qquad\text{for all }x,y\in F.
\]
Hence $\dimH\Psi_{\mathbf c}(F)\le\dimH F$ for every shift sequence
$\mathbf c\in\{0,1,2,3\}^{\N}$. Conversely, every point
$x=[0;d_1,d_2,\ldots]\in E$ has $(d_i\bmod4)$ eventually periodic; setting
$c_i:=(d_i-1)\bmod4\in\{0,1,2,3\}$ gives an eventually periodic sequence
$\mathbf c=(c_1,c_2,\dots)\in\{0,1,2,3\}^{\N}$ such that
$a_i:=d_i-c_i=4\lfloor(d_i-1)/4\rfloor+1\in\mathcal A$ for every $i\ge1$, so
$x=\Psi_{\mathbf c}([0;a_1,a_2,\dots])\in\Psi_{\mathbf c}(F)$. Thus $E$ is the
union of $\Psi_{\mathbf c}(F)$ over the countable set of eventually periodic
sequences $\mathbf c\in\{0,1,2,3\}^{\N}$ (each determined by a preperiod length,
a period length, and finitely many residues in $\{0,1,2,3\}$). Countable
stability gives $\dimH E\le\dimH F$, while $F\subset E$ (corresponding to
$c_i\equiv0$) gives $\dimH F\le\dimH E$. Therefore $\dimH E=\dimH F$.

\emph{Step 2: Transfer operator, Bowen's formula, and integral tail bounds.}
For the numerical bounds on $\dimH F$, let $f_a(x)=1/(a+x)$ for $0\le x\le1$.
For an alphabet $\mathcal D\subseteq\mathcal A$ and $s>1/2$ define the
Perron--Frobenius transfer operator on bounded positive functions $v:[0,1]\to(0,\infty)$
by
\[
 (\mathcal L_{s,\mathcal D}v)(x)
   =\sum_{a\in\mathcal D}(a+x)^{-2s}v(f_a(x)).
\]
The topological pressure $P(s,\mathcal D)$ is the limit of
$n^{-1}\log(\mathcal L_{s,\mathcal D}^n1)(0)$ as $n\to\infty$.
Indeed, the derivative of an $n$-letter cylinder map
$f_{a_1}\circ\cdots\circ f_{a_n}(x)$ has magnitude $(q_n+xq_{n-1})^{-2}$, which
is maximal at $x=0$ and varies by at most a factor of $4$ on $[0,1]$. The
pressure $s\mapsto P(s,\mathcal D)$ is continuous and strictly decreasing on
$(1/2,\infty)$, and Bowen's formula identifies its unique zero with the
Hausdorff dimension of the limit set coded by $\mathcal D$
\cite[\S2, Theorem~2.3]{CLU}. To satisfy the strict uniform contraction
convention in that reference (since $f_1'(0)=-1$), one passes to the system of
all two-letter branches $f_a\circ f_b$ ($a,b\in\mathcal D$): their derivatives
satisfy $|(f_a\circ f_b)'(x)|=(ab+ax+1)^{-2}\le(1\cdot1+0+1)^{-2}=1/4<1$ on
$[0,1]$, their cylinder interiors are disjoint, their transfer operator is
$\mathcal L_{s,\mathcal D}^2$ (so its pressure is $2P(s,\mathcal D)$, with the
same zero), and the limit set coded by paired words in $\mathcal D^2$ is
identical to the limit set coded by $\mathcal D$.

For any bounded positive function $v:[0,1]\to(0,\infty)$ with
$\inf_{[0,1]}v>0$, a uniform pointwise upper inequality
$\mathcal L_{s,\mathcal D}v\le(1-\eta)v$ on $[0,1]$ ($0<\eta<1$) forces
$P(s,\mathcal D)\le\log(1-\eta)<0$: by positivity of $\mathcal L_{s,\mathcal D}$,
iteration gives $\mathcal L_{s,\mathcal D}^n v\le(1-\eta)^n v$, and comparing $1$
with $v/\inf v$ yields
$(\mathcal L_{s,\mathcal D}^n1)(0)\le(1-\eta)^n v(0)/\inf v$.
Likewise, a uniform pointwise lower inequality
$\mathcal L_{s,\mathcal D}v\ge(1+\eta)v$ on $[0,1]$ ($\eta>0$) forces
$(\mathcal L_{s,\mathcal D}^n1)(0)\ge(1+\eta)^n v(0)/\sup v$ and hence
$P(s,\mathcal D)\ge\log(1+\eta)>0$. We apply the lower test at $s=357/512$ to
the finite subalphabet $\mathcal D_N:=\{4k+1:0\le k<N\}$ with $N=2^{24}$ (whose
limit set is contained in $F$, so $\dimH F\ge\dimH F_{\mathcal D_N}>357/512$),
and the upper test at $s=359/512$ to the full alphabet $\mathcal D=\mathcal A$
(certifying $\dimH F<359/512$).

In both tests, $v$ is chosen as a positive step function on the $256$ uniform
grid cells $C_j=[l_j,u_j)=[j/256,(j+1)/256)$ ($0\le j\le254$) and
$C_{255}=[255/256,1]$. On a cell $[l,u]$, for each explicit branch $a=4k+1$
($0\le k<64$) and $x\in[l,u]$, the weight $(a+x)^{-2s}$ is decreasing in $x$
and thus lies in $[(a+u)^{-2s},(a+l)^{-2s}]$, while $f_a(x)=1/(a+x)$ lies in the
closed image interval $[f_a(u),f_a(l)]=[1/(a+u),1/(a+l)]$; we bound $v(f_a(x))$
from below and above by the minimum and maximum of $v$ over all grid cells
intersected by $[f_a(u),f_a(l)]$. For all tail indices $k\ge64$, the digit is
$a=4k+1\ge257>256$, so for every $x\in[0,1]$ the image point satisfies
$0<f_a(x)=1/(4k+1+x)\le 1/257<1/256$ and therefore falls strictly inside the
first cell $C_0=[0,1/256)$, where $v$ takes the constant value $v_0$. Putting
$p:=2s>1$, we bound the tail sums over $k\ge64$ by monotone integral comparison:
for every $x\in[0,1]$ and every $t\in[k,k+1]$, we have
$4k+1+x\le 4k+2\le 4t+2$, so
\[
(4k+1+x)^{-p}\ge(4k+2)^{-p}\ge\int_k^{k+1}(4t+2)^{-p}\,dt,
\]
and summing from $k=64$ to $N-1$ (where $4\cdot64+2=258$) gives the lower tail
bound
\[
 \sum_{k=64}^{N-1}(4k+1+x)^{-p}v_0
 \ge v_0\int_{64}^N(4t+2)^{-p}\,dt
 =v_0\left[\frac{(4t+2)^{1-p}}{4(1-p)}\right]_{t=64}^{t=N}
 =v_0\,\frac{258^{1-p}-(4N+2)^{1-p}}{4(p-1)}.
\]
In the opposite direction, for every $x\in[0,1]$ and every $t\in[k-1,k]$, we
have $4k+1+x\ge 4k+1\ge 4t+1$, so
\[
(4k+1+x)^{-p}\le(4k+1)^{-p}\le\int_{k-1}^k(4t+1)^{-p}\,dt,
\]
and summing from $k=64$ to $\infty$ (where $4\cdot63+1=253$) gives the upper
tail bound
\[
 \sum_{k=64}^{\infty}(4k+1+x)^{-p}v_0
 \le v_0\int_{63}^\infty(4t+1)^{-p}\,dt
 =v_0\left[\frac{(4t+1)^{1-p}}{4(1-p)}\right]_{t=63}^{t=\infty}
 =v_0\,\frac{253^{1-p}}{4(p-1)} .
\]

\emph{Step 3: Exact integer-arithmetic verification.}
The companion script \texttt{residue\_certificate.py} constructs a
positive integer vector of cell values for each test. Every branch
weight and integral power is enclosed by rational endpoints. In
particular, to enclose $S(y/z)^{-m/d}$, where $S=10^{13}$ and
$y,z,m,d$ are positive integers, the script checks the exact
integer inequalities $L^d y^m\le S^d z^m\le U^d y^m$ for the proposed
integer endpoints $L,U$. Summation and comparison against the
corresponding cell value use exact integer arithmetic. At $s=357/512$, the
smallest relative lower margin $\min_j((\mathcal L_{s,\mathcal D_N}v-v)/v)|_{C_j}$
equals
$7204302042904417931/2465429597500000000000>29/10000>0$; at $s=359/512$, the
smallest relative upper margin $\min_j((v-\mathcal L_{s,\mathcal A}v)/v)|_{C_j}$
equals
$65040245891337248463/9860138780000000000000>65/10000>0$. The
strict signs certify the stated open interval $(357/512,359/512)$, of width
$1/256$.
\end{proof}

As a comparison, Table~1 of \cite{CLU} bounds the dimension for
$\{1,4,7,\ldots\}$ by $0.743586$. Its Proposition~2.7 transfers this upper bound
to $\{1,5,9,\ldots\}$ by the increasing bijection $1+3k\mapsto 1+4k$.

\paragraph{Lean status of Section~\ref{sec:windows-residues}.}
Fourteen core algebraic, conditional-probability, bi-Lipschitz, and
integer-certificate soundness lemmas underlying
Theorems~\ref{thm:windows}--\ref{thm:residues} are kernel-checked at Layer~A in
\texttt{EpiFormal/\allowbreak{}Window\allowbreak{}And\allowbreak{}Residue\allowbreak{}Core.lean}
(base axioms only, zero \texttt{sorryAx}), including the crossing bounds
$T/3<t<T/2$ and $4T/3<4t+t'<5T/2$, the conditional-probability and exponential
moment step inequalities, the tail ratio bound $10/3$ and combined prefix
bi-Lipschitz bound $40/3$, and the exact power-enclosure soundness
$L^d y^m\le S^d z^m\le U^d y^m\implies L/S\le(z/y)^{m/d}\le U/S$ for
\texttt{residue\_certificate.py}. The full continued-fraction measure
ergodicity and Hausdorff-dimension pressure arguments remain paper proofs.

\section{Numerical evidence}\label{sec:num}

All high-precision computations use \texttt{mpmath} with dynamic precision
$\mathrm{digits}(Q_j)+300$ (or 100 decimal digits for fixed-horizon checks up
to $j=35$) to survive the catastrophic cancellation in
$Q_jx-\mathrm{round}(Q_jx)$. These are corroborations, not proofs.

\paragraph{Spectral gap.} For $x=B\pi$ ($B=1,\dots,10$), $\sqrt2$, the golden
ratio $\varphi$, and random reals, $\max_{j\le60}j\|Q_jx\|$ lies in $22$--$29$
and grows $\sim j/4$ in this finite sample; the data do not establish an
infinite limsup. For $x\in L$ (e.g.\ $(\ee+1)/2$, $1+(\ee+1)/2$),
$j\|Q_jx\|<10^{-12}$ already for $j\ge10$ and $\sim10^{-117}$ at $j=60$,
confirming $\Lam=0$.

\paragraph{Dimension (box-counting).} Exact rational covering trees give the
following box-counting estimates for $E_c\cap[0,1]$:
\begin{center}
\begin{tabular}{c|cccc}
$c$ & $N=4$ & $N=12$ & path $N=400$ & bound $n=1500$\\\hline
$1/4$ & 0.371 & 0.323 & 0.171 & 0.180\\
$1/2$ & 0.630 & --- & 0.251 & 0.270\\
$1$   & 0.863 & --- & 0.343 & 0.360\\
$2$   & 1.000 & --- & 0.442 & 0.446
\end{tabular}
\end{center}
The finite counts suggest the branching scale $\lfloor8c\rfloor+1$ ($3,5,9,17$), and
the rigorous upper bound decays as $\log B/\log n\to0$. \textbf{Caveat:} the
decay is only $O(1/\log n)$; at accessible $N$ the estimates are far from $0$
(e.g.\ $s=0.01$ needs $n\approx10^{60}$). The numerics are consistent with
$\dimH=0$ but cannot resolve it directly; the proof is the covering argument of
Theorem~\ref{thm:finite}. Below the gap ($c=0.24$) the covering upper bound
still grows exponentially while the true branching tends to $0$: the covering
method is not sharp below the gap, which is why the exact classification
(Theorem~\ref{thm:class}) is needed there.

\paragraph{Discrete spectral ladder and three-way splitting.} The 100-digit
\texttt{mpmath} script \texttt{verify\_\allowbreak full\_\allowbreak spectrum\_\allowbreak discrete.py} verifies
the quantized levels $\Lam(x_k)=k/4$ for $k=1,2,3,4$ (with $|4ju_j(x_k)-k|\approx7.14\times10^{-3}\cdot k$ at $j=35$), confirms that two points in
$E_{\mathrm{lim}}^+(k/4)$ with distinct initial digits have difference in $L$ to
within $5.95\times10^{-97}$, and checks that the alternating binary digit choice
$d_j=(-1)^j k\in\{-k,+k\}$ in $E_{\mathrm{lim}}(k/4)$ outside $E_{\mathrm{lim}}^+(k/4)\cup E_{\mathrm{lim}}^-(k/4)$ satisfies $4j\|Q_jx\|\to k$ while $4ju_j(x)$ oscillates
between $+k$ and $-k$. The companion script \texttt{verify\_\allowbreak weighted\_\allowbreak lower.py}
confirms the Falconer tree-density ratio convergence toward $1-\alpha$ for
subcritical weights $\alpha\in\{0.25,0.50,0.75\}$.

\paragraph{General polynomial-growth recurrences.} The script
\texttt{verify\_\allowbreak poly\_\allowbreak recurrence\_\allowbreak 2b.py}
(at least $400$-digit effective \texttt{mpmath} precision for Experiments~1--3
and $50$-digit precision up to $N=10^5$ for Experiments~4--5) passes all
assertions across $r\in\{1,2,3\}$ and $c\in\{1,2,4\}$: it verifies the
non-asymptotic factorial lower bound of Lemma~\ref{lem:poly-growth}, the
$O(j^{-r})$ convergence of $j^r\|Q_jx\|\to k/c$ for $k\in\{1,3\}$, the
intrinsic lattice $L^{(r,c,d)}=\frac{1}{g}(\Z+\Z\theta)$ (including the
non-coprime case $(2,4,0)$ with $Q_0=3,Q_1=6,g=3$, where
$q_0\alpha_U+q_1\alpha_V=1/3$ to within $5.14\times10^{-1079}$ and $x_2-x_1\in
L^{(2,4,0)}$ to within $6.0\times10^{-1072}$) alongside alternating-digit
oscillation in $E_{\mathrm{lim}}^{(r,c,d)}(k/c)$, the subcritical Falconer
dimension trajectories approaching $1-\alpha/r$ (for both fixed $t>0$ and the
slowly decaying $t=0$ radius) versus decay toward $0$ at critical $\alpha=r$
($s_{10^5}\approx0.033$ for $(2,1,-1)$ and $0.056$ for $(3,2,-3)$), and the
exact gauge ratio $\log(2k+1)/\log2$ at $n=10^5$ (with the $17.6\%$--$39.5\%$
finite-$n$ excess over the asymptotic scale factor $1/r$ matching the
$O(\log\log n/\log n)$ subleading term $\frac{\log n+\log\log n+\log r}{\log n+(\log c-r)/r}$
to within $0.009$).

\section{Honesty statement, scope, and formal verification}\label{sec:verification}

\begin{itemize}[leftmargin=1.5em]
\item Theorems~\ref{thm:class} and~\ref{thm:gap} have short self-contained paper
proofs in Section~\ref{sec:class} and are kernel-checked at Layer~A in
\texttt{EpiFormal/\allowbreak Classification.lean}
(\texttt{theorem2\_classification} and \texttt{corollary3\_gap}), using only the
three Lean base axioms \texttt{propext}, \texttt{Classical.choice} and
\texttt{Quot.sound}; the former $\delta_j$ axioms are now proved theorems
(\texttt{DeltaSign.lean}).
\item Corollaries~\ref{thm:spectral} and~\ref{thm:sub} (for $c<1/4$), and the
countability and dimension-$0$ statements for $c<1/4$, are kernel-checked at
Layer~A in \texttt{EpiFormal/\allowbreak SpectralGap.lean} (five declarations,
no \texttt{sorryAx}). The complete spectrum discretization
$\Lam(\R)=\frac14\Z_{\ge0}\cup\{\infty\}$ and open-interval spectral-gap
exclusion (Theorem~\ref{thm:full-spectrum}), the gap constancy $E_c=E_{k/4}$
(Corollary~\ref{thm:sub}), the three-way digit rigidity and signed-limit
$L$-difference, countability, and dimension-$0$ lemmas
(Theorem~\ref{thm:trichotomy}), and the $r=1$ linear $(cj+d)$ defect
approximation lemmas are kernel-checked at Layer~A in
\texttt{EpiFormal/\allowbreak FullSpectrumDiscrete.lean} (15 audited
declarations, base axioms only, no \texttt{sorryAx}). The remaining signed-coset
collapse $E_{\mathrm{lim}}^\pm(k/4)=x_k^\pm+L$ in
Theorem~\ref{thm:trichotomy}(iii) is now likewise kernel-checked at Layer~A in
\texttt{EpiFormal/\allowbreak TrichotomyAndPolyExtension.lean} (8 declarations,
ending in \texttt{trichotomy\_\allowbreak part\_\allowbreak iii\_\allowbreak complete},
base axioms only, no \texttt{sorryAx}), completing Layer~A verification of all
three parts~(i)--(iii) of Theorem~\ref{thm:trichotomy}.
\item For the general polynomial-growth denominator recurrences ($r\ge1$) of
Section~\ref{sec:poly-rec}, Lemmas~\ref{lem:poly-growth}(i)--(iii) and
\ref{lem:poly-lattice}, Theorem~\ref{thm:general-rec} (full spectrum
$\frac{1}{c}\Z_{\ge0}\cup\{\infty\}$ for all $r\ge1$),
Theorem~\ref{thm:poly-gauge-trichotomy}(iii) (signed-coset collapse), and
Theorem~\ref{thm:poly-phase-transition}(i),(iii) (critical $\alpha=r$
dimension-$0$/full-measure complement and supercritical $\alpha>r$ collapse)
are kernel-checked at Layer~A in
\texttt{EpiFormal/\allowbreak TrichotomyAndPolyExtension.lean} (37 audited
\texttt{PolyRecData} declarations, base axioms only, no \texttt{sorryAx}).
The critical-gauge statements for general recurrences and the subcritical
dimension upper bound are verified by \texttt{PolyGaugeAndPhaseTransition.lean}
and \texttt{PolyWeightedUpper.lean}, respectively. The lower bound for arbitrary
$\texttt{PolyRecData}$ is verified in \texttt{PolyWeightedLower.lean}.
\item Theorem~\ref{thm:finite} ($\dimH E_c=0$ for all $c\in\R$) is \emph{fully
machine-checked at Layer~A} (\texttt{DimHCore.lean} + \texttt{DimHFinite.lean},
base axioms only: \texttt{propext}, \texttt{Classical.choice},
\texttt{Quot.sound}; no $\delta$ axioms, no \texttt{sorryAx}). The assembly uses
a surviving-center \texttt{Finset} recursion (\texttt{survSet}) rather than path
encoding; the growth estimate is explicit
(\texttt{tendsto\_pow\_div\_factorial\_atTop}), not asymptotic. Moreover, the
complementary infinite-level conclusions of
Corollary~\ref{cor:infinite-level}---Lebesgue-nullity of $\{\Lam<\infty\}$,
full Lebesgue measure on every interval $[a,b]$ and a.e.\ $\Lam=\infty$, and
$\dimH\{\Lam=\infty\}=1$---are likewise kernel-checked at Layer~A in
\texttt{EpiFormal/\allowbreak{}GaugeExact.lean}
(\texttt{volume\_finiteLevel\_zero}, \texttt{ae\_apprExponent\_eq\_top},
\texttt{volume\_infiniteLevel\_Icc}, and
\texttt{dimH\_infiniteLevel\_eq\_one}).
\item Theorems~\ref{thm:upper} ($0\le\alpha\le1$, all $c\in\R$)
and~\ref{conj:weight} ($0\le\alpha<1$, prescribed limits, sublevel sets for all
$c\ge0$, and the zero-level set) are \emph{fully machine-checked at Layer~A}
(both the upper and lower bounds). The upper bound is proved in
\texttt{Weighted\allowbreak Fiber.lean}, \texttt{Weighted\allowbreak DimHCore.lean},
\texttt{Weighted\allowbreak Upper.lean}, and \texttt{Weighted\allowbreak Final.lean}
(as \texttt{dimH\_\allowbreak levelSet\ensuremath{\alpha}\_le}); the
matching lower bound and exact dimension identities are proved in
\texttt{Weighted\allowbreak ExactDim.lean}
(as \texttt{one\_\allowbreak le\_\allowbreak hausdorff\allowbreak Measure\_\allowbreak cantorSet},
\texttt{dimH\_\allowbreak exact\allowbreak LimitSet\ensuremath{\alpha}\_ge},
\texttt{dimH\_\allowbreak exact\allowbreak LimitSet\ensuremath{\alpha}\_eq},
\texttt{dimH\_\allowbreak levelSet\ensuremath{\alpha}\_eq}, and
\texttt{dimH\_\allowbreak appr\allowbreak Exponent\ensuremath{\alpha}\_eq\_\allowbreak zero}; base
axioms only, no \texttt{sorryAx}). The existence of $x_0$ with $\Lam(x_0)=1/4$
(as \texttt{lambda\_\allowbreak x0\_\allowbreak eq\_\allowbreak quarter} in \texttt{Quarter\allowbreak Point\allowbreak Final.lean}) and the
unweighted-zero strictness witness
(\texttt{Unweighted\allowbreak Correction.\allowbreak witness} in Remark~\ref{rem:unweighted})
are likewise kernel-checked at Layer~A.
\item Theorem~\ref{thm:gauge}, Corollary~\ref{cor:exact-gauge},
Remark~\ref{rem:gauge-scale}, and Theorem~\ref{thm:trichotomy}(i)--(iii) are now
\emph{fully machine-checked at Layer~A} both on $\R$ and (for the metric gauge
statements of parts~(i)--(ii)) natively on the
quotient metric space $\R/\Z$ (\texttt{UnitAddCircle}) across
\texttt{EpiFormal/\allowbreak GaugeNull.lean} (3 declarations),
\texttt{EpiFormal/\allowbreak GaugeExact.lean} (22 declarations),
\texttt{EpiFormal/\allowbreak GaugeCircle.lean} (53 declarations),
\texttt{EpiFormal/\allowbreak FullSpectrumDiscrete.lean} (15 declarations), and
\texttt{EpiFormal/\allowbreak TrichotomyAndPolyExtension.lean} (8 part-(iii)
declarations), using only
the three Lean base axioms \texttt{propext}, \texttt{Classical.choice}, and
\texttt{Quot.sound}, with zero \texttt{sorryAx}. In particular,
\texttt{GaugeCircle.lean} explicitly constructs the piecewise-linear gauges
$h_k$ (\texttt{critGauge}), $g_2$ (\texttt{binaryGauge}), and $g_\beta$
(\texttt{interpGauge}), discharges all external continuity hypotheses for
Theorem~\ref{thm:gauge} and Remark~\ref{rem:gauge-scale}, constructively proves
open-interval and open-arc $\mathcal H^{h_k}$-infinitude
(Corollary~\ref{cor:exact-gauge}(i) on $\R$ and \texttt{UnitAddCircle}), and
verifies the binary Cantor-tree positivity, $\mathcal H^{g_2}$-infinitude on
every nonempty open subset of $\R/\Z$, $\sigma$-finiteness,
$\mathcal H^{h_k}$-nullity, and $g_\beta$ threshold dichotomy for
$E_{\mathrm{lim}}(k/4)$ in Theorem~\ref{thm:trichotomy}(ii), while
\texttt{TrichotomyAndPolyExtension.lean} proves the signed-limit coset collapse
$E_{\mathrm{lim}}^\pm(k/4)=x_k^\pm+L$ in Theorem~\ref{thm:trichotomy}(iii).
\item For the continued-fraction window and residue results
(Theorems~\ref{thm:windows}--\ref{thm:residues}), fourteen core algebraic,
conditional-probability, bi-Lipschitz, and integer-certificate soundness lemmas
are kernel-checked at Layer~A in
\texttt{EpiFormal/\allowbreak{}Window\allowbreak{}And\allowbreak{}Residue\allowbreak{}Core.lean};
the full continued-fraction measure ergodicity and Hausdorff-dimension pressure
arguments remain paper proofs, and the numerical interval $(357/512,359/512)$ in
Theorem~\ref{thm:residues} is certified by the reproducible exact
integer-arithmetic script \texttt{residue\_certificate.py}.
\item This article does \emph{not} prove the irrationality of $e+\pi$ (which
remains an open problem). It is a metric-theory spin-off; the connection is
that the zero-level classification machinery arose from the $e+\pi$ escape
program \cite{Phase1}.
\item The zero-level classification and initial gap (Theorems~\ref{thm:class}
and~\ref{thm:gap}) are a subsequence variant of the Kraaikamp--Liardet
criterion \cite{KL}, quoted as Lemma~4.1 of \cite{BDS}; see
Section~\ref{sec:prior}. The independent contributions concern the complete spectrum
discretization and three-way positive-level splitting
(Theorems~\ref{thm:full-spectrum}, \ref{thm:trichotomy},
and~\ref{thm:general-rec}--\ref{thm:poly-phase-transition}, with
Theorems~\ref{thm:full-spectrum}, \ref{thm:trichotomy}(i)--(iii),
\ref{thm:general-rec}, \ref{thm:poly-gauge-trichotomy}(iii), and
\ref{thm:poly-phase-transition}(i),(iii) kernel-checked in Lean~4), the critical gauge theorems
(Theorem~\ref{thm:gauge} and Corollaries~\ref{cor:exact-gauge}
and~\ref{cor:liminf-gauges}), the marked-component gauge
(Theorem~\ref{thm:finite-state-tail}), the machine-checked dimension-zero and subcritical dimension theorems
(Theorems~\ref{thm:finite}, \ref{thm:upper}, and~\ref{conj:weight}), and the
window and residue theorems (Theorems~\ref{thm:windows}--\ref{thm:residues}).
\end{itemize}

\paragraph{Limits of the current extensions.}
The finite-state gauge theorems distinguish zero, positive finite, and
infinite measures before taking the countable union over starting times.
The pressure and circulation theorems classify the gauges of the
stated frequency levels. Theorems~\ref{thm:frequency-strip} and
\ref{thm:markov-upper-class} constrain their speed of convergence.
The sharp upper and lower classes at the Perron frequency and
the cubic-root exponent at other interior frequencies are
established at their reference gauges.
Theorem~\ref{thm:sharp-cubic-amplitude} determines the transition
amplitude and isolates the only remaining pair of extremal
limits at equality; within it, Lean verifies only the three
algebraic lemmas of \texttt{SharpWindow.lean} (the cubic
factorization identity, the strict peak, and the residue-count
error bound), and its probabilistic and measure arguments are
paper proofs. The measure of this endpoint subset
and the classification at other gauges require additional estimates. Extending the finite-state construction to
arithmetic constraints in general recurrence families also requires a
realization argument. The two remaining polynomial-growth gauge cases
are paper proofs.

\section*{Use of AI and formal verification}
This work was carried out with substantial assistance from generative AI
systems (including ChatGPT and Codex), both for developing and drafting
mathematical arguments and text, preparing the numerical certificates, and
writing the Lean~4 formalization. The results marked as machine-checked (Layer~A)
were verified by the Lean~4 kernel, and their axiom dependencies were audited
with \texttt{\#print axioms} in \texttt{Audit.lean} (349 audited declarations
in total---327 depending only on the Lean base axioms \texttt{propext},
\texttt{Classical.choice}, and \texttt{Quot.sound}, and 22 phase-1 declarations
additionally using the transcendental/irrationality-measure whitelist axioms
\texttt{e\_transcendental}, \texttt{pi\_transcendental}, and
\texttt{mu\_e\_eq\_two}---with zero project-added axioms and zero
\texttt{sorryAx}). The additional modules in the verification package audit the general
critical-gauge results, the general subcritical dimension equalities,
and the upper-bound and countability arguments of
Theorem~\ref{thm:finite-state-tail} in \texttt{FiniteStateTail.lean}.
The remaining paper proofs---including the positive-measure lower bound of
Theorem~\ref{thm:finite-state-tail}, the full continued-fraction measure
ergodicity and Hausdorff-dimension pressure arguments in
Theorems~\ref{thm:windows}--\ref{thm:residues} (beyond the 14 core lemmas in
\texttt{WindowAndResidueCore.lean}), and the probabilistic and measure
arguments in Section~\ref{sec:finite-state}---are not machine-checked in Lean.
At the time of writing, these paper proofs have been checked by
line-by-line re-derivation within the same AI-assisted workflow and
cross-checked against the reproducible scripts and Lean audits described
above; they have not yet been examined by an independent human referee.

\begin{thebibliography}{99}
\bibitem{BDS} A.~B\'ir\'o, J.-M.~Deshouillers, V.~T.~S\'os, \emph{Good
approximation and characterization of subgroups of $\R/\Z$}, Studia Sci.\ Math.\ Hungar.\ 38 (2001), 97--113.

\bibitem{BV} V.~Beresnevich and S.~Velani, \emph{A mass transference principle
and the Duffin--Schaeffer conjecture}, Ann.\ of Math.\ (2) 164 (2006), 971--992.

\bibitem{BLR} Y.~Bugeaud, L.~Liao, M.~Rams, \emph{Metrical results on the
distribution of fractional parts of powers of real numbers}, Proc.\ Edinb.\ Math.\ Soc.\ (2) 62 (2019), 505--521, doi:10.1017/\allowbreak S0013091518000585; arXiv:1710.03596. (Cor.~8, the hypothesis is $Q_{j+1}/Q_j\to\infty$.)

\bibitem{CMMS} C.~Cabrelli, F.~Mendivil, U.~M.~Molter, R.~Shonkwiler,
\emph{On the Hausdorff $h$-measure of Cantor sets}, Pacific J.\ Math.\ 217
(2004), 45--59, doi:10.2140/\allowbreak pjm.2004.217.45.

\bibitem{CLU} V.~Chousionis, D.~Leykekhman, M.~Urba\'nski, \emph{On the
dimension spectrum of infinite subsystems of continued fractions}, Trans.\ Amer.\ Math.\ Soc.\ 373 (2020), 1009--1042, arXiv:1805.11904.

\bibitem{Falconer} K.~Falconer, \emph{Fractal Geometry: Mathematical
Foundations and Applications}, 2nd ed., Wiley, 2003. (Countable sets have
Hausdorff dimension~$0$.)

\bibitem{FLLW} A.-H.~Fan, L.-M.~Liao, B.-W.~Wang, J.~Wu, \emph{On Khintchine
exponents and Lyapunov exponents of continued fractions}, Ergodic Theory Dynam.\ Systems 29 (2009), 73--109.

\bibitem{FWW} D.-J.~Feng, Z.-Y.~Wen, J.~Wu, \emph{Some dimensional results
for homogeneous Moran sets}, Sci.\ China Ser.\ A 40 (1997), 475--482,
doi:10.1007/\allowbreak BF02896955.

\bibitem{KL} C.~Kraaikamp and P.~Liardet, \emph{Good approximations and
continued fractions}, Proc.\ Amer.\ Math.\ Soc.\ 112 (1991), 303--309.

\bibitem{LimaMoreira} D.~Lima and C.~G.~Moreira, \emph{Phase transitions on
the Markov and Lagrange dynamical spectra}, Ann.\ Inst.\ H.~Poincar\'e
C Anal.\ Non Lin\'eaire 38 (2021), 1429--1459,
doi:10.1016/\allowbreak j.anihpc.2020.11.007; arXiv:1801.04636.

\bibitem{Moreira} C.~G.~Moreira, \emph{Geometric properties of the Markov and
Lagrange spectra}, Ann.\ of Math.\ (2) 188 (2018), 145--170,
doi:10.4007/\allowbreak annals.2018.188.1.3.

\bibitem{MEA} C.~G.~Moreira, H.~Erazo, N.~Angelini,
\emph{Generalized Hausdorff measure of irrationals with Lagrange value exactly 3},
arXiv:2511.10328v3.

{\sloppy
\bibitem{MW} R.~D.~Mauldin and S.~C.~Williams,
\emph{Hausdorff dimension in graph directed constructions},
Trans.~Amer.~Math.~Soc.~309 (1988), no.~2, 811--829.
}

\bibitem{MarcusSofic} B.~Marcus, \emph{Sofic systems and encoding data},
IEEE Trans.~Inform.~Theory \textbf{31} (1985), no.~3, 366--377.

\bibitem{OlsenFrequency} L.~Olsen,
\emph{Distribution of digits in integers: fractal dimensions and zeta functions},
Acta Arith. \textbf{105} (2002), no.~3, 253--277,
doi:10.4064/\allowbreak aa105-3-4.

\bibitem{ElekesKeletiFrequency} M.~Elekes and T.~Keleti,
\emph{Borel sets which are null or non-$\sigma$-finite for every
translation invariant measure}, arXiv:1109.5309.


\bibitem{GenericZeroInfinity}
J.~Ma and Z.~Wen, \emph{Hausdorff and packing measure of sets of
generic points: a zero-infinity law}, J.~London Math.~Soc.
\textbf{69} (2004), 383--406.

\bibitem{ZiemianRotation}
K.~Ziemian, \emph{Rotation sets for subshifts of finite type},
Fund. Math. \textbf{146} (1995), 189--201.

\bibitem{Phase1} \emph{On the irrationality of $e+\pi$: proven results,
new reductions, and a systematic obstruction analysis},
companion preprint (Phase~1), November 2026. Theorem~2
(classification) and Corollary~3 (initial spectral gap) of that
preprint are kernel-checked in its Lean~4 formalization package
\texttt{EpiFormal}, modules \texttt{Classification.lean} and
\texttt{Spectral\allowbreak Gap.lean} (dated 2026-11-15).
\bibitem{Feller1943}
W.~Feller, \emph{The general form of the so-called law of the iterated
logarithm}, Trans.\ Amer.\ Math.\ Soc.\ \textbf{54} (1943), 373--402,
doi:10.1090/S0002-9947-1943-0009263-7.
\bibitem{CurtissMGF}
J.~H.~Curtiss, \emph{A note on the theory of moment generating
functions}, Ann.\ Math.\ Statist.\ \textbf{13} (1942), 430--433,
doi:10.1214/aoms/1177731541.
\bibitem{JaffuelBarrier}
B.~Jaffuel, \emph{The critical barrier for the survival of branching
random walk with absorption}, Ann.\ Inst.\ H.\ Poincar\'e Probab.\
Statist.\ \textbf{48} (2012), 989--1009, doi:10.1214/11-AIHP453.
\bibitem{WhittFCLT}
W.~Whitt, \emph{Proofs of the martingale FCLT},
Probab.\ Surv.\ \textbf{4} (2007), 268--302,
doi:10.1214/07-PS122.
\end{thebibliography}

\end{document}

